
%% $Id$

\documentclass[12pt,a4paper,fleqn]{report}
\usepackage{latexsym,graphicx}
\usepackage[refpage]{nomencl}
\usepackage{../iman,../extra,../isar,../proof}
\usepackage[nohyphen,strings]{../underscore}
\usepackage{Thy/document/isabelle,Thy/document/isabellesym}
\usepackage{style}
\usepackage{../pdfsetup}


\hyphenation{Isabelle}
\hyphenation{Isar}

\isadroptag{theory}
\title{\includegraphics[scale=0.5]{isabelle_isar}
  \\[4ex] The Isabelle/Isar Implementation}
\author{\emph{Makarius Wenzel}}

%FIXME
%\makeglossary

\makeindex


\begin{document}

\maketitle 

\begin{abstract}
  We describe the key concepts underlying the Isabelle/Isar
  implementation, including ML references for the most important
  functions.  The aim is to give some insight into the overall system
  architecture, and provide clues on implementing applications within
  this framework.
\end{abstract}

\vspace*{2.5cm}
\begin{quote}
  
  {\small\em Isabelle was not designed; it evolved.  Not everyone
    likes this idea.  Specification experts rightly abhor
    trial-and-error programming.  They suggest that no one should
    write a program without first writing a complete formal
    specification. But university departments are not software houses.
    Programs like Isabelle are not products: when they have served
    their purpose, they are discarded.}
  
  Lawrence C. Paulson, ``Isabelle: The Next 700 Theorem Provers''

  \vspace*{1cm}
  
  {\small\em As I did 20 years ago, I still fervently believe that the
    only way to make software secure, reliable, and fast is to make it
    small.  Fight features.}
  
  Andrew S. Tanenbaum

\end{quote}

\thispagestyle{empty}\clearpage

\pagenumbering{roman}
\tableofcontents
\listoffigures
\clearfirst

%
\chapter{Introduction}

\section{Quick start}

Isar is already part of Isabelle (as of version Isabelle99, or later).  The
\texttt{isabelle} binary provides option \texttt{-I} to run the Isar
interaction loop at startup, rather than the plain ML top-level.  Thus the
quickest way to do anything with Isabelle/Isar is as follows:
\begin{ttbox}
isabelle -I HOL\medskip
\out{> Welcome to Isabelle/HOL (Isabelle99)}\medskip
theory Foo = Main:
constdefs foo :: nat  "foo == 1";
lemma "0 < foo" by (simp add: foo_def);
end
\end{ttbox}
Note that any Isabelle/Isar command may be retracted by \texttt{undo}.

Plain TTY-based interaction like this used to be quite feasible with
traditional tactic based theorem proving, but developing Isar documents
demands some better user-interface support.  \emph{Proof~General}\index{Proof
  General} of LFCS Edinburgh \cite{proofgeneral} offers a generic Emacs-based
environment for interactive theorem provers that does all the cut-and-paste
and forward-backward walk through the document in a very neat way.  Note that
in Isabelle/Isar, the current position within a partial proof document is more
informative than the actual proof state.  Thus Proof~General provides the
canonical working environment for Isabelle/Isar, both for getting acquainted
(e.g.\ by replaying existing Isar documents) and serious production work.

\medskip

The easiest way to use Proof~General is to make it the default Isabelle user
interface.  Just say something like this in your Isabelle settings file (cf.\ 
\cite{isabelle-sys}):
\begin{ttbox}
ISABELLE_INTERFACE=\$ISABELLE_HOME/contrib/ProofGeneral/isar/interface
PROOFGENERAL_OPTIONS=""
\end{ttbox}
You may have to change \texttt{\$ISABELLE_HOME/contrib/ProofGeneral} to the
actual installation directory of Proof~General.  From now on, the capital
\texttt{Isabelle} executable refers to the \texttt{ProofGeneral/isar}
interface.  Its usage is as follows:
\begin{ttbox}
Usage: interface [OPTIONS] [FILES ...]

  Options are:
    -l NAME      logic image name (default $ISABELLE_LOGIC=HOL)
    -p NAME      Emacs program name (default xemacs)
    -u BOOL      use .emacs file (default true)
    -w BOOL      use window system (default true)

  Starts Proof General for Isabelle/Isar with proof documents FILES
  (default Scratch.thy).

  PROOFGENERAL_OPTIONS=
\end{ttbox}
Apart from the command line, the defaults for these options may be overridden
via the \texttt{PROOFGENERAL_OPTIONS} setting as well.  For example, plain GNU
Emacs may be configured as follows:
\begin{ttbox}
PROOFGENERAL_OPTIONS="-p emacs"
\end{ttbox}

With the proper Isabelle interface setup, Isar documents may now be edited by
visiting appropriate theory files, e.g.\ 
\begin{ttbox}
Isabelle \({\langle}isabellehome{\rangle}\)/src/HOL/Isar_examples/BasicLogic.thy
\end{ttbox}
Users of XEmacs may note the tool bar for navigating forward and backward
through the text.  Consult the Proof~General documentation \cite{proofgeneral}
for further basic commands, such as \texttt{c-c return} or \texttt{c-c u}.

\medskip

Occasionally, a user's \texttt{.emacs} contains material that is incompatible
with the version of (X)Emacs that Proof~General prefers.  Then proper startup
may be still achieved by using the \texttt{-u false} option.\footnote{Also
  note that the Emacs lisp files
  \texttt{\$ISABELLE_HOME/etc/proofgeneral-settings.el} and
  \texttt{\$ISABELLE_HOME_USER/etc/proofgeneral-settings.el} are automatically
  loaded by Proof~General if invoked via the interface wrapper script.}


\section{How to write Isar proofs anyway?}

This is one of the key questions, of course.  Isar offers a rather different
approach to formal proof documents than plain old tactic scripts.  Experienced
users of existing interactive theorem proving systems may have to learn
thinking different in order to make effective use of Isabelle/Isar.  On the
other hand, Isabelle/Isar comes much closer to existing mathematical practice
of formal proof, so users with less experience in old-style tactical proving,
but a good understanding of mathematical proof might cope with Isar even
better.

Unfortunately, there is no tutorial on Isabelle/Isar available yet.  This
document really is a \emph{reference manual}.  Nevertheless, we will give some
discussions of the general principles underlying Isar in
chapter~\ref{ch:basics}, and provide some clues of how these may be put into
practice.  Some more background information on Isar is given in
\cite{Wenzel:1999:TPHOL}.  Furthermore, there are several examples distributed
with Isabelle (see directory \texttt{HOL/Isar_examples}).


%%% Local Variables: 
%%% mode: latex
%%% TeX-master: "isar-ref"
%%% End: 

%
\begin{isabellebody}%
\def\isabellecontext{prelim}%
%
\isadelimtheory
\isanewline
\isanewline
\isanewline
%
\endisadelimtheory
%
\isatagtheory
\isacommand{theory}\isamarkupfalse%
\ prelim\ \isakeyword{imports}\ base\ \isakeyword{begin}%
\endisatagtheory
{\isafoldtheory}%
%
\isadelimtheory
%
\endisadelimtheory
%
\isamarkupchapter{Preliminaries%
}
\isamarkuptrue%
%
\isamarkupsection{Contexts \label{sec:context}%
}
\isamarkuptrue%
%
\begin{isamarkuptext}%
A logical context represents the background that is required for
  formulating statements and composing proofs.  It acts as a medium to
  produce formal content, depending on earlier material (declarations,
  results etc.).

  For example, derivations within the Isabelle/Pure logic can be
  described as a judgment \isa{{\isasymGamma}\ {\isasymturnstile}\isactrlsub {\isasymTheta}\ {\isasymphi}}, which means that a
  proposition \isa{{\isasymphi}} is derivable from hypotheses \isa{{\isasymGamma}}
  within the theory \isa{{\isasymTheta}}.  There are logical reasons for
  keeping \isa{{\isasymTheta}} and \isa{{\isasymGamma}} separate: theories can be
  liberal about supporting type constructors and schematic
  polymorphism of constants and axioms, while the inner calculus of
  \isa{{\isasymGamma}\ {\isasymturnstile}\ {\isasymphi}} is strictly limited to Simple Type Theory (with
  fixed type variables in the assumptions).

  \medskip Contexts and derivations are linked by the following key
  principles:

  \begin{itemize}

  \item Transfer: monotonicity of derivations admits results to be
  transferred into a \emph{larger} context, i.e.\ \isa{{\isasymGamma}\ {\isasymturnstile}\isactrlsub {\isasymTheta}\ {\isasymphi}} implies \isa{{\isasymGamma}{\isacharprime}\ {\isasymturnstile}\isactrlsub {\isasymTheta}\isactrlsub {\isacharprime}\ {\isasymphi}} for contexts \isa{{\isasymTheta}{\isacharprime}\ {\isasymsupseteq}\ {\isasymTheta}} and \isa{{\isasymGamma}{\isacharprime}\ {\isasymsupseteq}\ {\isasymGamma}}.

  \item Export: discharge of hypotheses admits results to be exported
  into a \emph{smaller} context, i.e.\ \isa{{\isasymGamma}{\isacharprime}\ {\isasymturnstile}\isactrlsub {\isasymTheta}\ {\isasymphi}}
  implies \isa{{\isasymGamma}\ {\isasymturnstile}\isactrlsub {\isasymTheta}\ {\isasymDelta}\ {\isasymLongrightarrow}\ {\isasymphi}} where \isa{{\isasymGamma}{\isacharprime}\ {\isasymsupseteq}\ {\isasymGamma}} and
  \isa{{\isasymDelta}\ {\isacharequal}\ {\isasymGamma}{\isacharprime}\ {\isacharminus}\ {\isasymGamma}}.  Note that \isa{{\isasymTheta}} remains unchanged here,
  only the \isa{{\isasymGamma}} part is affected.

  \end{itemize}

  \medskip By modeling the main characteristics of the primitive
  \isa{{\isasymTheta}} and \isa{{\isasymGamma}} above, and abstracting over any
  particular logical content, we arrive at the fundamental notions of
  \emph{theory context} and \emph{proof context} in Isabelle/Isar.
  These implement a certain policy to manage arbitrary \emph{context
  data}.  There is a strongly-typed mechanism to declare new kinds of
  data at compile time.

  The internal bootstrap process of Isabelle/Pure eventually reaches a
  stage where certain data slots provide the logical content of \isa{{\isasymTheta}} and \isa{{\isasymGamma}} sketched above, but this does not stop there!
  Various additional data slots support all kinds of mechanisms that
  are not necessarily part of the core logic.

  For example, there would be data for canonical introduction and
  elimination rules for arbitrary operators (depending on the
  object-logic and application), which enables users to perform
  standard proof steps implicitly (cf.\ the \isa{rule} method
  \cite{isabelle-isar-ref}).

  \medskip Thus Isabelle/Isar is able to bring forth more and more
  concepts successively.  In particular, an object-logic like
  Isabelle/HOL continues the Isabelle/Pure setup by adding specific
  components for automated reasoning (classical reasoner, tableau
  prover, structured induction etc.) and derived specification
  mechanisms (inductive predicates, recursive functions etc.).  All of
  this is ultimately based on the generic data management by theory
  and proof contexts introduced here.%
\end{isamarkuptext}%
\isamarkuptrue%
%
\isamarkupsubsection{Theory context \label{sec:context-theory}%
}
\isamarkuptrue%
%
\begin{isamarkuptext}%
\glossary{Theory}{FIXME}

  A \emph{theory} is a data container with explicit named and unique
  identifier.  Theories are related by a (nominal) sub-theory
  relation, which corresponds to the dependency graph of the original
  construction; each theory is derived from a certain sub-graph of
  ancestor theories.

  The \isa{merge} operation produces the least upper bound of two
  theories, which actually degenerates into absorption of one theory
  into the other (due to the nominal sub-theory relation).

  The \isa{begin} operation starts a new theory by importing
  several parent theories and entering a special \isa{draft} mode,
  which is sustained until the final \isa{end} operation.  A draft
  theory acts like a linear type, where updates invalidate earlier
  versions.  An invalidated draft is called ``stale''.

  The \isa{checkpoint} operation produces an intermediate stepping
  stone that will survive the next update: both the original and the
  changed theory remain valid and are related by the sub-theory
  relation.  Checkpointing essentially recovers purely functional
  theory values, at the expense of some extra internal bookkeeping.

  The \isa{copy} operation produces an auxiliary version that has
  the same data content, but is unrelated to the original: updates of
  the copy do not affect the original, neither does the sub-theory
  relation hold.

  \medskip The example in \figref{fig:ex-theory} below shows a theory
  graph derived from \isa{Pure}, with theory \isa{Length}
  importing \isa{Nat} and \isa{List}.  The body of \isa{Length} consists of a sequence of updates, working mostly on
  drafts.  Intermediate checkpoints may occur as well, due to the
  history mechanism provided by the Isar top-level, cf.\
  \secref{sec:isar-toplevel}.

  \begin{figure}[htb]
  \begin{center}
  \begin{tabular}{rcccl}
        &            & \isa{Pure} \\
        &            & \isa{{\isasymdown}} \\
        &            & \isa{FOL} \\
        & $\swarrow$ &              & $\searrow$ & \\
  \isa{Nat} &    &              &            & \isa{List} \\
        & $\searrow$ &              & $\swarrow$ \\
        &            & \isa{Length} \\
        &            & \multicolumn{3}{l}{~~\hyperlink{keyword.imports}{\mbox{\isa{\isakeyword{imports}}}}} \\
        &            & \multicolumn{3}{l}{~~\hyperlink{keyword.begin}{\mbox{\isa{\isakeyword{begin}}}}} \\
        &            & $\vdots$~~ \\
        &            & \isa{{\isasymbullet}}~~ \\
        &            & $\vdots$~~ \\
        &            & \isa{{\isasymbullet}}~~ \\
        &            & $\vdots$~~ \\
        &            & \multicolumn{3}{l}{~~\hyperlink{command.end}{\mbox{\isa{\isacommand{end}}}}} \\
  \end{tabular}
  \caption{A theory definition depending on ancestors}\label{fig:ex-theory}
  \end{center}
  \end{figure}

  \medskip There is a separate notion of \emph{theory reference} for
  maintaining a live link to an evolving theory context: updates on
  drafts are propagated automatically.  Dynamic updating stops after
  an explicit \isa{end} only.

  Derived entities may store a theory reference in order to indicate
  the context they belong to.  This implicitly assumes monotonic
  reasoning, because the referenced context may become larger without
  further notice.%
\end{isamarkuptext}%
\isamarkuptrue%
%
\isadelimmlref
%
\endisadelimmlref
%
\isatagmlref
%
\begin{isamarkuptext}%
\begin{mldecls}
  \indexmltype{theory}\verb|type theory| \\
  \indexml{Theory.subthy}\verb|Theory.subthy: theory * theory -> bool| \\
  \indexml{Theory.merge}\verb|Theory.merge: theory * theory -> theory| \\
  \indexml{Theory.checkpoint}\verb|Theory.checkpoint: theory -> theory| \\
  \indexml{Theory.copy}\verb|Theory.copy: theory -> theory| \\
  \end{mldecls}
  \begin{mldecls}
  \indexmltype{theory\_ref}\verb|type theory_ref| \\
  \indexml{Theory.deref}\verb|Theory.deref: theory_ref -> theory| \\
  \indexml{Theory.check\_thy}\verb|Theory.check_thy: theory -> theory_ref| \\
  \end{mldecls}

  \begin{description}

  \item \verb|theory| represents theory contexts.  This is
  essentially a linear type!  Most operations destroy the original
  version, which then becomes ``stale''.

  \item \verb|Theory.subthy|~\isa{{\isacharparenleft}thy\isactrlsub {\isadigit{1}}{\isacharcomma}\ thy\isactrlsub {\isadigit{2}}{\isacharparenright}}
  compares theories according to the inherent graph structure of the
  construction.  This sub-theory relation is a nominal approximation
  of inclusion (\isa{{\isasymsubseteq}}) of the corresponding content.

  \item \verb|Theory.merge|~\isa{{\isacharparenleft}thy\isactrlsub {\isadigit{1}}{\isacharcomma}\ thy\isactrlsub {\isadigit{2}}{\isacharparenright}}
  absorbs one theory into the other.  This fails for unrelated
  theories!

  \item \verb|Theory.checkpoint|~\isa{thy} produces a safe
  stepping stone in the linear development of \isa{thy}.  The next
  update will result in two related, valid theories.

  \item \verb|Theory.copy|~\isa{thy} produces a variant of \isa{thy} that holds a copy of the same data.  The result is not
  related to the original; the original is unchanched.

  \item \verb|theory_ref| represents a sliding reference to an
  always valid theory; updates on the original are propagated
  automatically.

  \item \verb|Theory.deref|~\isa{thy{\isacharunderscore}ref} turns a \verb|theory_ref| into an \verb|theory| value.  As the referenced
  theory evolves monotonically over time, later invocations of \verb|Theory.deref| may refer to a larger context.

  \item \verb|Theory.check_thy|~\isa{thy} produces a \verb|theory_ref| from a valid \verb|theory| value.

  \end{description}%
\end{isamarkuptext}%
\isamarkuptrue%
%
\endisatagmlref
{\isafoldmlref}%
%
\isadelimmlref
%
\endisadelimmlref
%
\isamarkupsubsection{Proof context \label{sec:context-proof}%
}
\isamarkuptrue%
%
\begin{isamarkuptext}%
\glossary{Proof context}{The static context of a structured proof,
  acts like a local ``theory'' of the current portion of Isar proof
  text, generalizes the idea of local hypotheses \isa{{\isasymGamma}} in
  judgments \isa{{\isasymGamma}\ {\isasymturnstile}\ {\isasymphi}} of natural deduction calculi.  There is a
  generic notion of introducing and discharging hypotheses.
  Arbritrary auxiliary context data may be adjoined.}

  A proof context is a container for pure data with a back-reference
  to the theory it belongs to.  The \isa{init} operation creates a
  proof context from a given theory.  Modifications to draft theories
  are propagated to the proof context as usual, but there is also an
  explicit \isa{transfer} operation to force resynchronization
  with more substantial updates to the underlying theory.  The actual
  context data does not require any special bookkeeping, thanks to the
  lack of destructive features.

  Entities derived in a proof context need to record inherent logical
  requirements explicitly, since there is no separate context
  identification as for theories.  For example, hypotheses used in
  primitive derivations (cf.\ \secref{sec:thms}) are recorded
  separately within the sequent \isa{{\isasymGamma}\ {\isasymturnstile}\ {\isasymphi}}, just to make double
  sure.  Results could still leak into an alien proof context do to
  programming errors, but Isabelle/Isar includes some extra validity
  checks in critical positions, notably at the end of a sub-proof.

  Proof contexts may be manipulated arbitrarily, although the common
  discipline is to follow block structure as a mental model: a given
  context is extended consecutively, and results are exported back
  into the original context.  Note that the Isar proof states model
  block-structured reasoning explicitly, using a stack of proof
  contexts internally, cf.\ \secref{sec:isar-proof-state}.%
\end{isamarkuptext}%
\isamarkuptrue%
%
\isadelimmlref
%
\endisadelimmlref
%
\isatagmlref
%
\begin{isamarkuptext}%
\begin{mldecls}
  \indexmltype{Proof.context}\verb|type Proof.context| \\
  \indexml{ProofContext.init}\verb|ProofContext.init: theory -> Proof.context| \\
  \indexml{ProofContext.theory\_of}\verb|ProofContext.theory_of: Proof.context -> theory| \\
  \indexml{ProofContext.transfer}\verb|ProofContext.transfer: theory -> Proof.context -> Proof.context| \\
  \end{mldecls}

  \begin{description}

  \item \verb|Proof.context| represents proof contexts.  Elements
  of this type are essentially pure values, with a sliding reference
  to the background theory.

  \item \verb|ProofContext.init|~\isa{thy} produces a proof context
  derived from \isa{thy}, initializing all data.

  \item \verb|ProofContext.theory_of|~\isa{ctxt} selects the
  background theory from \isa{ctxt}, dereferencing its internal
  \verb|theory_ref|.

  \item \verb|ProofContext.transfer|~\isa{thy\ ctxt} promotes the
  background theory of \isa{ctxt} to the super theory \isa{thy}.

  \end{description}%
\end{isamarkuptext}%
\isamarkuptrue%
%
\endisatagmlref
{\isafoldmlref}%
%
\isadelimmlref
%
\endisadelimmlref
%
\isamarkupsubsection{Generic contexts \label{sec:generic-context}%
}
\isamarkuptrue%
%
\begin{isamarkuptext}%
A generic context is the disjoint sum of either a theory or proof
  context.  Occasionally, this enables uniform treatment of generic
  context data, typically extra-logical information.  Operations on
  generic contexts include the usual injections, partial selections,
  and combinators for lifting operations on either component of the
  disjoint sum.

  Moreover, there are total operations \isa{theory{\isacharunderscore}of} and \isa{proof{\isacharunderscore}of} to convert a generic context into either kind: a theory
  can always be selected from the sum, while a proof context might
  have to be constructed by an ad-hoc \isa{init} operation.%
\end{isamarkuptext}%
\isamarkuptrue%
%
\isadelimmlref
%
\endisadelimmlref
%
\isatagmlref
%
\begin{isamarkuptext}%
\begin{mldecls}
  \indexmltype{Context.generic}\verb|type Context.generic| \\
  \indexml{Context.theory\_of}\verb|Context.theory_of: Context.generic -> theory| \\
  \indexml{Context.proof\_of}\verb|Context.proof_of: Context.generic -> Proof.context| \\
  \end{mldecls}

  \begin{description}

  \item \verb|Context.generic| is the direct sum of \verb|theory| and \verb|Proof.context|, with the datatype
  constructors \verb|Context.Theory| and \verb|Context.Proof|.

  \item \verb|Context.theory_of|~\isa{context} always produces a
  theory from the generic \isa{context}, using \verb|ProofContext.theory_of| as required.

  \item \verb|Context.proof_of|~\isa{context} always produces a
  proof context from the generic \isa{context}, using \verb|ProofContext.init| as required (note that this re-initializes the
  context data with each invocation).

  \end{description}%
\end{isamarkuptext}%
\isamarkuptrue%
%
\endisatagmlref
{\isafoldmlref}%
%
\isadelimmlref
%
\endisadelimmlref
%
\isamarkupsubsection{Context data \label{sec:context-data}%
}
\isamarkuptrue%
%
\begin{isamarkuptext}%
The main purpose of theory and proof contexts is to manage arbitrary
  data.  New data types can be declared incrementally at compile time.
  There are separate declaration mechanisms for any of the three kinds
  of contexts: theory, proof, generic.

  \paragraph{Theory data} may refer to destructive entities, which are
  maintained in direct correspondence to the linear evolution of
  theory values, including explicit copies.\footnote{Most existing
  instances of destructive theory data are merely historical relics
  (e.g.\ the destructive theorem storage, and destructive hints for
  the Simplifier and Classical rules).}  A theory data declaration
  needs to implement the following SML signature:

  \medskip
  \begin{tabular}{ll}
  \isa{{\isasymtype}\ T} & representing type \\
  \isa{{\isasymval}\ empty{\isacharcolon}\ T} & empty default value \\
  \isa{{\isasymval}\ copy{\isacharcolon}\ T\ {\isasymrightarrow}\ T} & refresh impure data \\
  \isa{{\isasymval}\ extend{\isacharcolon}\ T\ {\isasymrightarrow}\ T} & re-initialize on import \\
  \isa{{\isasymval}\ merge{\isacharcolon}\ T\ {\isasymtimes}\ T\ {\isasymrightarrow}\ T} & join on import \\
  \end{tabular}
  \medskip

  \noindent The \isa{empty} value acts as initial default for
  \emph{any} theory that does not declare actual data content; \isa{copy} maintains persistent integrity for impure data, it is just
  the identity for pure values; \isa{extend} is acts like a
  unitary version of \isa{merge}, both operations should also
  include the functionality of \isa{copy} for impure data.

  \paragraph{Proof context data} is purely functional.  A declaration
  needs to implement the following SML signature:

  \medskip
  \begin{tabular}{ll}
  \isa{{\isasymtype}\ T} & representing type \\
  \isa{{\isasymval}\ init{\isacharcolon}\ theory\ {\isasymrightarrow}\ T} & produce initial value \\
  \end{tabular}
  \medskip

  \noindent The \isa{init} operation is supposed to produce a pure
  value from the given background theory.

  \paragraph{Generic data} provides a hybrid interface for both theory
  and proof data.  The declaration is essentially the same as for
  (pure) theory data, without \isa{copy}.  The \isa{init}
  operation for proof contexts merely selects the current data value
  from the background theory.

  \bigskip A data declaration of type \isa{T} results in the
  following interface:

  \medskip
  \begin{tabular}{ll}
  \isa{init{\isacharcolon}\ theory\ {\isasymrightarrow}\ theory} \\
  \isa{get{\isacharcolon}\ context\ {\isasymrightarrow}\ T} \\
  \isa{put{\isacharcolon}\ T\ {\isasymrightarrow}\ context\ {\isasymrightarrow}\ context} \\
  \isa{map{\isacharcolon}\ {\isacharparenleft}T\ {\isasymrightarrow}\ T{\isacharparenright}\ {\isasymrightarrow}\ context\ {\isasymrightarrow}\ context} \\
  \end{tabular}
  \medskip

  \noindent Here \isa{init} is only applicable to impure theory
  data to install a fresh copy persistently (destructive update on
  uninitialized has no permanent effect).  The other operations provide
  access for the particular kind of context (theory, proof, or generic
  context).  Note that this is a safe interface: there is no other way
  to access the corresponding data slot of a context.  By keeping
  these operations private, a component may maintain abstract values
  authentically, without other components interfering.%
\end{isamarkuptext}%
\isamarkuptrue%
%
\isadelimmlref
%
\endisadelimmlref
%
\isatagmlref
%
\begin{isamarkuptext}%
\begin{mldecls}
  \indexmlfunctor{TheoryDataFun}\verb|functor TheoryDataFun| \\
  \indexmlfunctor{ProofDataFun}\verb|functor ProofDataFun| \\
  \indexmlfunctor{GenericDataFun}\verb|functor GenericDataFun| \\
  \end{mldecls}

  \begin{description}

  \item \verb|TheoryDataFun|\isa{{\isacharparenleft}spec{\isacharparenright}} declares data for
  type \verb|theory| according to the specification provided as
  argument structure.  The resulting structure provides data init and
  access operations as described above.

  \item \verb|ProofDataFun|\isa{{\isacharparenleft}spec{\isacharparenright}} is analogous to
  \verb|TheoryDataFun| for type \verb|Proof.context|.

  \item \verb|GenericDataFun|\isa{{\isacharparenleft}spec{\isacharparenright}} is analogous to
  \verb|TheoryDataFun| for type \verb|Context.generic|.

  \end{description}%
\end{isamarkuptext}%
\isamarkuptrue%
%
\endisatagmlref
{\isafoldmlref}%
%
\isadelimmlref
%
\endisadelimmlref
%
\isamarkupsection{Names \label{sec:names}%
}
\isamarkuptrue%
%
\begin{isamarkuptext}%
In principle, a name is just a string, but there are various
  convention for encoding additional structure.  For example, ``\isa{Foo{\isachardot}bar{\isachardot}baz}'' is considered as a qualified name consisting of
  three basic name components.  The individual constituents of a name
  may have further substructure, e.g.\ the string
  ``\verb,\,\verb,<alpha>,'' encodes as a single symbol.%
\end{isamarkuptext}%
\isamarkuptrue%
%
\isamarkupsubsection{Strings of symbols%
}
\isamarkuptrue%
%
\begin{isamarkuptext}%
\glossary{Symbol}{The smallest unit of text in Isabelle, subsumes
  plain ASCII characters as well as an infinite collection of named
  symbols (for greek, math etc.).}

  A \emph{symbol} constitutes the smallest textual unit in Isabelle
  --- raw characters are normally not encountered at all.  Isabelle
  strings consist of a sequence of symbols, represented as a packed
  string or a list of strings.  Each symbol is in itself a small
  string, which has either one of the following forms:

  \begin{enumerate}

  \item a single ASCII character ``\isa{c}'', for example
  ``\verb,a,'',

  \item a regular symbol ``\verb,\,\verb,<,\isa{ident}\verb,>,'',
  for example ``\verb,\,\verb,<alpha>,'',

  \item a control symbol ``\verb,\,\verb,<^,\isa{ident}\verb,>,'',
  for example ``\verb,\,\verb,<^bold>,'',

  \item a raw symbol ``\verb,\,\verb,<^raw:,\isa{text}\verb,>,''
  where \isa{text} constists of printable characters excluding
  ``\verb,.,'' and ``\verb,>,'', for example
  ``\verb,\,\verb,<^raw:$\sum_{i = 1}^n$>,'',

  \item a numbered raw control symbol ``\verb,\,\verb,<^raw,\isa{n}\verb,>, where \isa{n} consists of digits, for example
  ``\verb,\,\verb,<^raw42>,''.

  \end{enumerate}

  \noindent The \isa{ident} syntax for symbol names is \isa{letter\ {\isacharparenleft}letter\ {\isacharbar}\ digit{\isacharparenright}\isactrlsup {\isacharasterisk}}, where \isa{letter\ {\isacharequal}\ A{\isachardot}{\isachardot}Za{\isachardot}{\isachardot}z} and \isa{digit\ {\isacharequal}\ {\isadigit{0}}{\isachardot}{\isachardot}{\isadigit{9}}}.  There are infinitely many
  regular symbols and control symbols, but a fixed collection of
  standard symbols is treated specifically.  For example,
  ``\verb,\,\verb,<alpha>,'' is classified as a letter, which means it
  may occur within regular Isabelle identifiers.

  Since the character set underlying Isabelle symbols is 7-bit ASCII
  and 8-bit characters are passed through transparently, Isabelle may
  also process Unicode/UCS data in UTF-8 encoding.  Unicode provides
  its own collection of mathematical symbols, but there is no built-in
  link to the standard collection of Isabelle.

  \medskip Output of Isabelle symbols depends on the print mode
  (\secref{FIXME}).  For example, the standard {\LaTeX} setup of the
  Isabelle document preparation system would present
  ``\verb,\,\verb,<alpha>,'' as \isa{{\isasymalpha}}, and
  ``\verb,\,\verb,<^bold>,\verb,\,\verb,<alpha>,'' as \isa{\isactrlbold {\isasymalpha}}.%
\end{isamarkuptext}%
\isamarkuptrue%
%
\isadelimmlref
%
\endisadelimmlref
%
\isatagmlref
%
\begin{isamarkuptext}%
\begin{mldecls}
  \indexmltype{Symbol.symbol}\verb|type Symbol.symbol| \\
  \indexml{Symbol.explode}\verb|Symbol.explode: string -> Symbol.symbol list| \\
  \indexml{Symbol.is\_letter}\verb|Symbol.is_letter: Symbol.symbol -> bool| \\
  \indexml{Symbol.is\_digit}\verb|Symbol.is_digit: Symbol.symbol -> bool| \\
  \indexml{Symbol.is\_quasi}\verb|Symbol.is_quasi: Symbol.symbol -> bool| \\
  \indexml{Symbol.is\_blank}\verb|Symbol.is_blank: Symbol.symbol -> bool| \\
  \end{mldecls}
  \begin{mldecls}
  \indexmltype{Symbol.sym}\verb|type Symbol.sym| \\
  \indexml{Symbol.decode}\verb|Symbol.decode: Symbol.symbol -> Symbol.sym| \\
  \end{mldecls}

  \begin{description}

  \item \verb|Symbol.symbol| represents individual Isabelle
  symbols; this is an alias for \verb|string|.

  \item \verb|Symbol.explode|~\isa{str} produces a symbol list
  from the packed form.  This function supercedes \verb|String.explode| for virtually all purposes of manipulating text in
  Isabelle!

  \item \verb|Symbol.is_letter|, \verb|Symbol.is_digit|, \verb|Symbol.is_quasi|, \verb|Symbol.is_blank| classify standard
  symbols according to fixed syntactic conventions of Isabelle, cf.\
  \cite{isabelle-isar-ref}.

  \item \verb|Symbol.sym| is a concrete datatype that represents
  the different kinds of symbols explicitly, with constructors \verb|Symbol.Char|, \verb|Symbol.Sym|, \verb|Symbol.Ctrl|, \verb|Symbol.Raw|.

  \item \verb|Symbol.decode| converts the string representation of a
  symbol into the datatype version.

  \end{description}%
\end{isamarkuptext}%
\isamarkuptrue%
%
\endisatagmlref
{\isafoldmlref}%
%
\isadelimmlref
%
\endisadelimmlref
%
\isamarkupsubsection{Basic names \label{sec:basic-names}%
}
\isamarkuptrue%
%
\begin{isamarkuptext}%
A \emph{basic name} essentially consists of a single Isabelle
  identifier.  There are conventions to mark separate classes of basic
  names, by attaching a suffix of underscores (\isa{{\isacharunderscore}}): one
  underscore means \emph{internal name}, two underscores means
  \emph{Skolem name}, three underscores means \emph{internal Skolem
  name}.

  For example, the basic name \isa{foo} has the internal version
  \isa{foo{\isacharunderscore}}, with Skolem versions \isa{foo{\isacharunderscore}{\isacharunderscore}} and \isa{foo{\isacharunderscore}{\isacharunderscore}{\isacharunderscore}}, respectively.

  These special versions provide copies of the basic name space, apart
  from anything that normally appears in the user text.  For example,
  system generated variables in Isar proof contexts are usually marked
  as internal, which prevents mysterious name references like \isa{xaa} to appear in the text.

  \medskip Manipulating binding scopes often requires on-the-fly
  renamings.  A \emph{name context} contains a collection of already
  used names.  The \isa{declare} operation adds names to the
  context.

  The \isa{invents} operation derives a number of fresh names from
  a given starting point.  For example, the first three names derived
  from \isa{a} are \isa{a}, \isa{b}, \isa{c}.

  The \isa{variants} operation produces fresh names by
  incrementing tentative names as base-26 numbers (with digits \isa{a{\isachardot}{\isachardot}z}) until all clashes are resolved.  For example, name \isa{foo} results in variants \isa{fooa}, \isa{foob}, \isa{fooc}, \dots, \isa{fooaa}, \isa{fooab} etc.; each renaming
  step picks the next unused variant from this sequence.%
\end{isamarkuptext}%
\isamarkuptrue%
%
\isadelimmlref
%
\endisadelimmlref
%
\isatagmlref
%
\begin{isamarkuptext}%
\begin{mldecls}
  \indexml{Name.internal}\verb|Name.internal: string -> string| \\
  \indexml{Name.skolem}\verb|Name.skolem: string -> string| \\
  \end{mldecls}
  \begin{mldecls}
  \indexmltype{Name.context}\verb|type Name.context| \\
  \indexml{Name.context}\verb|Name.context: Name.context| \\
  \indexml{Name.declare}\verb|Name.declare: string -> Name.context -> Name.context| \\
  \indexml{Name.invents}\verb|Name.invents: Name.context -> string -> int -> string list| \\
  \indexml{Name.variants}\verb|Name.variants: string list -> Name.context -> string list * Name.context| \\
  \end{mldecls}

  \begin{description}

  \item \verb|Name.internal|~\isa{name} produces an internal name
  by adding one underscore.

  \item \verb|Name.skolem|~\isa{name} produces a Skolem name by
  adding two underscores.

  \item \verb|Name.context| represents the context of already used
  names; the initial value is \verb|Name.context|.

  \item \verb|Name.declare|~\isa{name} enters a used name into the
  context.

  \item \verb|Name.invents|~\isa{context\ name\ n} produces \isa{n} fresh names derived from \isa{name}.

  \item \verb|Name.variants|~\isa{names\ context} produces fresh
  varians of \isa{names}; the result is entered into the context.

  \end{description}%
\end{isamarkuptext}%
\isamarkuptrue%
%
\endisatagmlref
{\isafoldmlref}%
%
\isadelimmlref
%
\endisadelimmlref
%
\isamarkupsubsection{Indexed names%
}
\isamarkuptrue%
%
\begin{isamarkuptext}%
An \emph{indexed name} (or \isa{indexname}) is a pair of a basic
  name and a natural number.  This representation allows efficient
  renaming by incrementing the second component only.  The canonical
  way to rename two collections of indexnames apart from each other is
  this: determine the maximum index \isa{maxidx} of the first
  collection, then increment all indexes of the second collection by
  \isa{maxidx\ {\isacharplus}\ {\isadigit{1}}}; the maximum index of an empty collection is
  \isa{{\isacharminus}{\isadigit{1}}}.

  Occasionally, basic names and indexed names are injected into the
  same pair type: the (improper) indexname \isa{{\isacharparenleft}x{\isacharcomma}\ {\isacharminus}{\isadigit{1}}{\isacharparenright}} is used
  to encode basic names.

  \medskip Isabelle syntax observes the following rules for
  representing an indexname \isa{{\isacharparenleft}x{\isacharcomma}\ i{\isacharparenright}} as a packed string:

  \begin{itemize}

  \item \isa{{\isacharquery}x} if \isa{x} does not end with a digit and \isa{i\ {\isacharequal}\ {\isadigit{0}}},

  \item \isa{{\isacharquery}xi} if \isa{x} does not end with a digit,

  \item \isa{{\isacharquery}x{\isachardot}i} otherwise.

  \end{itemize}

  Indexnames may acquire large index numbers over time.  Results are
  normalized towards \isa{{\isadigit{0}}} at certain checkpoints, notably at
  the end of a proof.  This works by producing variants of the
  corresponding basic name components.  For example, the collection
  \isa{{\isacharquery}x{\isadigit{1}}{\isacharcomma}\ {\isacharquery}x{\isadigit{7}}{\isacharcomma}\ {\isacharquery}x{\isadigit{4}}{\isadigit{2}}} becomes \isa{{\isacharquery}x{\isacharcomma}\ {\isacharquery}xa{\isacharcomma}\ {\isacharquery}xb}.%
\end{isamarkuptext}%
\isamarkuptrue%
%
\isadelimmlref
%
\endisadelimmlref
%
\isatagmlref
%
\begin{isamarkuptext}%
\begin{mldecls}
  \indexmltype{indexname}\verb|type indexname| \\
  \end{mldecls}

  \begin{description}

  \item \verb|indexname| represents indexed names.  This is an
  abbreviation for \verb|string * int|.  The second component is
  usually non-negative, except for situations where \isa{{\isacharparenleft}x{\isacharcomma}\ {\isacharminus}{\isadigit{1}}{\isacharparenright}}
  is used to embed basic names into this type.

  \end{description}%
\end{isamarkuptext}%
\isamarkuptrue%
%
\endisatagmlref
{\isafoldmlref}%
%
\isadelimmlref
%
\endisadelimmlref
%
\isamarkupsubsection{Qualified names and name spaces%
}
\isamarkuptrue%
%
\begin{isamarkuptext}%
A \emph{qualified name} consists of a non-empty sequence of basic
  name components.  The packed representation uses a dot as separator,
  as in ``\isa{A{\isachardot}b{\isachardot}c}''.  The last component is called \emph{base}
  name, the remaining prefix \emph{qualifier} (which may be empty).
  The idea of qualified names is to encode nested structures by
  recording the access paths as qualifiers.  For example, an item
  named ``\isa{A{\isachardot}b{\isachardot}c}'' may be understood as a local entity \isa{c}, within a local structure \isa{b}, within a global
  structure \isa{A}.  Typically, name space hierarchies consist of
  1--2 levels of qualification, but this need not be always so.

  The empty name is commonly used as an indication of unnamed
  entities, whenever this makes any sense.  The basic operations on
  qualified names are smart enough to pass through such improper names
  unchanged.

  \medskip A \isa{naming} policy tells how to turn a name
  specification into a fully qualified internal name (by the \isa{full} operation), and how fully qualified names may be accessed
  externally.  For example, the default naming policy is to prefix an
  implicit path: \isa{full\ x} produces \isa{path{\isachardot}x}, and the
  standard accesses for \isa{path{\isachardot}x} include both \isa{x} and
  \isa{path{\isachardot}x}.  Normally, the naming is implicit in the theory or
  proof context; there are separate versions of the corresponding.

  \medskip A \isa{name\ space} manages a collection of fully
  internalized names, together with a mapping between external names
  and internal names (in both directions).  The corresponding \isa{intern} and \isa{extern} operations are mostly used for
  parsing and printing only!  The \isa{declare} operation augments
  a name space according to the accesses determined by the naming
  policy.

  \medskip As a general principle, there is a separate name space for
  each kind of formal entity, e.g.\ logical constant, type
  constructor, type class, theorem.  It is usually clear from the
  occurrence in concrete syntax (or from the scope) which kind of
  entity a name refers to.  For example, the very same name \isa{c} may be used uniformly for a constant, type constructor, and
  type class.

  There are common schemes to name theorems systematically, according
  to the name of the main logical entity involved, e.g.\ \isa{c{\isachardot}intro} for a canonical theorem related to constant \isa{c}.
  This technique of mapping names from one space into another requires
  some care in order to avoid conflicts.  In particular, theorem names
  derived from a type constructor or type class are better suffixed in
  addition to the usual qualification, e.g.\ \isa{c{\isacharunderscore}type{\isachardot}intro}
  and \isa{c{\isacharunderscore}class{\isachardot}intro} for theorems related to type \isa{c}
  and class \isa{c}, respectively.%
\end{isamarkuptext}%
\isamarkuptrue%
%
\isadelimmlref
%
\endisadelimmlref
%
\isatagmlref
%
\begin{isamarkuptext}%
\begin{mldecls}
  \indexml{NameSpace.base}\verb|NameSpace.base: string -> string| \\
  \indexml{NameSpace.qualifier}\verb|NameSpace.qualifier: string -> string| \\
  \indexml{NameSpace.append}\verb|NameSpace.append: string -> string -> string| \\
  \indexml{NameSpace.implode}\verb|NameSpace.implode: string list -> string| \\
  \indexml{NameSpace.explode}\verb|NameSpace.explode: string -> string list| \\
  \end{mldecls}
  \begin{mldecls}
  \indexmltype{NameSpace.naming}\verb|type NameSpace.naming| \\
  \indexml{NameSpace.default\_naming}\verb|NameSpace.default_naming: NameSpace.naming| \\
  \indexml{NameSpace.add\_path}\verb|NameSpace.add_path: string -> NameSpace.naming -> NameSpace.naming| \\
  \indexml{NameSpace.full}\verb|NameSpace.full: NameSpace.naming -> string -> string| \\
  \end{mldecls}
  \begin{mldecls}
  \indexmltype{NameSpace.T}\verb|type NameSpace.T| \\
  \indexml{NameSpace.empty}\verb|NameSpace.empty: NameSpace.T| \\
  \indexml{NameSpace.merge}\verb|NameSpace.merge: NameSpace.T * NameSpace.T -> NameSpace.T| \\
  \indexml{NameSpace.declare}\verb|NameSpace.declare: NameSpace.naming -> string -> NameSpace.T -> NameSpace.T| \\
  \indexml{NameSpace.intern}\verb|NameSpace.intern: NameSpace.T -> string -> string| \\
  \indexml{NameSpace.extern}\verb|NameSpace.extern: NameSpace.T -> string -> string| \\
  \end{mldecls}

  \begin{description}

  \item \verb|NameSpace.base|~\isa{name} returns the base name of a
  qualified name.

  \item \verb|NameSpace.qualifier|~\isa{name} returns the qualifier
  of a qualified name.

  \item \verb|NameSpace.append|~\isa{name\isactrlisub {\isadigit{1}}\ name\isactrlisub {\isadigit{2}}}
  appends two qualified names.

  \item \verb|NameSpace.implode|~\isa{name} and \verb|NameSpace.explode|~\isa{names} convert between the packed string
  representation and the explicit list form of qualified names.

  \item \verb|NameSpace.naming| represents the abstract concept of
  a naming policy.

  \item \verb|NameSpace.default_naming| is the default naming policy.
  In a theory context, this is usually augmented by a path prefix
  consisting of the theory name.

  \item \verb|NameSpace.add_path|~\isa{path\ naming} augments the
  naming policy by extending its path component.

  \item \verb|NameSpace.full|\isa{naming\ name} turns a name
  specification (usually a basic name) into the fully qualified
  internal version, according to the given naming policy.

  \item \verb|NameSpace.T| represents name spaces.

  \item \verb|NameSpace.empty| and \verb|NameSpace.merge|~\isa{{\isacharparenleft}space\isactrlisub {\isadigit{1}}{\isacharcomma}\ space\isactrlisub {\isadigit{2}}{\isacharparenright}} are the canonical operations for
  maintaining name spaces according to theory data management
  (\secref{sec:context-data}).

  \item \verb|NameSpace.declare|~\isa{naming\ name\ space} enters a
  fully qualified name into the name space, with external accesses
  determined by the naming policy.

  \item \verb|NameSpace.intern|~\isa{space\ name} internalizes a
  (partially qualified) external name.

  This operation is mostly for parsing!  Note that fully qualified
  names stemming from declarations are produced via \verb|NameSpace.full| (or its derivatives for \verb|theory| and
  \verb|Proof.context|).

  \item \verb|NameSpace.extern|~\isa{space\ name} externalizes a
  (fully qualified) internal name.

  This operation is mostly for printing!  Note unqualified names are
  produced via \verb|NameSpace.base|.

  \end{description}%
\end{isamarkuptext}%
\isamarkuptrue%
%
\endisatagmlref
{\isafoldmlref}%
%
\isadelimmlref
%
\endisadelimmlref
%
\isadelimtheory
%
\endisadelimtheory
%
\isatagtheory
\isacommand{end}\isamarkupfalse%
%
\endisatagtheory
{\isafoldtheory}%
%
\isadelimtheory
%
\endisadelimtheory
\isanewline
\end{isabellebody}%
%%% Local Variables:
%%% mode: latex
%%% TeX-master: "root"
%%% End:

%
\begin{isabellebody}%
\def\isabellecontext{logic}%
%
\isadelimtheory
\isanewline
\isanewline
\isanewline
%
\endisadelimtheory
%
\isatagtheory
\isacommand{theory}\isamarkupfalse%
\ logic\ \isakeyword{imports}\ base\ \isakeyword{begin}%
\endisatagtheory
{\isafoldtheory}%
%
\isadelimtheory
%
\endisadelimtheory
%
\isamarkupchapter{Primitive logic \label{ch:logic}%
}
\isamarkuptrue%
%
\begin{isamarkuptext}%
The logical foundations of Isabelle/Isar are that of the Pure logic,
  which has been introduced as a natural-deduction framework in
  \cite{paulson700}.  This is essentially the same logic as ``\isa{{\isasymlambda}HOL}'' in the more abstract setting of Pure Type Systems (PTS)
  \cite{Barendregt-Geuvers:2001}, although there are some key
  differences in the specific treatment of simple types in
  Isabelle/Pure.

  Following type-theoretic parlance, the Pure logic consists of three
  levels of \isa{{\isasymlambda}}-calculus with corresponding arrows, \isa{{\isasymRightarrow}} for syntactic function space (terms depending on terms), \isa{{\isasymAnd}} for universal quantification (proofs depending on terms), and
  \isa{{\isasymLongrightarrow}} for implication (proofs depending on proofs).

  Derivations are relative to a logical theory, which declares type
  constructors, constants, and axioms.  Theory declarations support
  schematic polymorphism, which is strictly speaking outside the
  logic.\footnote{This is the deeper logical reason, why the theory
  context \isa{{\isasymTheta}} is separate from the proof context \isa{{\isasymGamma}}
  of the core calculus.}%
\end{isamarkuptext}%
\isamarkuptrue%
%
\isamarkupsection{Types \label{sec:types}%
}
\isamarkuptrue%
%
\begin{isamarkuptext}%
The language of types is an uninterpreted order-sorted first-order
  algebra; types are qualified by ordered type classes.

  \medskip A \emph{type class} is an abstract syntactic entity
  declared in the theory context.  The \emph{subclass relation} \isa{c\isactrlisub {\isadigit{1}}\ {\isasymsubseteq}\ c\isactrlisub {\isadigit{2}}} is specified by stating an acyclic
  generating relation; the transitive closure is maintained
  internally.  The resulting relation is an ordering: reflexive,
  transitive, and antisymmetric.

  A \emph{sort} is a list of type classes written as \isa{s\ {\isacharequal}\ {\isacharbraceleft}c\isactrlisub {\isadigit{1}}{\isacharcomma}\ {\isasymdots}{\isacharcomma}\ c\isactrlisub m{\isacharbraceright}}, which represents symbolic
  intersection.  Notationally, the curly braces are omitted for
  singleton intersections, i.e.\ any class \isa{c} may be read as
  a sort \isa{{\isacharbraceleft}c{\isacharbraceright}}.  The ordering on type classes is extended to
  sorts according to the meaning of intersections: \isa{{\isacharbraceleft}c\isactrlisub {\isadigit{1}}{\isacharcomma}\ {\isasymdots}\ c\isactrlisub m{\isacharbraceright}\ {\isasymsubseteq}\ {\isacharbraceleft}d\isactrlisub {\isadigit{1}}{\isacharcomma}\ {\isasymdots}{\isacharcomma}\ d\isactrlisub n{\isacharbraceright}} iff
  \isa{{\isasymforall}j{\isachardot}\ {\isasymexists}i{\isachardot}\ c\isactrlisub i\ {\isasymsubseteq}\ d\isactrlisub j}.  The empty intersection
  \isa{{\isacharbraceleft}{\isacharbraceright}} refers to the universal sort, which is the largest
  element wrt.\ the sort order.  The intersections of all (finitely
  many) classes declared in the current theory are the minimal
  elements wrt.\ the sort order.

  \medskip A \emph{fixed type variable} is a pair of a basic name
  (starting with a \isa{{\isacharprime}} character) and a sort constraint, e.g.\
  \isa{{\isacharparenleft}{\isacharprime}a{\isacharcomma}\ s{\isacharparenright}} which is usually printed as \isa{{\isasymalpha}\isactrlisub s}.
  A \emph{schematic type variable} is a pair of an indexname and a
  sort constraint, e.g.\ \isa{{\isacharparenleft}{\isacharparenleft}{\isacharprime}a{\isacharcomma}\ {\isadigit{0}}{\isacharparenright}{\isacharcomma}\ s{\isacharparenright}} which is usually
  printed as \isa{{\isacharquery}{\isasymalpha}\isactrlisub s}.

  Note that \emph{all} syntactic components contribute to the identity
  of type variables, including the sort constraint.  The core logic
  handles type variables with the same name but different sorts as
  different, although some outer layers of the system make it hard to
  produce anything like this.

  A \emph{type constructor} \isa{{\isasymkappa}} is a \isa{k}-ary operator
  on types declared in the theory.  Type constructor application is
  written postfix as \isa{{\isacharparenleft}{\isasymalpha}\isactrlisub {\isadigit{1}}{\isacharcomma}\ {\isasymdots}{\isacharcomma}\ {\isasymalpha}\isactrlisub k{\isacharparenright}{\isasymkappa}}.  For
  \isa{k\ {\isacharequal}\ {\isadigit{0}}} the argument tuple is omitted, e.g.\ \isa{prop}
  instead of \isa{{\isacharparenleft}{\isacharparenright}prop}.  For \isa{k\ {\isacharequal}\ {\isadigit{1}}} the parentheses
  are omitted, e.g.\ \isa{{\isasymalpha}\ list} instead of \isa{{\isacharparenleft}{\isasymalpha}{\isacharparenright}list}.
  Further notation is provided for specific constructors, notably the
  right-associative infix \isa{{\isasymalpha}\ {\isasymRightarrow}\ {\isasymbeta}} instead of \isa{{\isacharparenleft}{\isasymalpha}{\isacharcomma}\ {\isasymbeta}{\isacharparenright}fun}.
  
  A \emph{type} is defined inductively over type variables and type
  constructors as follows: \isa{{\isasymtau}\ {\isacharequal}\ {\isasymalpha}\isactrlisub s\ {\isacharbar}\ {\isacharquery}{\isasymalpha}\isactrlisub s\ {\isacharbar}\ {\isacharparenleft}{\isasymtau}\isactrlsub {\isadigit{1}}{\isacharcomma}\ {\isasymdots}{\isacharcomma}\ {\isasymtau}\isactrlsub k{\isacharparenright}{\isasymkappa}}.

  A \emph{type abbreviation} is a syntactic definition \isa{{\isacharparenleft}\isactrlvec {\isasymalpha}{\isacharparenright}{\isasymkappa}\ {\isacharequal}\ {\isasymtau}} of an arbitrary type expression \isa{{\isasymtau}} over
  variables \isa{\isactrlvec {\isasymalpha}}.  Type abbreviations appear as type
  constructors in the syntax, but are expanded before entering the
  logical core.

  A \emph{type arity} declares the image behavior of a type
  constructor wrt.\ the algebra of sorts: \isa{{\isasymkappa}\ {\isacharcolon}{\isacharcolon}\ {\isacharparenleft}s\isactrlisub {\isadigit{1}}{\isacharcomma}\ {\isasymdots}{\isacharcomma}\ s\isactrlisub k{\isacharparenright}s} means that \isa{{\isacharparenleft}{\isasymtau}\isactrlisub {\isadigit{1}}{\isacharcomma}\ {\isasymdots}{\isacharcomma}\ {\isasymtau}\isactrlisub k{\isacharparenright}{\isasymkappa}} is
  of sort \isa{s} if every argument type \isa{{\isasymtau}\isactrlisub i} is
  of sort \isa{s\isactrlisub i}.  Arity declarations are implicitly
  completed, i.e.\ \isa{{\isasymkappa}\ {\isacharcolon}{\isacharcolon}\ {\isacharparenleft}\isactrlvec s{\isacharparenright}c} entails \isa{{\isasymkappa}\ {\isacharcolon}{\isacharcolon}\ {\isacharparenleft}\isactrlvec s{\isacharparenright}c{\isacharprime}} for any \isa{c{\isacharprime}\ {\isasymsupseteq}\ c}.

  \medskip The sort algebra is always maintained as \emph{coregular},
  which means that type arities are consistent with the subclass
  relation: for any type constructor \isa{{\isasymkappa}}, and classes \isa{c\isactrlisub {\isadigit{1}}\ {\isasymsubseteq}\ c\isactrlisub {\isadigit{2}}}, and arities \isa{{\isasymkappa}\ {\isacharcolon}{\isacharcolon}\ {\isacharparenleft}\isactrlvec s\isactrlisub {\isadigit{1}}{\isacharparenright}c\isactrlisub {\isadigit{1}}} and \isa{{\isasymkappa}\ {\isacharcolon}{\isacharcolon}\ {\isacharparenleft}\isactrlvec s\isactrlisub {\isadigit{2}}{\isacharparenright}c\isactrlisub {\isadigit{2}}} holds \isa{\isactrlvec s\isactrlisub {\isadigit{1}}\ {\isasymsubseteq}\ \isactrlvec s\isactrlisub {\isadigit{2}}} component-wise.

  The key property of a coregular order-sorted algebra is that sort
  constraints can be solved in a most general fashion: for each type
  constructor \isa{{\isasymkappa}} and sort \isa{s} there is a most general
  vector of argument sorts \isa{{\isacharparenleft}s\isactrlisub {\isadigit{1}}{\isacharcomma}\ {\isasymdots}{\isacharcomma}\ s\isactrlisub k{\isacharparenright}} such
  that a type scheme \isa{{\isacharparenleft}{\isasymalpha}\isactrlbsub s\isactrlisub {\isadigit{1}}\isactrlesub {\isacharcomma}\ {\isasymdots}{\isacharcomma}\ {\isasymalpha}\isactrlbsub s\isactrlisub k\isactrlesub {\isacharparenright}{\isasymkappa}} is of sort \isa{s}.
  Consequently, type unification has most general solutions (modulo
  equivalence of sorts), so type-inference produces primary types as
  expected \cite{nipkow-prehofer}.%
\end{isamarkuptext}%
\isamarkuptrue%
%
\isadelimmlref
%
\endisadelimmlref
%
\isatagmlref
%
\begin{isamarkuptext}%
\begin{mldecls}
  \indexmltype{class}\verb|type class| \\
  \indexmltype{sort}\verb|type sort| \\
  \indexmltype{arity}\verb|type arity| \\
  \indexmltype{typ}\verb|type typ| \\
  \indexml{map\_atyps}\verb|map_atyps: (typ -> typ) -> typ -> typ| \\
  \indexml{fold\_atyps}\verb|fold_atyps: (typ -> 'a -> 'a) -> typ -> 'a -> 'a| \\
  \end{mldecls}
  \begin{mldecls}
  \indexml{Sign.subsort}\verb|Sign.subsort: theory -> sort * sort -> bool| \\
  \indexml{Sign.of\_sort}\verb|Sign.of_sort: theory -> typ * sort -> bool| \\
  \indexml{Sign.add\_types}\verb|Sign.add_types: (string * int * mixfix) list -> theory -> theory| \\
  \indexml{Sign.add\_tyabbrs\_i}\verb|Sign.add_tyabbrs_i: |\isasep\isanewline%
\verb|  (string * string list * typ * mixfix) list -> theory -> theory| \\
  \indexml{Sign.primitive\_class}\verb|Sign.primitive_class: string * class list -> theory -> theory| \\
  \indexml{Sign.primitive\_classrel}\verb|Sign.primitive_classrel: class * class -> theory -> theory| \\
  \indexml{Sign.primitive\_arity}\verb|Sign.primitive_arity: arity -> theory -> theory| \\
  \end{mldecls}

  \begin{description}

  \item \verb|class| represents type classes; this is an alias for
  \verb|string|.

  \item \verb|sort| represents sorts; this is an alias for
  \verb|class list|.

  \item \verb|arity| represents type arities; this is an alias for
  triples of the form \isa{{\isacharparenleft}{\isasymkappa}{\isacharcomma}\ \isactrlvec s{\isacharcomma}\ s{\isacharparenright}} for \isa{{\isasymkappa}\ {\isacharcolon}{\isacharcolon}\ {\isacharparenleft}\isactrlvec s{\isacharparenright}s} described above.

  \item \verb|typ| represents types; this is a datatype with
  constructors \verb|TFree|, \verb|TVar|, \verb|Type|.

  \item \verb|map_atyps|~\isa{f\ {\isasymtau}} applies the mapping \isa{f}
  to all atomic types (\verb|TFree|, \verb|TVar|) occurring in \isa{{\isasymtau}}.

  \item \verb|fold_atyps|~\isa{f\ {\isasymtau}} iterates the operation \isa{f} over all occurrences of atomic types (\verb|TFree|, \verb|TVar|)
  in \isa{{\isasymtau}}; the type structure is traversed from left to right.

  \item \verb|Sign.subsort|~\isa{thy\ {\isacharparenleft}s\isactrlisub {\isadigit{1}}{\isacharcomma}\ s\isactrlisub {\isadigit{2}}{\isacharparenright}}
  tests the subsort relation \isa{s\isactrlisub {\isadigit{1}}\ {\isasymsubseteq}\ s\isactrlisub {\isadigit{2}}}.

  \item \verb|Sign.of_sort|~\isa{thy\ {\isacharparenleft}{\isasymtau}{\isacharcomma}\ s{\isacharparenright}} tests whether type
  \isa{{\isasymtau}} is of sort \isa{s}.

  \item \verb|Sign.add_types|~\isa{{\isacharbrackleft}{\isacharparenleft}{\isasymkappa}{\isacharcomma}\ k{\isacharcomma}\ mx{\isacharparenright}{\isacharcomma}\ {\isasymdots}{\isacharbrackright}} declares a new
  type constructors \isa{{\isasymkappa}} with \isa{k} arguments and
  optional mixfix syntax.

  \item \verb|Sign.add_tyabbrs_i|~\isa{{\isacharbrackleft}{\isacharparenleft}{\isasymkappa}{\isacharcomma}\ \isactrlvec {\isasymalpha}{\isacharcomma}\ {\isasymtau}{\isacharcomma}\ mx{\isacharparenright}{\isacharcomma}\ {\isasymdots}{\isacharbrackright}}
  defines a new type abbreviation \isa{{\isacharparenleft}\isactrlvec {\isasymalpha}{\isacharparenright}{\isasymkappa}\ {\isacharequal}\ {\isasymtau}} with
  optional mixfix syntax.

  \item \verb|Sign.primitive_class|~\isa{{\isacharparenleft}c{\isacharcomma}\ {\isacharbrackleft}c\isactrlisub {\isadigit{1}}{\isacharcomma}\ {\isasymdots}{\isacharcomma}\ c\isactrlisub n{\isacharbrackright}{\isacharparenright}} declares a new class \isa{c}, together with class
  relations \isa{c\ {\isasymsubseteq}\ c\isactrlisub i}, for \isa{i\ {\isacharequal}\ {\isadigit{1}}{\isacharcomma}\ {\isasymdots}{\isacharcomma}\ n}.

  \item \verb|Sign.primitive_classrel|~\isa{{\isacharparenleft}c\isactrlisub {\isadigit{1}}{\isacharcomma}\ c\isactrlisub {\isadigit{2}}{\isacharparenright}} declares the class relation \isa{c\isactrlisub {\isadigit{1}}\ {\isasymsubseteq}\ c\isactrlisub {\isadigit{2}}}.

  \item \verb|Sign.primitive_arity|~\isa{{\isacharparenleft}{\isasymkappa}{\isacharcomma}\ \isactrlvec s{\isacharcomma}\ s{\isacharparenright}} declares
  the arity \isa{{\isasymkappa}\ {\isacharcolon}{\isacharcolon}\ {\isacharparenleft}\isactrlvec s{\isacharparenright}s}.

  \end{description}%
\end{isamarkuptext}%
\isamarkuptrue%
%
\endisatagmlref
{\isafoldmlref}%
%
\isadelimmlref
%
\endisadelimmlref
%
\isamarkupsection{Terms \label{sec:terms}%
}
\isamarkuptrue%
%
\begin{isamarkuptext}%
\glossary{Term}{FIXME}

  The language of terms is that of simply-typed \isa{{\isasymlambda}}-calculus
  with de-Bruijn indices for bound variables (cf.\ \cite{debruijn72}
  or \cite{paulson-ml2}), with the types being determined determined
  by the corresponding binders.  In contrast, free variables and
  constants are have an explicit name and type in each occurrence.

  \medskip A \emph{bound variable} is a natural number \isa{b},
  which accounts for the number of intermediate binders between the
  variable occurrence in the body and its binding position.  For
  example, the de-Bruijn term \isa{{\isasymlambda}\isactrlbsub nat\isactrlesub {\isachardot}\ {\isasymlambda}\isactrlbsub nat\isactrlesub {\isachardot}\ {\isadigit{1}}\ {\isacharplus}\ {\isadigit{0}}} would
  correspond to \isa{{\isasymlambda}x\isactrlbsub nat\isactrlesub {\isachardot}\ {\isasymlambda}y\isactrlbsub nat\isactrlesub {\isachardot}\ x\ {\isacharplus}\ y} in a named
  representation.  Note that a bound variable may be represented by
  different de-Bruijn indices at different occurrences, depending on
  the nesting of abstractions.

  A \emph{loose variable} is a bound variable that is outside the
  scope of local binders.  The types (and names) for loose variables
  can be managed as a separate context, that is maintained as a stack
  of hypothetical binders.  The core logic operates on closed terms,
  without any loose variables.

  A \emph{fixed variable} is a pair of a basic name and a type, e.g.\
  \isa{{\isacharparenleft}x{\isacharcomma}\ {\isasymtau}{\isacharparenright}} which is usually printed \isa{x\isactrlisub {\isasymtau}}.  A
  \emph{schematic variable} is a pair of an indexname and a type,
  e.g.\ \isa{{\isacharparenleft}{\isacharparenleft}x{\isacharcomma}\ {\isadigit{0}}{\isacharparenright}{\isacharcomma}\ {\isasymtau}{\isacharparenright}} which is usually printed as \isa{{\isacharquery}x\isactrlisub {\isasymtau}}.

  \medskip A \emph{constant} is a pair of a basic name and a type,
  e.g.\ \isa{{\isacharparenleft}c{\isacharcomma}\ {\isasymtau}{\isacharparenright}} which is usually printed as \isa{c\isactrlisub {\isasymtau}}.  Constants are declared in the context as polymorphic
  families \isa{c\ {\isacharcolon}{\isacharcolon}\ {\isasymsigma}}, meaning that all substitution instances
  \isa{c\isactrlisub {\isasymtau}} for \isa{{\isasymtau}\ {\isacharequal}\ {\isasymsigma}{\isasymvartheta}} are valid.

  The vector of \emph{type arguments} of constant \isa{c\isactrlisub {\isasymtau}}
  wrt.\ the declaration \isa{c\ {\isacharcolon}{\isacharcolon}\ {\isasymsigma}} is defined as the codomain of
  the matcher \isa{{\isasymvartheta}\ {\isacharequal}\ {\isacharbraceleft}{\isacharquery}{\isasymalpha}\isactrlisub {\isadigit{1}}\ {\isasymmapsto}\ {\isasymtau}\isactrlisub {\isadigit{1}}{\isacharcomma}\ {\isasymdots}{\isacharcomma}\ {\isacharquery}{\isasymalpha}\isactrlisub n\ {\isasymmapsto}\ {\isasymtau}\isactrlisub n{\isacharbraceright}} presented in canonical order \isa{{\isacharparenleft}{\isasymtau}\isactrlisub {\isadigit{1}}{\isacharcomma}\ {\isasymdots}{\isacharcomma}\ {\isasymtau}\isactrlisub n{\isacharparenright}}.  Within a given theory context,
  there is a one-to-one correspondence between any constant \isa{c\isactrlisub {\isasymtau}} and the application \isa{c{\isacharparenleft}{\isasymtau}\isactrlisub {\isadigit{1}}{\isacharcomma}\ {\isasymdots}{\isacharcomma}\ {\isasymtau}\isactrlisub n{\isacharparenright}} of its type arguments.  For example, with \isa{plus\ {\isacharcolon}{\isacharcolon}\ {\isasymalpha}\ {\isasymRightarrow}\ {\isasymalpha}\ {\isasymRightarrow}\ {\isasymalpha}}, the instance \isa{plus\isactrlbsub nat\ {\isasymRightarrow}\ nat\ {\isasymRightarrow}\ nat\isactrlesub } corresponds to \isa{plus{\isacharparenleft}nat{\isacharparenright}}.

  Constant declarations \isa{c\ {\isacharcolon}{\isacharcolon}\ {\isasymsigma}} may contain sort constraints
  for type variables in \isa{{\isasymsigma}}.  These are observed by
  type-inference as expected, but \emph{ignored} by the core logic.
  This means the primitive logic is able to reason with instances of
  polymorphic constants that the user-level type-checker would reject
  due to violation of type class restrictions.

  \medskip An \emph{atomic} term is either a variable or constant.  A
  \emph{term} is defined inductively over atomic terms, with
  abstraction and application as follows: \isa{t\ {\isacharequal}\ b\ {\isacharbar}\ x\isactrlisub {\isasymtau}\ {\isacharbar}\ {\isacharquery}x\isactrlisub {\isasymtau}\ {\isacharbar}\ c\isactrlisub {\isasymtau}\ {\isacharbar}\ {\isasymlambda}\isactrlisub {\isasymtau}{\isachardot}\ t\ {\isacharbar}\ t\isactrlisub {\isadigit{1}}\ t\isactrlisub {\isadigit{2}}}.
  Parsing and printing takes care of converting between an external
  representation with named bound variables.  Subsequently, we shall
  use the latter notation instead of internal de-Bruijn
  representation.

  The inductive relation \isa{t\ {\isacharcolon}{\isacharcolon}\ {\isasymtau}} assigns a (unique) type to a
  term according to the structure of atomic terms, abstractions, and
  applicatins:
  \[
  \infer{\isa{a\isactrlisub {\isasymtau}\ {\isacharcolon}{\isacharcolon}\ {\isasymtau}}}{}
  \qquad
  \infer{\isa{{\isacharparenleft}{\isasymlambda}x\isactrlsub {\isasymtau}{\isachardot}\ t{\isacharparenright}\ {\isacharcolon}{\isacharcolon}\ {\isasymtau}\ {\isasymRightarrow}\ {\isasymsigma}}}{\isa{t\ {\isacharcolon}{\isacharcolon}\ {\isasymsigma}}}
  \qquad
  \infer{\isa{t\ u\ {\isacharcolon}{\isacharcolon}\ {\isasymsigma}}}{\isa{t\ {\isacharcolon}{\isacharcolon}\ {\isasymtau}\ {\isasymRightarrow}\ {\isasymsigma}} & \isa{u\ {\isacharcolon}{\isacharcolon}\ {\isasymtau}}}
  \]
  A \emph{well-typed term} is a term that can be typed according to these rules.

  Typing information can be omitted: type-inference is able to
  reconstruct the most general type of a raw term, while assigning
  most general types to all of its variables and constants.
  Type-inference depends on a context of type constraints for fixed
  variables, and declarations for polymorphic constants.

  The identity of atomic terms consists both of the name and the type
  component.  This means that different variables \isa{x\isactrlbsub {\isasymtau}\isactrlisub {\isadigit{1}}\isactrlesub } and \isa{x\isactrlbsub {\isasymtau}\isactrlisub {\isadigit{2}}\isactrlesub } may become the same after type
  instantiation.  Some outer layers of the system make it hard to
  produce variables of the same name, but different types.  In
  contrast, mixed instances of polymorphic constants occur frequently.

  \medskip The \emph{hidden polymorphism} of a term \isa{t\ {\isacharcolon}{\isacharcolon}\ {\isasymsigma}}
  is the set of type variables occurring in \isa{t}, but not in
  \isa{{\isasymsigma}}.  This means that the term implicitly depends on type
  arguments that are not accounted in the result type, i.e.\ there are
  different type instances \isa{t{\isasymvartheta}\ {\isacharcolon}{\isacharcolon}\ {\isasymsigma}} and \isa{t{\isasymvartheta}{\isacharprime}\ {\isacharcolon}{\isacharcolon}\ {\isasymsigma}} with the same type.  This slightly
  pathological situation notoriously demands additional care.

  \medskip A \emph{term abbreviation} is a syntactic definition \isa{c\isactrlisub {\isasymsigma}\ {\isasymequiv}\ t} of a closed term \isa{t} of type \isa{{\isasymsigma}},
  without any hidden polymorphism.  A term abbreviation looks like a
  constant in the syntax, but is expanded before entering the logical
  core.  Abbreviations are usually reverted when printing terms, using
  \isa{t\ {\isasymrightarrow}\ c\isactrlisub {\isasymsigma}} as rules for higher-order rewriting.

  \medskip Canonical operations on \isa{{\isasymlambda}}-terms include \isa{{\isasymalpha}{\isasymbeta}{\isasymeta}}-conversion: \isa{{\isasymalpha}}-conversion refers to capture-free
  renaming of bound variables; \isa{{\isasymbeta}}-conversion contracts an
  abstraction applied to an argument term, substituting the argument
  in the body: \isa{{\isacharparenleft}{\isasymlambda}x{\isachardot}\ b{\isacharparenright}a} becomes \isa{b{\isacharbrackleft}a{\isacharslash}x{\isacharbrackright}}; \isa{{\isasymeta}}-conversion contracts vacuous application-abstraction: \isa{{\isasymlambda}x{\isachardot}\ f\ x} becomes \isa{f}, provided that the bound variable
  does not occur in \isa{f}.

  Terms are normally treated modulo \isa{{\isasymalpha}}-conversion, which is
  implicit in the de-Bruijn representation.  Names for bound variables
  in abstractions are maintained separately as (meaningless) comments,
  mostly for parsing and printing.  Full \isa{{\isasymalpha}{\isasymbeta}{\isasymeta}}-conversion is
  commonplace in various standard operations (\secref{sec:rules}) that
  are based on higher-order unification and matching.%
\end{isamarkuptext}%
\isamarkuptrue%
%
\isadelimmlref
%
\endisadelimmlref
%
\isatagmlref
%
\begin{isamarkuptext}%
\begin{mldecls}
  \indexmltype{term}\verb|type term| \\
  \indexml{op aconv}\verb|op aconv: term * term -> bool| \\
  \indexml{map\_types}\verb|map_types: (typ -> typ) -> term -> term| \\
  \indexml{fold\_types}\verb|fold_types: (typ -> 'a -> 'a) -> term -> 'a -> 'a| \\
  \indexml{map\_aterms}\verb|map_aterms: (term -> term) -> term -> term| \\
  \indexml{fold\_aterms}\verb|fold_aterms: (term -> 'a -> 'a) -> term -> 'a -> 'a| \\
  \end{mldecls}
  \begin{mldecls}
  \indexml{fastype\_of}\verb|fastype_of: term -> typ| \\
  \indexml{lambda}\verb|lambda: term -> term -> term| \\
  \indexml{betapply}\verb|betapply: term * term -> term| \\
  \indexml{Sign.declare\_const}\verb|Sign.declare_const: Properties.T -> (Name.binding * typ) * mixfix ->|\isasep\isanewline%
\verb|  theory -> term * theory| \\
  \indexml{Sign.add\_abbrev}\verb|Sign.add_abbrev: string -> Properties.T -> bstring * term ->|\isasep\isanewline%
\verb|  theory -> (term * term) * theory| \\
  \indexml{Sign.const\_typargs}\verb|Sign.const_typargs: theory -> string * typ -> typ list| \\
  \indexml{Sign.const\_instance}\verb|Sign.const_instance: theory -> string * typ list -> typ| \\
  \end{mldecls}

  \begin{description}

  \item \verb|term| represents de-Bruijn terms, with comments in
  abstractions, and explicitly named free variables and constants;
  this is a datatype with constructors \verb|Bound|, \verb|Free|, \verb|Var|, \verb|Const|, \verb|Abs|, \verb|op $|.

  \item \isa{t}~\verb|aconv|~\isa{u} checks \isa{{\isasymalpha}}-equivalence of two terms.  This is the basic equality relation
  on type \verb|term|; raw datatype equality should only be used
  for operations related to parsing or printing!

  \item \verb|map_types|~\isa{f\ t} applies the mapping \isa{f} to all types occurring in \isa{t}.

  \item \verb|fold_types|~\isa{f\ t} iterates the operation \isa{f} over all occurrences of types in \isa{t}; the term
  structure is traversed from left to right.

  \item \verb|map_aterms|~\isa{f\ t} applies the mapping \isa{f}
  to all atomic terms (\verb|Bound|, \verb|Free|, \verb|Var|, \verb|Const|) occurring in \isa{t}.

  \item \verb|fold_aterms|~\isa{f\ t} iterates the operation \isa{f} over all occurrences of atomic terms (\verb|Bound|, \verb|Free|,
  \verb|Var|, \verb|Const|) in \isa{t}; the term structure is
  traversed from left to right.

  \item \verb|fastype_of|~\isa{t} determines the type of a
  well-typed term.  This operation is relatively slow, despite the
  omission of any sanity checks.

  \item \verb|lambda|~\isa{a\ b} produces an abstraction \isa{{\isasymlambda}a{\isachardot}\ b}, where occurrences of the atomic term \isa{a} in the
  body \isa{b} are replaced by bound variables.

  \item \verb|betapply|~\isa{{\isacharparenleft}t{\isacharcomma}\ u{\isacharparenright}} produces an application \isa{t\ u}, with topmost \isa{{\isasymbeta}}-conversion if \isa{t} is an
  abstraction.

  \item \verb|Sign.declare_const|~\isa{properties\ {\isacharparenleft}{\isacharparenleft}c{\isacharcomma}\ {\isasymsigma}{\isacharparenright}{\isacharcomma}\ mx{\isacharparenright}}
  declares a new constant \isa{c\ {\isacharcolon}{\isacharcolon}\ {\isasymsigma}} with optional mixfix
  syntax.

  \item \verb|Sign.add_abbrev|~\isa{print{\isacharunderscore}mode\ properties\ {\isacharparenleft}c{\isacharcomma}\ t{\isacharparenright}}
  introduces a new term abbreviation \isa{c\ {\isasymequiv}\ t}.

  \item \verb|Sign.const_typargs|~\isa{thy\ {\isacharparenleft}c{\isacharcomma}\ {\isasymtau}{\isacharparenright}} and \verb|Sign.const_instance|~\isa{thy\ {\isacharparenleft}c{\isacharcomma}\ {\isacharbrackleft}{\isasymtau}\isactrlisub {\isadigit{1}}{\isacharcomma}\ {\isasymdots}{\isacharcomma}\ {\isasymtau}\isactrlisub n{\isacharbrackright}{\isacharparenright}}
  convert between two representations of polymorphic constants: full
  type instance vs.\ compact type arguments form.

  \end{description}%
\end{isamarkuptext}%
\isamarkuptrue%
%
\endisatagmlref
{\isafoldmlref}%
%
\isadelimmlref
%
\endisadelimmlref
%
\isamarkupsection{Theorems \label{sec:thms}%
}
\isamarkuptrue%
%
\begin{isamarkuptext}%
\glossary{Proposition}{FIXME A \seeglossary{term} of
  \seeglossary{type} \isa{prop}.  Internally, there is nothing
  special about propositions apart from their type, but the concrete
  syntax enforces a clear distinction.  Propositions are structured
  via implication \isa{A\ {\isasymLongrightarrow}\ B} or universal quantification \isa{{\isasymAnd}x{\isachardot}\ B\ x} --- anything else is considered atomic.  The canonical
  form for propositions is that of a \seeglossary{Hereditary Harrop
  Formula}. FIXME}

  \glossary{Theorem}{A proven proposition within a certain theory and
  proof context, formally \isa{{\isasymGamma}\ {\isasymturnstile}\isactrlsub {\isasymTheta}\ {\isasymphi}}; both contexts are
  rarely spelled out explicitly.  Theorems are usually normalized
  according to the \seeglossary{HHF} format. FIXME}

  \glossary{Fact}{Sometimes used interchangeably for
  \seeglossary{theorem}.  Strictly speaking, a list of theorems,
  essentially an extra-logical conjunction.  Facts emerge either as
  local assumptions, or as results of local goal statements --- both
  may be simultaneous, hence the list representation. FIXME}

  \glossary{Schematic variable}{FIXME}

  \glossary{Fixed variable}{A variable that is bound within a certain
  proof context; an arbitrary-but-fixed entity within a portion of
  proof text. FIXME}

  \glossary{Free variable}{Synonymous for \seeglossary{fixed
  variable}. FIXME}

  \glossary{Bound variable}{FIXME}

  \glossary{Variable}{See \seeglossary{schematic variable},
  \seeglossary{fixed variable}, \seeglossary{bound variable}, or
  \seeglossary{type variable}.  The distinguishing feature of
  different variables is their binding scope. FIXME}

  A \emph{proposition} is a well-typed term of type \isa{prop}, a
  \emph{theorem} is a proven proposition (depending on a context of
  hypotheses and the background theory).  Primitive inferences include
  plain natural deduction rules for the primary connectives \isa{{\isasymAnd}} and \isa{{\isasymLongrightarrow}} of the framework.  There is also a builtin
  notion of equality/equivalence \isa{{\isasymequiv}}.%
\end{isamarkuptext}%
\isamarkuptrue%
%
\isamarkupsubsection{Primitive connectives and rules \label{sec:prim-rules}%
}
\isamarkuptrue%
%
\begin{isamarkuptext}%
The theory \isa{Pure} contains constant declarations for the
  primitive connectives \isa{{\isasymAnd}}, \isa{{\isasymLongrightarrow}}, and \isa{{\isasymequiv}} of
  the logical framework, see \figref{fig:pure-connectives}.  The
  derivability judgment \isa{A\isactrlisub {\isadigit{1}}{\isacharcomma}\ {\isasymdots}{\isacharcomma}\ A\isactrlisub n\ {\isasymturnstile}\ B} is
  defined inductively by the primitive inferences given in
  \figref{fig:prim-rules}, with the global restriction that the
  hypotheses must \emph{not} contain any schematic variables.  The
  builtin equality is conceptually axiomatized as shown in
  \figref{fig:pure-equality}, although the implementation works
  directly with derived inferences.

  \begin{figure}[htb]
  \begin{center}
  \begin{tabular}{ll}
  \isa{all\ {\isacharcolon}{\isacharcolon}\ {\isacharparenleft}{\isasymalpha}\ {\isasymRightarrow}\ prop{\isacharparenright}\ {\isasymRightarrow}\ prop} & universal quantification (binder \isa{{\isasymAnd}}) \\
  \isa{{\isasymLongrightarrow}\ {\isacharcolon}{\isacharcolon}\ prop\ {\isasymRightarrow}\ prop\ {\isasymRightarrow}\ prop} & implication (right associative infix) \\
  \isa{{\isasymequiv}\ {\isacharcolon}{\isacharcolon}\ {\isasymalpha}\ {\isasymRightarrow}\ {\isasymalpha}\ {\isasymRightarrow}\ prop} & equality relation (infix) \\
  \end{tabular}
  \caption{Primitive connectives of Pure}\label{fig:pure-connectives}
  \end{center}
  \end{figure}

  \begin{figure}[htb]
  \begin{center}
  \[
  \infer[\isa{{\isacharparenleft}axiom{\isacharparenright}}]{\isa{{\isasymturnstile}\ A}}{\isa{A\ {\isasymin}\ {\isasymTheta}}}
  \qquad
  \infer[\isa{{\isacharparenleft}assume{\isacharparenright}}]{\isa{A\ {\isasymturnstile}\ A}}{}
  \]
  \[
  \infer[\isa{{\isacharparenleft}{\isasymAnd}{\isacharunderscore}intro{\isacharparenright}}]{\isa{{\isasymGamma}\ {\isasymturnstile}\ {\isasymAnd}x{\isachardot}\ b{\isacharbrackleft}x{\isacharbrackright}}}{\isa{{\isasymGamma}\ {\isasymturnstile}\ b{\isacharbrackleft}x{\isacharbrackright}} & \isa{x\ {\isasymnotin}\ {\isasymGamma}}}
  \qquad
  \infer[\isa{{\isacharparenleft}{\isasymAnd}{\isacharunderscore}elim{\isacharparenright}}]{\isa{{\isasymGamma}\ {\isasymturnstile}\ b{\isacharbrackleft}a{\isacharbrackright}}}{\isa{{\isasymGamma}\ {\isasymturnstile}\ {\isasymAnd}x{\isachardot}\ b{\isacharbrackleft}x{\isacharbrackright}}}
  \]
  \[
  \infer[\isa{{\isacharparenleft}{\isasymLongrightarrow}{\isacharunderscore}intro{\isacharparenright}}]{\isa{{\isasymGamma}\ {\isacharminus}\ A\ {\isasymturnstile}\ A\ {\isasymLongrightarrow}\ B}}{\isa{{\isasymGamma}\ {\isasymturnstile}\ B}}
  \qquad
  \infer[\isa{{\isacharparenleft}{\isasymLongrightarrow}{\isacharunderscore}elim{\isacharparenright}}]{\isa{{\isasymGamma}\isactrlsub {\isadigit{1}}\ {\isasymunion}\ {\isasymGamma}\isactrlsub {\isadigit{2}}\ {\isasymturnstile}\ B}}{\isa{{\isasymGamma}\isactrlsub {\isadigit{1}}\ {\isasymturnstile}\ A\ {\isasymLongrightarrow}\ B} & \isa{{\isasymGamma}\isactrlsub {\isadigit{2}}\ {\isasymturnstile}\ A}}
  \]
  \caption{Primitive inferences of Pure}\label{fig:prim-rules}
  \end{center}
  \end{figure}

  \begin{figure}[htb]
  \begin{center}
  \begin{tabular}{ll}
  \isa{{\isasymturnstile}\ {\isacharparenleft}{\isasymlambda}x{\isachardot}\ b{\isacharbrackleft}x{\isacharbrackright}{\isacharparenright}\ a\ {\isasymequiv}\ b{\isacharbrackleft}a{\isacharbrackright}} & \isa{{\isasymbeta}}-conversion \\
  \isa{{\isasymturnstile}\ x\ {\isasymequiv}\ x} & reflexivity \\
  \isa{{\isasymturnstile}\ x\ {\isasymequiv}\ y\ {\isasymLongrightarrow}\ P\ x\ {\isasymLongrightarrow}\ P\ y} & substitution \\
  \isa{{\isasymturnstile}\ {\isacharparenleft}{\isasymAnd}x{\isachardot}\ f\ x\ {\isasymequiv}\ g\ x{\isacharparenright}\ {\isasymLongrightarrow}\ f\ {\isasymequiv}\ g} & extensionality \\
  \isa{{\isasymturnstile}\ {\isacharparenleft}A\ {\isasymLongrightarrow}\ B{\isacharparenright}\ {\isasymLongrightarrow}\ {\isacharparenleft}B\ {\isasymLongrightarrow}\ A{\isacharparenright}\ {\isasymLongrightarrow}\ A\ {\isasymequiv}\ B} & logical equivalence \\
  \end{tabular}
  \caption{Conceptual axiomatization of Pure equality}\label{fig:pure-equality}
  \end{center}
  \end{figure}

  The introduction and elimination rules for \isa{{\isasymAnd}} and \isa{{\isasymLongrightarrow}} are analogous to formation of dependently typed \isa{{\isasymlambda}}-terms representing the underlying proof objects.  Proof terms
  are irrelevant in the Pure logic, though; they cannot occur within
  propositions.  The system provides a runtime option to record
  explicit proof terms for primitive inferences.  Thus all three
  levels of \isa{{\isasymlambda}}-calculus become explicit: \isa{{\isasymRightarrow}} for
  terms, and \isa{{\isasymAnd}{\isacharslash}{\isasymLongrightarrow}} for proofs (cf.\
  \cite{Berghofer-Nipkow:2000:TPHOL}).

  Observe that locally fixed parameters (as in \isa{{\isasymAnd}{\isacharunderscore}intro}) need
  not be recorded in the hypotheses, because the simple syntactic
  types of Pure are always inhabitable.  ``Assumptions'' \isa{x\ {\isacharcolon}{\isacharcolon}\ {\isasymtau}} for type-membership are only present as long as some \isa{x\isactrlisub {\isasymtau}} occurs in the statement body.\footnote{This is the key
  difference to ``\isa{{\isasymlambda}HOL}'' in the PTS framework
  \cite{Barendregt-Geuvers:2001}, where hypotheses \isa{x\ {\isacharcolon}\ A} are
  treated uniformly for propositions and types.}

  \medskip The axiomatization of a theory is implicitly closed by
  forming all instances of type and term variables: \isa{{\isasymturnstile}\ A{\isasymvartheta}} holds for any substitution instance of an axiom
  \isa{{\isasymturnstile}\ A}.  By pushing substitutions through derivations
  inductively, we also get admissible \isa{generalize} and \isa{instance} rules as shown in \figref{fig:subst-rules}.

  \begin{figure}[htb]
  \begin{center}
  \[
  \infer{\isa{{\isasymGamma}\ {\isasymturnstile}\ B{\isacharbrackleft}{\isacharquery}{\isasymalpha}{\isacharbrackright}}}{\isa{{\isasymGamma}\ {\isasymturnstile}\ B{\isacharbrackleft}{\isasymalpha}{\isacharbrackright}} & \isa{{\isasymalpha}\ {\isasymnotin}\ {\isasymGamma}}}
  \quad
  \infer[\quad\isa{{\isacharparenleft}generalize{\isacharparenright}}]{\isa{{\isasymGamma}\ {\isasymturnstile}\ B{\isacharbrackleft}{\isacharquery}x{\isacharbrackright}}}{\isa{{\isasymGamma}\ {\isasymturnstile}\ B{\isacharbrackleft}x{\isacharbrackright}} & \isa{x\ {\isasymnotin}\ {\isasymGamma}}}
  \]
  \[
  \infer{\isa{{\isasymGamma}\ {\isasymturnstile}\ B{\isacharbrackleft}{\isasymtau}{\isacharbrackright}}}{\isa{{\isasymGamma}\ {\isasymturnstile}\ B{\isacharbrackleft}{\isacharquery}{\isasymalpha}{\isacharbrackright}}}
  \quad
  \infer[\quad\isa{{\isacharparenleft}instantiate{\isacharparenright}}]{\isa{{\isasymGamma}\ {\isasymturnstile}\ B{\isacharbrackleft}t{\isacharbrackright}}}{\isa{{\isasymGamma}\ {\isasymturnstile}\ B{\isacharbrackleft}{\isacharquery}x{\isacharbrackright}}}
  \]
  \caption{Admissible substitution rules}\label{fig:subst-rules}
  \end{center}
  \end{figure}

  Note that \isa{instantiate} does not require an explicit
  side-condition, because \isa{{\isasymGamma}} may never contain schematic
  variables.

  In principle, variables could be substituted in hypotheses as well,
  but this would disrupt the monotonicity of reasoning: deriving
  \isa{{\isasymGamma}{\isasymvartheta}\ {\isasymturnstile}\ B{\isasymvartheta}} from \isa{{\isasymGamma}\ {\isasymturnstile}\ B} is
  correct, but \isa{{\isasymGamma}{\isasymvartheta}\ {\isasymsupseteq}\ {\isasymGamma}} does not necessarily hold:
  the result belongs to a different proof context.

  \medskip An \emph{oracle} is a function that produces axioms on the
  fly.  Logically, this is an instance of the \isa{axiom} rule
  (\figref{fig:prim-rules}), but there is an operational difference.
  The system always records oracle invocations within derivations of
  theorems.  Tracing plain axioms (and named theorems) is optional.

  Axiomatizations should be limited to the bare minimum, typically as
  part of the initial logical basis of an object-logic formalization.
  Later on, theories are usually developed in a strictly definitional
  fashion, by stating only certain equalities over new constants.

  A \emph{simple definition} consists of a constant declaration \isa{c\ {\isacharcolon}{\isacharcolon}\ {\isasymsigma}} together with an axiom \isa{{\isasymturnstile}\ c\ {\isasymequiv}\ t}, where \isa{t\ {\isacharcolon}{\isacharcolon}\ {\isasymsigma}} is a closed term without any hidden polymorphism.  The RHS
  may depend on further defined constants, but not \isa{c} itself.
  Definitions of functions may be presented as \isa{c\ \isactrlvec x\ {\isasymequiv}\ t} instead of the puristic \isa{c\ {\isasymequiv}\ {\isasymlambda}\isactrlvec x{\isachardot}\ t}.

  An \emph{overloaded definition} consists of a collection of axioms
  for the same constant, with zero or one equations \isa{c{\isacharparenleft}{\isacharparenleft}\isactrlvec {\isasymalpha}{\isacharparenright}{\isasymkappa}{\isacharparenright}\ {\isasymequiv}\ t} for each type constructor \isa{{\isasymkappa}} (for
  distinct variables \isa{\isactrlvec {\isasymalpha}}).  The RHS may mention
  previously defined constants as above, or arbitrary constants \isa{d{\isacharparenleft}{\isasymalpha}\isactrlisub i{\isacharparenright}} for some \isa{{\isasymalpha}\isactrlisub i} projected from \isa{\isactrlvec {\isasymalpha}}.  Thus overloaded definitions essentially work by
  primitive recursion over the syntactic structure of a single type
  argument.%
\end{isamarkuptext}%
\isamarkuptrue%
%
\isadelimmlref
%
\endisadelimmlref
%
\isatagmlref
%
\begin{isamarkuptext}%
\begin{mldecls}
  \indexmltype{ctyp}\verb|type ctyp| \\
  \indexmltype{cterm}\verb|type cterm| \\
  \indexml{Thm.ctyp\_of}\verb|Thm.ctyp_of: theory -> typ -> ctyp| \\
  \indexml{Thm.cterm\_of}\verb|Thm.cterm_of: theory -> term -> cterm| \\
  \end{mldecls}
  \begin{mldecls}
  \indexmltype{thm}\verb|type thm| \\
  \indexml{proofs}\verb|proofs: int ref| \\
  \indexml{Thm.assume}\verb|Thm.assume: cterm -> thm| \\
  \indexml{Thm.forall\_intr}\verb|Thm.forall_intr: cterm -> thm -> thm| \\
  \indexml{Thm.forall\_elim}\verb|Thm.forall_elim: cterm -> thm -> thm| \\
  \indexml{Thm.implies\_intr}\verb|Thm.implies_intr: cterm -> thm -> thm| \\
  \indexml{Thm.implies\_elim}\verb|Thm.implies_elim: thm -> thm -> thm| \\
  \indexml{Thm.generalize}\verb|Thm.generalize: string list * string list -> int -> thm -> thm| \\
  \indexml{Thm.instantiate}\verb|Thm.instantiate: (ctyp * ctyp) list * (cterm * cterm) list -> thm -> thm| \\
  \indexml{Thm.get\_axiom\_i}\verb|Thm.get_axiom_i: theory -> string -> thm| \\
  \indexml{Thm.invoke\_oracle\_i}\verb|Thm.invoke_oracle_i: theory -> string -> theory * Object.T -> thm| \\
  \end{mldecls}
  \begin{mldecls}
  \indexml{Theory.add\_axioms\_i}\verb|Theory.add_axioms_i: (string * term) list -> theory -> theory| \\
  \indexml{Theory.add\_deps}\verb|Theory.add_deps: string -> string * typ -> (string * typ) list -> theory -> theory| \\
  \indexml{Theory.add\_oracle}\verb|Theory.add_oracle: string * (theory * Object.T -> term) -> theory -> theory| \\
  \indexml{Theory.add\_defs\_i}\verb|Theory.add_defs_i: bool -> bool -> (bstring * term) list -> theory -> theory| \\
  \end{mldecls}

  \begin{description}

  \item \verb|ctyp| and \verb|cterm| represent certified types
  and terms, respectively.  These are abstract datatypes that
  guarantee that its values have passed the full well-formedness (and
  well-typedness) checks, relative to the declarations of type
  constructors, constants etc. in the theory.

  \item \verb|ctyp_of|~\isa{thy\ {\isasymtau}} and \verb|cterm_of|~\isa{thy\ t} explicitly checks types and terms, respectively.  This also
  involves some basic normalizations, such expansion of type and term
  abbreviations from the theory context.

  Re-certification is relatively slow and should be avoided in tight
  reasoning loops.  There are separate operations to decompose
  certified entities (including actual theorems).

  \item \verb|thm| represents proven propositions.  This is an
  abstract datatype that guarantees that its values have been
  constructed by basic principles of the \verb|Thm| module.
  Every \verb|thm| value contains a sliding back-reference to the
  enclosing theory, cf.\ \secref{sec:context-theory}.

  \item \verb|proofs| determines the detail of proof recording within
  \verb|thm| values: \verb|0| records only oracles, \verb|1| records
  oracles, axioms and named theorems, \verb|2| records full proof
  terms.

  \item \verb|Thm.assume|, \verb|Thm.forall_intr|, \verb|Thm.forall_elim|, \verb|Thm.implies_intr|, and \verb|Thm.implies_elim|
  correspond to the primitive inferences of \figref{fig:prim-rules}.

  \item \verb|Thm.generalize|~\isa{{\isacharparenleft}\isactrlvec {\isasymalpha}{\isacharcomma}\ \isactrlvec x{\isacharparenright}}
  corresponds to the \isa{generalize} rules of
  \figref{fig:subst-rules}.  Here collections of type and term
  variables are generalized simultaneously, specified by the given
  basic names.

  \item \verb|Thm.instantiate|~\isa{{\isacharparenleft}\isactrlvec {\isasymalpha}\isactrlisub s{\isacharcomma}\ \isactrlvec x\isactrlisub {\isasymtau}{\isacharparenright}} corresponds to the \isa{instantiate} rules
  of \figref{fig:subst-rules}.  Type variables are substituted before
  term variables.  Note that the types in \isa{\isactrlvec x\isactrlisub {\isasymtau}}
  refer to the instantiated versions.

  \item \verb|Thm.get_axiom_i|~\isa{thy\ name} retrieves a named
  axiom, cf.\ \isa{axiom} in \figref{fig:prim-rules}.

  \item \verb|Thm.invoke_oracle_i|~\isa{thy\ name\ arg} invokes a
  named oracle function, cf.\ \isa{axiom} in
  \figref{fig:prim-rules}.

  \item \verb|Theory.add_axioms_i|~\isa{{\isacharbrackleft}{\isacharparenleft}name{\isacharcomma}\ A{\isacharparenright}{\isacharcomma}\ {\isasymdots}{\isacharbrackright}} declares
  arbitrary propositions as axioms.

  \item \verb|Theory.add_oracle|~\isa{{\isacharparenleft}name{\isacharcomma}\ f{\isacharparenright}} declares an oracle
  function for generating arbitrary axioms on the fly.

  \item \verb|Theory.add_deps|~\isa{name\ c\isactrlisub {\isasymtau}\ \isactrlvec d\isactrlisub {\isasymsigma}} declares dependencies of a named specification
  for constant \isa{c\isactrlisub {\isasymtau}}, relative to existing
  specifications for constants \isa{\isactrlvec d\isactrlisub {\isasymsigma}}.

  \item \verb|Theory.add_defs_i|~\isa{unchecked\ overloaded\ {\isacharbrackleft}{\isacharparenleft}name{\isacharcomma}\ c\ \isactrlvec x\ {\isasymequiv}\ t{\isacharparenright}{\isacharcomma}\ {\isasymdots}{\isacharbrackright}} states a definitional axiom for an existing
  constant \isa{c}.  Dependencies are recorded (cf.\ \verb|Theory.add_deps|), unless the \isa{unchecked} option is set.

  \end{description}%
\end{isamarkuptext}%
\isamarkuptrue%
%
\endisatagmlref
{\isafoldmlref}%
%
\isadelimmlref
%
\endisadelimmlref
%
\isamarkupsubsection{Auxiliary definitions%
}
\isamarkuptrue%
%
\begin{isamarkuptext}%
Theory \isa{Pure} provides a few auxiliary definitions, see
  \figref{fig:pure-aux}.  These special constants are normally not
  exposed to the user, but appear in internal encodings.

  \begin{figure}[htb]
  \begin{center}
  \begin{tabular}{ll}
  \isa{conjunction\ {\isacharcolon}{\isacharcolon}\ prop\ {\isasymRightarrow}\ prop\ {\isasymRightarrow}\ prop} & (infix \isa{{\isacharampersand}}) \\
  \isa{{\isasymturnstile}\ A\ {\isacharampersand}\ B\ {\isasymequiv}\ {\isacharparenleft}{\isasymAnd}C{\isachardot}\ {\isacharparenleft}A\ {\isasymLongrightarrow}\ B\ {\isasymLongrightarrow}\ C{\isacharparenright}\ {\isasymLongrightarrow}\ C{\isacharparenright}} \\[1ex]
  \isa{prop\ {\isacharcolon}{\isacharcolon}\ prop\ {\isasymRightarrow}\ prop} & (prefix \isa{{\isacharhash}}, suppressed) \\
  \isa{{\isacharhash}A\ {\isasymequiv}\ A} \\[1ex]
  \isa{term\ {\isacharcolon}{\isacharcolon}\ {\isasymalpha}\ {\isasymRightarrow}\ prop} & (prefix \isa{TERM}) \\
  \isa{term\ x\ {\isasymequiv}\ {\isacharparenleft}{\isasymAnd}A{\isachardot}\ A\ {\isasymLongrightarrow}\ A{\isacharparenright}} \\[1ex]
  \isa{TYPE\ {\isacharcolon}{\isacharcolon}\ {\isasymalpha}\ itself} & (prefix \isa{TYPE}) \\
  \isa{{\isacharparenleft}unspecified{\isacharparenright}} \\
  \end{tabular}
  \caption{Definitions of auxiliary connectives}\label{fig:pure-aux}
  \end{center}
  \end{figure}

  Derived conjunction rules include introduction \isa{A\ {\isasymLongrightarrow}\ B\ {\isasymLongrightarrow}\ A\ {\isacharampersand}\ B}, and destructions \isa{A\ {\isacharampersand}\ B\ {\isasymLongrightarrow}\ A} and \isa{A\ {\isacharampersand}\ B\ {\isasymLongrightarrow}\ B}.
  Conjunction allows to treat simultaneous assumptions and conclusions
  uniformly.  For example, multiple claims are intermediately
  represented as explicit conjunction, but this is refined into
  separate sub-goals before the user continues the proof; the final
  result is projected into a list of theorems (cf.\
  \secref{sec:tactical-goals}).

  The \isa{prop} marker (\isa{{\isacharhash}}) makes arbitrarily complex
  propositions appear as atomic, without changing the meaning: \isa{{\isasymGamma}\ {\isasymturnstile}\ A} and \isa{{\isasymGamma}\ {\isasymturnstile}\ {\isacharhash}A} are interchangeable.  See
  \secref{sec:tactical-goals} for specific operations.

  The \isa{term} marker turns any well-typed term into a derivable
  proposition: \isa{{\isasymturnstile}\ TERM\ t} holds unconditionally.  Although
  this is logically vacuous, it allows to treat terms and proofs
  uniformly, similar to a type-theoretic framework.

  The \isa{TYPE} constructor is the canonical representative of
  the unspecified type \isa{{\isasymalpha}\ itself}; it essentially injects the
  language of types into that of terms.  There is specific notation
  \isa{TYPE{\isacharparenleft}{\isasymtau}{\isacharparenright}} for \isa{TYPE\isactrlbsub {\isasymtau}\ itself\isactrlesub }.
  Although being devoid of any particular meaning, the \isa{TYPE{\isacharparenleft}{\isasymtau}{\isacharparenright}} accounts for the type \isa{{\isasymtau}} within the term
  language.  In particular, \isa{TYPE{\isacharparenleft}{\isasymalpha}{\isacharparenright}} may be used as formal
  argument in primitive definitions, in order to circumvent hidden
  polymorphism (cf.\ \secref{sec:terms}).  For example, \isa{c\ TYPE{\isacharparenleft}{\isasymalpha}{\isacharparenright}\ {\isasymequiv}\ A{\isacharbrackleft}{\isasymalpha}{\isacharbrackright}} defines \isa{c\ {\isacharcolon}{\isacharcolon}\ {\isasymalpha}\ itself\ {\isasymRightarrow}\ prop} in terms of
  a proposition \isa{A} that depends on an additional type
  argument, which is essentially a predicate on types.%
\end{isamarkuptext}%
\isamarkuptrue%
%
\isadelimmlref
%
\endisadelimmlref
%
\isatagmlref
%
\begin{isamarkuptext}%
\begin{mldecls}
  \indexml{Conjunction.intr}\verb|Conjunction.intr: thm -> thm -> thm| \\
  \indexml{Conjunction.elim}\verb|Conjunction.elim: thm -> thm * thm| \\
  \indexml{Drule.mk\_term}\verb|Drule.mk_term: cterm -> thm| \\
  \indexml{Drule.dest\_term}\verb|Drule.dest_term: thm -> cterm| \\
  \indexml{Logic.mk\_type}\verb|Logic.mk_type: typ -> term| \\
  \indexml{Logic.dest\_type}\verb|Logic.dest_type: term -> typ| \\
  \end{mldecls}

  \begin{description}

  \item \verb|Conjunction.intr| derives \isa{A\ {\isacharampersand}\ B} from \isa{A} and \isa{B}.

  \item \verb|Conjunction.elim| derives \isa{A} and \isa{B}
  from \isa{A\ {\isacharampersand}\ B}.

  \item \verb|Drule.mk_term| derives \isa{TERM\ t}.

  \item \verb|Drule.dest_term| recovers term \isa{t} from \isa{TERM\ t}.

  \item \verb|Logic.mk_type|~\isa{{\isasymtau}} produces the term \isa{TYPE{\isacharparenleft}{\isasymtau}{\isacharparenright}}.

  \item \verb|Logic.dest_type|~\isa{TYPE{\isacharparenleft}{\isasymtau}{\isacharparenright}} recovers the type
  \isa{{\isasymtau}}.

  \end{description}%
\end{isamarkuptext}%
\isamarkuptrue%
%
\endisatagmlref
{\isafoldmlref}%
%
\isadelimmlref
%
\endisadelimmlref
%
\isamarkupsection{Rules \label{sec:rules}%
}
\isamarkuptrue%
%
\isadelimFIXME
%
\endisadelimFIXME
%
\isatagFIXME
%
\begin{isamarkuptext}%
FIXME

  A \emph{rule} is any Pure theorem in HHF normal form; there is a
  separate calculus for rule composition, which is modeled after
  Gentzen's Natural Deduction \cite{Gentzen:1935}, but allows
  rules to be nested arbitrarily, similar to \cite{extensions91}.

  Normally, all theorems accessible to the user are proper rules.
  Low-level inferences are occasional required internally, but the
  result should be always presented in canonical form.  The higher
  interfaces of Isabelle/Isar will always produce proper rules.  It is
  important to maintain this invariant in add-on applications!

  There are two main principles of rule composition: \isa{resolution} (i.e.\ backchaining of rules) and \isa{by{\isacharminus}assumption} (i.e.\ closing a branch); both principles are
  combined in the variants of \isa{elim{\isacharminus}resolution} and \isa{dest{\isacharminus}resolution}.  Raw \isa{composition} is occasionally
  useful as well, also it is strictly speaking outside of the proper
  rule calculus.

  Rules are treated modulo general higher-order unification, which is
  unification modulo the equational theory of \isa{{\isasymalpha}{\isasymbeta}{\isasymeta}}-conversion
  on \isa{{\isasymlambda}}-terms.  Moreover, propositions are understood modulo
  the (derived) equivalence \isa{{\isacharparenleft}A\ {\isasymLongrightarrow}\ {\isacharparenleft}{\isasymAnd}x{\isachardot}\ B\ x{\isacharparenright}{\isacharparenright}\ {\isasymequiv}\ {\isacharparenleft}{\isasymAnd}x{\isachardot}\ A\ {\isasymLongrightarrow}\ B\ x{\isacharparenright}}.

  This means that any operations within the rule calculus may be
  subject to spontaneous \isa{{\isasymalpha}{\isasymbeta}{\isasymeta}}-HHF conversions.  It is common
  practice not to contract or expand unnecessarily.  Some mechanisms
  prefer an one form, others the opposite, so there is a potential
  danger to produce some oscillation!

  Only few operations really work \emph{modulo} HHF conversion, but
  expect a normal form: quantifiers \isa{{\isasymAnd}} before implications
  \isa{{\isasymLongrightarrow}} at each level of nesting.

\glossary{Hereditary Harrop Formula}{The set of propositions in HHF
format is defined inductively as \isa{H\ {\isacharequal}\ {\isacharparenleft}{\isasymAnd}x\isactrlsup {\isacharasterisk}{\isachardot}\ H\isactrlsup {\isacharasterisk}\ {\isasymLongrightarrow}\ A{\isacharparenright}}, for variables \isa{x} and atomic propositions \isa{A}.
Any proposition may be put into HHF form by normalizing with the rule
\isa{{\isacharparenleft}A\ {\isasymLongrightarrow}\ {\isacharparenleft}{\isasymAnd}x{\isachardot}\ B\ x{\isacharparenright}{\isacharparenright}\ {\isasymequiv}\ {\isacharparenleft}{\isasymAnd}x{\isachardot}\ A\ {\isasymLongrightarrow}\ B\ x{\isacharparenright}}.  In Isabelle, the outermost
quantifier prefix is represented via \seeglossary{schematic
variables}, such that the top-level structure is merely that of a
\seeglossary{Horn Clause}}.

\glossary{HHF}{See \seeglossary{Hereditary Harrop Formula}.}


  \[
  \infer[\isa{{\isacharparenleft}assumption{\isacharparenright}}]{\isa{C{\isasymvartheta}}}
  {\isa{{\isacharparenleft}{\isasymAnd}\isactrlvec x{\isachardot}\ \isactrlvec H\ \isactrlvec x\ {\isasymLongrightarrow}\ A\ \isactrlvec x{\isacharparenright}\ {\isasymLongrightarrow}\ C} & \isa{A{\isasymvartheta}\ {\isacharequal}\ H\isactrlsub i{\isasymvartheta}}~~\text{(for some~\isa{i})}}
  \]


  \[
  \infer[\isa{{\isacharparenleft}compose{\isacharparenright}}]{\isa{\isactrlvec A{\isasymvartheta}\ {\isasymLongrightarrow}\ C{\isasymvartheta}}}
  {\isa{\isactrlvec A\ {\isasymLongrightarrow}\ B} & \isa{B{\isacharprime}\ {\isasymLongrightarrow}\ C} & \isa{B{\isasymvartheta}\ {\isacharequal}\ B{\isacharprime}{\isasymvartheta}}}
  \]


  \[
  \infer[\isa{{\isacharparenleft}{\isasymAnd}{\isacharunderscore}lift{\isacharparenright}}]{\isa{{\isacharparenleft}{\isasymAnd}\isactrlvec x{\isachardot}\ \isactrlvec A\ {\isacharparenleft}{\isacharquery}\isactrlvec a\ \isactrlvec x{\isacharparenright}{\isacharparenright}\ {\isasymLongrightarrow}\ {\isacharparenleft}{\isasymAnd}\isactrlvec x{\isachardot}\ B\ {\isacharparenleft}{\isacharquery}\isactrlvec a\ \isactrlvec x{\isacharparenright}{\isacharparenright}}}{\isa{\isactrlvec A\ {\isacharquery}\isactrlvec a\ {\isasymLongrightarrow}\ B\ {\isacharquery}\isactrlvec a}}
  \]
  \[
  \infer[\isa{{\isacharparenleft}{\isasymLongrightarrow}{\isacharunderscore}lift{\isacharparenright}}]{\isa{{\isacharparenleft}\isactrlvec H\ {\isasymLongrightarrow}\ \isactrlvec A{\isacharparenright}\ {\isasymLongrightarrow}\ {\isacharparenleft}\isactrlvec H\ {\isasymLongrightarrow}\ B{\isacharparenright}}}{\isa{\isactrlvec A\ {\isasymLongrightarrow}\ B}}
  \]

  The \isa{resolve} scheme is now acquired from \isa{{\isasymAnd}{\isacharunderscore}lift},
  \isa{{\isasymLongrightarrow}{\isacharunderscore}lift}, and \isa{compose}.

  \[
  \infer[\isa{{\isacharparenleft}resolution{\isacharparenright}}]
  {\isa{{\isacharparenleft}{\isasymAnd}\isactrlvec x{\isachardot}\ \isactrlvec H\ \isactrlvec x\ {\isasymLongrightarrow}\ \isactrlvec A\ {\isacharparenleft}{\isacharquery}\isactrlvec a\ \isactrlvec x{\isacharparenright}{\isacharparenright}{\isasymvartheta}\ {\isasymLongrightarrow}\ C{\isasymvartheta}}}
  {\begin{tabular}{l}
    \isa{\isactrlvec A\ {\isacharquery}\isactrlvec a\ {\isasymLongrightarrow}\ B\ {\isacharquery}\isactrlvec a} \\
    \isa{{\isacharparenleft}{\isasymAnd}\isactrlvec x{\isachardot}\ \isactrlvec H\ \isactrlvec x\ {\isasymLongrightarrow}\ B{\isacharprime}\ \isactrlvec x{\isacharparenright}\ {\isasymLongrightarrow}\ C} \\
    \isa{{\isacharparenleft}{\isasymlambda}\isactrlvec x{\isachardot}\ B\ {\isacharparenleft}{\isacharquery}\isactrlvec a\ \isactrlvec x{\isacharparenright}{\isacharparenright}{\isasymvartheta}\ {\isacharequal}\ B{\isacharprime}{\isasymvartheta}} \\
   \end{tabular}}
  \]


  FIXME \isa{elim{\isacharunderscore}resolution}, \isa{dest{\isacharunderscore}resolution}%
\end{isamarkuptext}%
\isamarkuptrue%
%
\endisatagFIXME
{\isafoldFIXME}%
%
\isadelimFIXME
%
\endisadelimFIXME
%
\isadelimtheory
%
\endisadelimtheory
%
\isatagtheory
\isacommand{end}\isamarkupfalse%
%
\endisatagtheory
{\isafoldtheory}%
%
\isadelimtheory
%
\endisadelimtheory
\isanewline
\end{isabellebody}%
%%% Local Variables:
%%% mode: latex
%%% TeX-master: "root"
%%% End:

%% $Id$
\chapter{Tactics} \label{tactics}
\index{tactics|(} Tactics have type \mltydx{tactic}.  This is just an
abbreviation for functions from theorems to theorem sequences, where
the theorems represent states of a backward proof.  Tactics seldom
need to be coded from scratch, as functions; instead they are
expressed using basic tactics and tacticals.

This chapter only presents the primitive tactics.  Substantial proofs
require the power of automatic tools like simplification
(Chapter~\ref{chap:simplification}) and classical tableau reasoning
(Chapter~\ref{chap:classical}).

\section{Resolution and assumption tactics}
{\bf Resolution} is Isabelle's basic mechanism for refining a subgoal using
a rule.  {\bf Elim-resolution} is particularly suited for elimination
rules, while {\bf destruct-resolution} is particularly suited for
destruction rules.  The {\tt r}, {\tt e}, {\tt d} naming convention is
maintained for several different kinds of resolution tactics, as well as
the shortcuts in the subgoal module.

All the tactics in this section act on a subgoal designated by a positive
integer~$i$.  They fail (by returning the empty sequence) if~$i$ is out of
range.

\subsection{Resolution tactics}
\index{resolution!tactics}
\index{tactics!resolution|bold}
\begin{ttbox} 
resolve_tac  : thm list -> int -> tactic
eresolve_tac : thm list -> int -> tactic
dresolve_tac : thm list -> int -> tactic
forward_tac  : thm list -> int -> tactic 
\end{ttbox}
These perform resolution on a list of theorems, $thms$, representing a list
of object-rules.  When generating next states, they take each of the rules
in the order given.  Each rule may yield several next states, or none:
higher-order resolution may yield multiple resolvents.
\begin{ttdescription}
\item[\ttindexbold{resolve_tac} {\it thms} {\it i}] 
  refines the proof state using the rules, which should normally be
  introduction rules.  It resolves a rule's conclusion with
  subgoal~$i$ of the proof state.

\item[\ttindexbold{eresolve_tac} {\it thms} {\it i}] 
  \index{elim-resolution}
  performs elim-resolution with the rules, which should normally be
  elimination rules.  It resolves with a rule, proves its first premise by
  assumption, and finally \emph{deletes} that assumption from any new
  subgoals.  (To rotate a rule's premises,
  see \texttt{rotate_prems} in~\S\ref{MiscellaneousForwardRules}.)

\item[\ttindexbold{dresolve_tac} {\it thms} {\it i}] 
  \index{forward proof}\index{destruct-resolution}
  performs destruct-resolution with the rules, which normally should
  be destruction rules.  This replaces an assumption by the result of
  applying one of the rules.

\item[\ttindexbold{forward_tac}]\index{forward proof}
  is like {\tt dresolve_tac} except that the selected assumption is not
  deleted.  It applies a rule to an assumption, adding the result as a new
  assumption.
\end{ttdescription}

\subsection{Assumption tactics}
\index{tactics!assumption|bold}\index{assumptions!tactics for}
\begin{ttbox} 
assume_tac    : int -> tactic
eq_assume_tac : int -> tactic
\end{ttbox} 
\begin{ttdescription}
\item[\ttindexbold{assume_tac} {\it i}] 
attempts to solve subgoal~$i$ by assumption.

\item[\ttindexbold{eq_assume_tac}] 
is like {\tt assume_tac} but does not use unification.  It succeeds (with a
\emph{unique} next state) if one of the assumptions is identical to the
subgoal's conclusion.  Since it does not instantiate variables, it cannot
make other subgoals unprovable.  It is intended to be called from proof
strategies, not interactively.
\end{ttdescription}

\subsection{Matching tactics} \label{match_tac}
\index{tactics!matching}
\begin{ttbox} 
match_tac  : thm list -> int -> tactic
ematch_tac : thm list -> int -> tactic
dmatch_tac : thm list -> int -> tactic
\end{ttbox}
These are just like the resolution tactics except that they never
instantiate unknowns in the proof state.  Flexible subgoals are not updated
willy-nilly, but are left alone.  Matching --- strictly speaking --- means
treating the unknowns in the proof state as constants; these tactics merely
discard unifiers that would update the proof state.
\begin{ttdescription}
\item[\ttindexbold{match_tac} {\it thms} {\it i}] 
refines the proof state using the rules, matching a rule's
conclusion with subgoal~$i$ of the proof state.

\item[\ttindexbold{ematch_tac}] 
is like {\tt match_tac}, but performs elim-resolution.

\item[\ttindexbold{dmatch_tac}] 
is like {\tt match_tac}, but performs destruct-resolution.
\end{ttdescription}


\subsection{Resolution with instantiation} \label{res_inst_tac}
\index{tactics!instantiation}\index{instantiation}
\begin{ttbox} 
res_inst_tac  : (string*string)list -> thm -> int -> tactic
eres_inst_tac : (string*string)list -> thm -> int -> tactic
dres_inst_tac : (string*string)list -> thm -> int -> tactic
forw_inst_tac : (string*string)list -> thm -> int -> tactic
\end{ttbox}
These tactics are designed for applying rules such as substitution and
induction, which cause difficulties for higher-order unification.  The
tactics accept explicit instantiations for unknowns in the rule ---
typically, in the rule's conclusion.  Each instantiation is a pair
{\tt($v$,$e$)}, where $v$ is an unknown \emph{without} its leading
question mark!
\begin{itemize}
\item If $v$ is the type unknown {\tt'a}, then
the rule must contain a type unknown \verb$?'a$ of some
sort~$s$, and $e$ should be a type of sort $s$.

\item If $v$ is the unknown {\tt P}, then
the rule must contain an unknown \verb$?P$ of some type~$\tau$,
and $e$ should be a term of some type~$\sigma$ such that $\tau$ and
$\sigma$ are unifiable.  If the unification of $\tau$ and $\sigma$
instantiates any type unknowns in $\tau$, these instantiations
are recorded for application to the rule.
\end{itemize}
Types are instantiated before terms.  Because type instantiations are
inferred from term instantiations, explicit type instantiations are seldom
necessary --- if \verb$?t$ has type \verb$?'a$, then the instantiation list
\verb$[("'a","bool"),("t","True")]$ may be simplified to
\verb$[("t","True")]$.  Type unknowns in the proof state may cause
failure because the tactics cannot instantiate them.

The instantiation tactics act on a given subgoal.  Terms in the
instantiations are type-checked in the context of that subgoal --- in
particular, they may refer to that subgoal's parameters.  Any unknowns in
the terms receive subscripts and are lifted over the parameters; thus, you
may not refer to unknowns in the subgoal.

\begin{ttdescription}
\item[\ttindexbold{res_inst_tac} {\it insts} {\it thm} {\it i}]
instantiates the rule {\it thm} with the instantiations {\it insts}, as
described above, and then performs resolution on subgoal~$i$.  Resolution
typically causes further instantiations; you need not give explicit
instantiations for every unknown in the rule.

\item[\ttindexbold{eres_inst_tac}] 
is like {\tt res_inst_tac}, but performs elim-resolution.

\item[\ttindexbold{dres_inst_tac}] 
is like {\tt res_inst_tac}, but performs destruct-resolution.

\item[\ttindexbold{forw_inst_tac}] 
is like {\tt dres_inst_tac} except that the selected assumption is not
deleted.  It applies the instantiated rule to an assumption, adding the
result as a new assumption.
\end{ttdescription}


\section{Other basic tactics}
\subsection{Tactic shortcuts}
\index{shortcuts!for tactics}
\index{tactics!resolution}\index{tactics!assumption}
\index{tactics!meta-rewriting}
\begin{ttbox} 
rtac     :      thm -> int -> tactic
etac     :      thm -> int -> tactic
dtac     :      thm -> int -> tactic
atac     :             int -> tactic
ares_tac : thm list -> int -> tactic
rewtac   :      thm ->        tactic
\end{ttbox}
These abbreviate common uses of tactics.
\begin{ttdescription}
\item[\ttindexbold{rtac} {\it thm} {\it i}] 
abbreviates \hbox{\tt resolve_tac [{\it thm}] {\it i}}, doing resolution.

\item[\ttindexbold{etac} {\it thm} {\it i}] 
abbreviates \hbox{\tt eresolve_tac [{\it thm}] {\it i}}, doing elim-resolution.

\item[\ttindexbold{dtac} {\it thm} {\it i}] 
abbreviates \hbox{\tt dresolve_tac [{\it thm}] {\it i}}, doing
destruct-resolution.

\item[\ttindexbold{atac} {\it i}] 
abbreviates \hbox{\tt assume_tac {\it i}}, doing proof by assumption.

\item[\ttindexbold{ares_tac} {\it thms} {\it i}] 
tries proof by assumption and resolution; it abbreviates
\begin{ttbox}
assume_tac {\it i} ORELSE resolve_tac {\it thms} {\it i}
\end{ttbox}

\item[\ttindexbold{rewtac} {\it def}] 
abbreviates \hbox{\tt rewrite_goals_tac [{\it def}]}, unfolding a definition.
\end{ttdescription}


\subsection{Inserting premises and facts}\label{cut_facts_tac}
\index{tactics!for inserting facts}\index{assumptions!inserting}
\begin{ttbox} 
cut_facts_tac : thm list -> int -> tactic
cut_inst_tac  : (string*string)list -> thm -> int -> tactic
subgoal_tac   : string -> int -> tactic
subgoal_tacs  : string list -> int -> tactic
\end{ttbox}
These tactics add assumptions to a subgoal.
\begin{ttdescription}
\item[\ttindexbold{cut_facts_tac} {\it thms} {\it i}] 
  adds the {\it thms} as new assumptions to subgoal~$i$.  Once they have
  been inserted as assumptions, they become subject to tactics such as {\tt
    eresolve_tac} and {\tt rewrite_goals_tac}.  Only rules with no premises
  are inserted: Isabelle cannot use assumptions that contain $\Imp$
  or~$\Forall$.  Sometimes the theorems are premises of a rule being
  derived, returned by~{\tt goal}; instead of calling this tactic, you
  could state the goal with an outermost meta-quantifier.

\item[\ttindexbold{cut_inst_tac} {\it insts} {\it thm} {\it i}]
  instantiates the {\it thm} with the instantiations {\it insts}, as
  described in \S\ref{res_inst_tac}.  It adds the resulting theorem as a
  new assumption to subgoal~$i$. 

\item[\ttindexbold{subgoal_tac} {\it formula} {\it i}] 
adds the {\it formula} as a assumption to subgoal~$i$, and inserts the same
{\it formula} as a new subgoal, $i+1$.

\item[\ttindexbold{subgoals_tac} {\it formulae} {\it i}] 
  uses {\tt subgoal_tac} to add the members of the list of {\it
    formulae} as assumptions to subgoal~$i$. 
\end{ttdescription}


\subsection{``Putting off'' a subgoal}
\begin{ttbox} 
defer_tac : int -> tactic
\end{ttbox}
\begin{ttdescription}
\item[\ttindexbold{defer_tac} {\it i}] 
  moves subgoal~$i$ to the last position in the proof state.  It can be
  useful when correcting a proof script: if the tactic given for subgoal~$i$
  fails, calling {\tt defer_tac} instead will let you continue with the rest
  of the script.

  The tactic fails if subgoal~$i$ does not exist or if the proof state
  contains type unknowns. 
\end{ttdescription}


\subsection{Definitions and meta-level rewriting} \label{sec:rewrite_goals}
\index{tactics!meta-rewriting|bold}\index{meta-rewriting|bold}
\index{definitions}

Definitions in Isabelle have the form $t\equiv u$, where $t$ is typically a
constant or a constant applied to a list of variables, for example $\it
sqr(n)\equiv n\times n$.  Conditional definitions, $\phi\Imp t\equiv u$,
are also supported.  {\bf Unfolding} the definition ${t\equiv u}$ means using
it as a rewrite rule, replacing~$t$ by~$u$ throughout a theorem.  {\bf
Folding} $t\equiv u$ means replacing~$u$ by~$t$.  Rewriting continues until
no rewrites are applicable to any subterm.

There are rules for unfolding and folding definitions; Isabelle does not do
this automatically.  The corresponding tactics rewrite the proof state,
yielding a single next state.  See also the {\tt goalw} command, which is the
easiest way of handling definitions.
\begin{ttbox} 
rewrite_goals_tac : thm list -> tactic
rewrite_tac       : thm list -> tactic
fold_goals_tac    : thm list -> tactic
fold_tac          : thm list -> tactic
\end{ttbox}
\begin{ttdescription}
\item[\ttindexbold{rewrite_goals_tac} {\it defs}]  
unfolds the {\it defs} throughout the subgoals of the proof state, while
leaving the main goal unchanged.  Use \ttindex{SELECT_GOAL} to restrict it to a
particular subgoal.

\item[\ttindexbold{rewrite_tac} {\it defs}]  
unfolds the {\it defs} throughout the proof state, including the main goal
--- not normally desirable!

\item[\ttindexbold{fold_goals_tac} {\it defs}]  
folds the {\it defs} throughout the subgoals of the proof state, while
leaving the main goal unchanged.

\item[\ttindexbold{fold_tac} {\it defs}]  
folds the {\it defs} throughout the proof state.
\end{ttdescription}

\begin{warn}
  These tactics only cope with definitions expressed as meta-level
  equalities ($\equiv$).  More general equivalences are handled by the
  simplifier, provided that it is set up appropriately for your logic
  (see Chapter~\ref{chap:simplification}).
\end{warn}

\subsection{Theorems useful with tactics}
\index{theorems!of pure theory}
\begin{ttbox} 
asm_rl: thm 
cut_rl: thm 
\end{ttbox}
\begin{ttdescription}
\item[\tdx{asm_rl}] 
is $\psi\Imp\psi$.  Under elim-resolution it does proof by assumption, and
\hbox{\tt eresolve_tac (asm_rl::{\it thms}) {\it i}} is equivalent to
\begin{ttbox} 
assume_tac {\it i}  ORELSE  eresolve_tac {\it thms} {\it i}
\end{ttbox}

\item[\tdx{cut_rl}] 
is $\List{\psi\Imp\theta,\psi}\Imp\theta$.  It is useful for inserting
assumptions; it underlies {\tt forward_tac}, {\tt cut_facts_tac}
and {\tt subgoal_tac}.
\end{ttdescription}


\section{Obscure tactics}

\subsection{Renaming parameters in a goal} \index{parameters!renaming}
\begin{ttbox} 
rename_tac        : string -> int -> tactic
rename_last_tac   : string -> string list -> int -> tactic
Logic.set_rename_prefix : string -> unit
Logic.auto_rename       : bool ref      \hfill{\bf initially false}
\end{ttbox}
When creating a parameter, Isabelle chooses its name by matching variable
names via the object-rule.  Given the rule $(\forall I)$ formalized as
$\left(\Forall x. P(x)\right) \Imp \forall x.P(x)$, Isabelle will note that
the $\Forall$-bound variable in the premise has the same name as the
$\forall$-bound variable in the conclusion.  

Sometimes there is insufficient information and Isabelle chooses an
arbitrary name.  The renaming tactics let you override Isabelle's choice.
Because renaming parameters has no logical effect on the proof state, the
{\tt by} command prints the message {\tt Warning:\ same as previous
level}.

Alternatively, you can suppress the naming mechanism described above and
have Isabelle generate uniform names for parameters.  These names have the
form $p${\tt a}, $p${\tt b}, $p${\tt c},~\ldots, where $p$ is any desired
prefix.  They are ugly but predictable.

\begin{ttdescription}
\item[\ttindexbold{rename_tac} {\it str} {\it i}] 
interprets the string {\it str} as a series of blank-separated variable
names, and uses them to rename the parameters of subgoal~$i$.  The names
must be distinct.  If there are fewer names than parameters, then the
tactic renames the innermost parameters and may modify the remaining ones
to ensure that all the parameters are distinct.

\item[\ttindexbold{rename_last_tac} {\it prefix} {\it suffixes} {\it i}] 
generates a list of names by attaching each of the {\it suffixes\/} to the 
{\it prefix}.  It is intended for coding structural induction tactics,
where several of the new parameters should have related names.

\item[\ttindexbold{Logic.set_rename_prefix} {\it prefix};] 
sets the prefix for uniform renaming to~{\it prefix}.  The default prefix
is {\tt"k"}.

\item[set \ttindexbold{Logic.auto_rename};] 
makes Isabelle generate uniform names for parameters. 
\end{ttdescription}


\subsection{Manipulating assumptions}
\index{assumptions!rotating}
\begin{ttbox} 
thin_tac   : string -> int -> tactic
rotate_tac : int -> int -> tactic
\end{ttbox}
\begin{ttdescription}
\item[\ttindexbold{thin_tac} {\it formula} $i$]  
\index{assumptions!deleting}
deletes the specified assumption from subgoal $i$.  Often the assumption
can be abbreviated, replacing subformul{\ae} by unknowns; the first matching
assumption will be deleted.  Removing useless assumptions from a subgoal
increases its readability and can make search tactics run faster.

\item[\ttindexbold{rotate_tac} $n$ $i$]  
\index{assumptions!rotating}
rotates the assumptions of subgoal $i$ by $n$ positions: from right to left
if $n$ is positive, and from left to right if $n$ is negative.  This is 
sometimes necessary in connection with \ttindex{asm_full_simp_tac}, which 
processes assumptions from left to right.
\end{ttdescription}


\subsection{Tidying the proof state}
\index{duplicate subgoals!removing}
\index{parameters!removing unused}
\index{flex-flex constraints}
\begin{ttbox} 
distinct_subgoals_tac : tactic
prune_params_tac      : tactic
flexflex_tac          : tactic
\end{ttbox}
\begin{ttdescription}
\item[\ttindexbold{distinct_subgoals_tac}]  
  removes duplicate subgoals from a proof state.  (These arise especially
  in \ZF{}, where the subgoals are essentially type constraints.)

\item[\ttindexbold{prune_params_tac}]  
  removes unused parameters from all subgoals of the proof state.  It works
  by rewriting with the theorem $(\Forall x. V)\equiv V$.  This tactic can
  make the proof state more readable.  It is used with
  \ttindex{rule_by_tactic} to simplify the resulting theorem.

\item[\ttindexbold{flexflex_tac}]  
  removes all flex-flex pairs from the proof state by applying the trivial
  unifier.  This drastic step loses information, and should only be done as
  the last step of a proof.

  Flex-flex constraints arise from difficult cases of higher-order
  unification.  To prevent this, use \ttindex{res_inst_tac} to instantiate
  some variables in a rule~(\S\ref{res_inst_tac}).  Normally flex-flex
  constraints can be ignored; they often disappear as unknowns get
  instantiated.
\end{ttdescription}


\subsection{Composition: resolution without lifting}
\index{tactics!for composition}
\begin{ttbox}
compose_tac: (bool * thm * int) -> int -> tactic
\end{ttbox}
{\bf Composing} two rules means resolving them without prior lifting or
renaming of unknowns.  This low-level operation, which underlies the
resolution tactics, may occasionally be useful for special effects.
A typical application is \ttindex{res_inst_tac}, which lifts and instantiates a
rule, then passes the result to {\tt compose_tac}.
\begin{ttdescription}
\item[\ttindexbold{compose_tac} ($flag$, $rule$, $m$) $i$] 
refines subgoal~$i$ using $rule$, without lifting.  The $rule$ is taken to
have the form $\List{\psi@1; \ldots; \psi@m} \Imp \psi$, where $\psi$ need
not be atomic; thus $m$ determines the number of new subgoals.  If
$flag$ is {\tt true} then it performs elim-resolution --- it solves the
first premise of~$rule$ by assumption and deletes that assumption.
\end{ttdescription}


\section{*Managing lots of rules}
These operations are not intended for interactive use.  They are concerned
with the processing of large numbers of rules in automatic proof
strategies.  Higher-order resolution involving a long list of rules is
slow.  Filtering techniques can shorten the list of rules given to
resolution, and can also detect whether a subgoal is too flexible,
with too many rules applicable.

\subsection{Combined resolution and elim-resolution} \label{biresolve_tac}
\index{tactics!resolution}
\begin{ttbox} 
biresolve_tac   : (bool*thm)list -> int -> tactic
bimatch_tac     : (bool*thm)list -> int -> tactic
subgoals_of_brl : bool*thm -> int
lessb           : (bool*thm) * (bool*thm) -> bool
\end{ttbox}
{\bf Bi-resolution} takes a list of $\it (flag,rule)$ pairs.  For each
pair, it applies resolution if the flag is~{\tt false} and
elim-resolution if the flag is~{\tt true}.  A single tactic call handles a
mixture of introduction and elimination rules.

\begin{ttdescription}
\item[\ttindexbold{biresolve_tac} {\it brls} {\it i}] 
refines the proof state by resolution or elim-resolution on each rule, as
indicated by its flag.  It affects subgoal~$i$ of the proof state.

\item[\ttindexbold{bimatch_tac}] 
is like {\tt biresolve_tac}, but performs matching: unknowns in the
proof state are never updated (see~\S\ref{match_tac}).

\item[\ttindexbold{subgoals_of_brl}({\it flag},{\it rule})] 
returns the number of new subgoals that bi-res\-o\-lu\-tion would yield for the
pair (if applied to a suitable subgoal).  This is $n$ if the flag is
{\tt false} and $n-1$ if the flag is {\tt true}, where $n$ is the number
of premises of the rule.  Elim-resolution yields one fewer subgoal than
ordinary resolution because it solves the major premise by assumption.

\item[\ttindexbold{lessb} ({\it brl1},{\it brl2})] 
returns the result of 
\begin{ttbox}
subgoals_of_brl{\it brl1} < subgoals_of_brl{\it brl2}
\end{ttbox}
\end{ttdescription}
Note that \hbox{\tt sort lessb {\it brls}} sorts a list of $\it
(flag,rule)$ pairs by the number of new subgoals they will yield.  Thus,
those that yield the fewest subgoals should be tried first.


\subsection{Discrimination nets for fast resolution}\label{filt_resolve_tac}
\index{discrimination nets|bold}
\index{tactics!resolution}
\begin{ttbox} 
net_resolve_tac  : thm list -> int -> tactic
net_match_tac    : thm list -> int -> tactic
net_biresolve_tac: (bool*thm) list -> int -> tactic
net_bimatch_tac  : (bool*thm) list -> int -> tactic
filt_resolve_tac : thm list -> int -> int -> tactic
could_unify      : term*term->bool
filter_thms      : (term*term->bool) -> int*term*thm list -> thm list
\end{ttbox}
The module {\tt Net} implements a discrimination net data structure for
fast selection of rules \cite[Chapter 14]{charniak80}.  A term is
classified by the symbol list obtained by flattening it in preorder.
The flattening takes account of function applications, constants, and free
and bound variables; it identifies all unknowns and also regards
\index{lambda abs@$\lambda$-abstractions}
$\lambda$-abstractions as unknowns, since they could $\eta$-contract to
anything.  

A discrimination net serves as a polymorphic dictionary indexed by terms.
The module provides various functions for inserting and removing items from
nets.  It provides functions for returning all items whose term could match
or unify with a target term.  The matching and unification tests are
overly lax (due to the identifications mentioned above) but they serve as
useful filters.

A net can store introduction rules indexed by their conclusion, and
elimination rules indexed by their major premise.  Isabelle provides
several functions for `compiling' long lists of rules into fast
resolution tactics.  When supplied with a list of theorems, these functions
build a discrimination net; the net is used when the tactic is applied to a
goal.  To avoid repeatedly constructing the nets, use currying: bind the
resulting tactics to \ML{} identifiers.

\begin{ttdescription}
\item[\ttindexbold{net_resolve_tac} {\it thms}] 
builds a discrimination net to obtain the effect of a similar call to {\tt
resolve_tac}.

\item[\ttindexbold{net_match_tac} {\it thms}] 
builds a discrimination net to obtain the effect of a similar call to {\tt
match_tac}.

\item[\ttindexbold{net_biresolve_tac} {\it brls}] 
builds a discrimination net to obtain the effect of a similar call to {\tt
biresolve_tac}.

\item[\ttindexbold{net_bimatch_tac} {\it brls}] 
builds a discrimination net to obtain the effect of a similar call to {\tt
bimatch_tac}.

\item[\ttindexbold{filt_resolve_tac} {\it thms} {\it maxr} {\it i}] 
uses discrimination nets to extract the {\it thms} that are applicable to
subgoal~$i$.  If more than {\it maxr\/} theorems are applicable then the
tactic fails.  Otherwise it calls {\tt resolve_tac}.  

This tactic helps avoid runaway instantiation of unknowns, for example in
type inference.

\item[\ttindexbold{could_unify} ({\it t},{\it u})] 
returns {\tt false} if~$t$ and~$u$ are `obviously' non-unifiable, and
otherwise returns~{\tt true}.  It assumes all variables are distinct,
reporting that {\tt ?a=?a} may unify with {\tt 0=1}.

\item[\ttindexbold{filter_thms} $could\; (limit,prem,thms)$] 
returns the list of potentially resolvable rules (in {\it thms\/}) for the
subgoal {\it prem}, using the predicate {\it could\/} to compare the
conclusion of the subgoal with the conclusion of each rule.  The resulting list
is no longer than {\it limit}.
\end{ttdescription}


\section{Programming tools for proof strategies}
Do not consider using the primitives discussed in this section unless you
really need to code tactics from scratch.

\subsection{Operations on type {\tt tactic}}
\index{tactics!primitives for coding} A tactic maps theorems to sequences of
theorems.  The type constructor for sequences (lazy lists) is called
\mltydx{Seq.seq}.  To simplify the types of tactics and tacticals,
Isabelle defines a type abbreviation:
\begin{ttbox} 
type tactic = thm -> thm Seq.seq
\end{ttbox} 
The following operations provide means for coding tactics in a clean style.
\begin{ttbox} 
PRIMITIVE :                  (thm -> thm) -> tactic  
SUBGOAL   : ((term*int) -> tactic) -> int -> tactic
\end{ttbox} 
\begin{ttdescription}
\item[\ttindexbold{PRIMITIVE} $f$] packages the meta-rule~$f$ as a tactic that
  applies $f$ to the proof state and returns the result as a one-element
  sequence.  If $f$ raises an exception, then the tactic's result is the empty
  sequence.

\item[\ttindexbold{SUBGOAL} $f$ $i$] 
extracts subgoal~$i$ from the proof state as a term~$t$, and computes a
tactic by calling~$f(t,i)$.  It applies the resulting tactic to the same
state.  The tactic body is expressed using tactics and tacticals, but may
peek at a particular subgoal:
\begin{ttbox} 
SUBGOAL (fn (t,i) => {\it tactic-valued expression})
\end{ttbox} 
\end{ttdescription}


\subsection{Tracing}
\index{tactics!tracing}
\index{tracing!of tactics}
\begin{ttbox} 
pause_tac: tactic
print_tac: tactic
\end{ttbox}
These tactics print tracing information when they are applied to a proof
state.  Their output may be difficult to interpret.  Note that certain of
the searching tacticals, such as {\tt REPEAT}, have built-in tracing
options.
\begin{ttdescription}
\item[\ttindexbold{pause_tac}] 
prints {\footnotesize\tt** Press RETURN to continue:} and then reads a line
from the terminal.  If this line is blank then it returns the proof state
unchanged; otherwise it fails (which may terminate a repetition).

\item[\ttindexbold{print_tac}] 
returns the proof state unchanged, with the side effect of printing it at
the terminal.
\end{ttdescription}


\section{*Sequences}
\index{sequences (lazy lists)|bold}
The module {\tt Seq} declares a type of lazy lists.  It uses
Isabelle's type \mltydx{option} to represent the possible presence
(\ttindexbold{Some}) or absence (\ttindexbold{None}) of
a value:
\begin{ttbox}
datatype 'a option = None  |  Some of 'a;
\end{ttbox}
The {\tt Seq} structure is supposed to be accessed via fully qualified
names and should not be \texttt{open}.

\subsection{Basic operations on sequences}
\begin{ttbox} 
Seq.empty   : 'a seq
Seq.make    : (unit -> ('a * 'a seq) option) -> 'a seq
Seq.single  : 'a -> 'a seq
Seq.pull    : 'a seq -> ('a * 'a seq) option
\end{ttbox}
\begin{ttdescription}
\item[Seq.empty] is the empty sequence.

\item[\tt Seq.make (fn () => Some ($x$, $xq$))] constructs the
  sequence with head~$x$ and tail~$xq$, neither of which is evaluated.

\item[Seq.single $x$] 
constructs the sequence containing the single element~$x$.

\item[Seq.pull $xq$] returns {\tt None} if the sequence is empty and
  {\tt Some ($x$, $xq'$)} if the sequence has head~$x$ and tail~$xq'$.
  Warning: calling \hbox{Seq.pull $xq$} again will {\it recompute\/}
  the value of~$x$; it is not stored!
\end{ttdescription}


\subsection{Converting between sequences and lists}
\begin{ttbox} 
Seq.chop    : int * 'a seq -> 'a list * 'a seq
Seq.list_of : 'a seq -> 'a list
Seq.of_list : 'a list -> 'a seq
\end{ttbox}
\begin{ttdescription}
\item[Seq.chop ($n$, $xq$)] returns the first~$n$ elements of~$xq$ as a
  list, paired with the remaining elements of~$xq$.  If $xq$ has fewer
  than~$n$ elements, then so will the list.
  
\item[Seq.list_of $xq$] returns the elements of~$xq$, which must be
  finite, as a list.
  
\item[Seq.of_list $xs$] creates a sequence containing the elements
  of~$xs$.
\end{ttdescription}


\subsection{Combining sequences}
\begin{ttbox} 
Seq.append      : 'a seq * 'a seq -> 'a seq
Seq.interleave  : 'a seq * 'a seq -> 'a seq
Seq.flat        : 'a seq seq -> 'a seq
Seq.map         : ('a -> 'b) -> 'a seq -> 'b seq
Seq.filter      : ('a -> bool) -> 'a seq -> 'a seq
\end{ttbox} 
\begin{ttdescription}
\item[Seq.append ($xq$, $yq$)] concatenates $xq$ to $yq$.
  
\item[Seq.interleave ($xq$, $yq$)] joins $xq$ with $yq$ by
  interleaving their elements.  The result contains all the elements
  of the sequences, even if both are infinite.
  
\item[Seq.flat $xqq$] concatenates a sequence of sequences.
  
\item[Seq.map $f$ $xq$] applies $f$ to every element
  of~$xq=x@1,x@2,\ldots$, yielding the sequence $f(x@1),f(x@2),\ldots$.
  
\item[Seq.filter $p$ $xq$] returns the sequence consisting of all
  elements~$x$ of~$xq$ such that $p(x)$ is {\tt true}.
\end{ttdescription}

\index{tactics|)}

%
\begin{isabellebody}%
\def\isabellecontext{Proof}%
%
\isadelimtheory
%
\endisadelimtheory
%
\isatagtheory
\isacommand{theory}\isamarkupfalse%
\ Proof\isanewline
\isakeyword{imports}\ Base\isanewline
\isakeyword{begin}%
\endisatagtheory
{\isafoldtheory}%
%
\isadelimtheory
%
\endisadelimtheory
%
\isamarkupchapter{Structured proofs%
}
\isamarkuptrue%
%
\isamarkupsection{Variables \label{sec:variables}%
}
\isamarkuptrue%
%
\begin{isamarkuptext}%
Any variable that is not explicitly bound by \isa{{\isasymlambda}}-abstraction
  is considered as ``free''.  Logically, free variables act like
  outermost universal quantification at the sequent level: \isa{A\isactrlisub {\isadigit{1}}{\isacharparenleft}x{\isacharparenright}{\isacharcomma}\ {\isasymdots}{\isacharcomma}\ A\isactrlisub n{\isacharparenleft}x{\isacharparenright}\ {\isasymturnstile}\ B{\isacharparenleft}x{\isacharparenright}} means that the result
  holds \emph{for all} values of \isa{x}.  Free variables for
  terms (not types) can be fully internalized into the logic: \isa{{\isasymturnstile}\ B{\isacharparenleft}x{\isacharparenright}} and \isa{{\isasymturnstile}\ {\isasymAnd}x{\isachardot}\ B{\isacharparenleft}x{\isacharparenright}} are interchangeable, provided
  that \isa{x} does not occur elsewhere in the context.
  Inspecting \isa{{\isasymturnstile}\ {\isasymAnd}x{\isachardot}\ B{\isacharparenleft}x{\isacharparenright}} more closely, we see that inside the
  quantifier, \isa{x} is essentially ``arbitrary, but fixed'',
  while from outside it appears as a place-holder for instantiation
  (thanks to \isa{{\isasymAnd}} elimination).

  The Pure logic represents the idea of variables being either inside
  or outside the current scope by providing separate syntactic
  categories for \emph{fixed variables} (e.g.\ \isa{x}) vs.\
  \emph{schematic variables} (e.g.\ \isa{{\isacharquery}x}).  Incidently, a
  universal result \isa{{\isasymturnstile}\ {\isasymAnd}x{\isachardot}\ B{\isacharparenleft}x{\isacharparenright}} has the HHF normal form \isa{{\isasymturnstile}\ B{\isacharparenleft}{\isacharquery}x{\isacharparenright}}, which represents its generality nicely without requiring
  an explicit quantifier.  The same principle works for type
  variables: \isa{{\isasymturnstile}\ B{\isacharparenleft}{\isacharquery}{\isasymalpha}{\isacharparenright}} represents the idea of ``\isa{{\isasymturnstile}\ {\isasymforall}{\isasymalpha}{\isachardot}\ B{\isacharparenleft}{\isasymalpha}{\isacharparenright}}'' without demanding a truly polymorphic framework.

  \medskip Additional care is required to treat type variables in a
  way that facilitates type-inference.  In principle, term variables
  depend on type variables, which means that type variables would have
  to be declared first.  For example, a raw type-theoretic framework
  would demand the context to be constructed in stages as follows:
  \isa{{\isasymGamma}\ {\isacharequal}\ {\isasymalpha}{\isacharcolon}\ type{\isacharcomma}\ x{\isacharcolon}\ {\isasymalpha}{\isacharcomma}\ a{\isacharcolon}\ A{\isacharparenleft}x\isactrlisub {\isasymalpha}{\isacharparenright}}.

  We allow a slightly less formalistic mode of operation: term
  variables \isa{x} are fixed without specifying a type yet
  (essentially \emph{all} potential occurrences of some instance
  \isa{x\isactrlisub {\isasymtau}} are fixed); the first occurrence of \isa{x}
  within a specific term assigns its most general type, which is then
  maintained consistently in the context.  The above example becomes
  \isa{{\isasymGamma}\ {\isacharequal}\ x{\isacharcolon}\ term{\isacharcomma}\ {\isasymalpha}{\isacharcolon}\ type{\isacharcomma}\ A{\isacharparenleft}x\isactrlisub {\isasymalpha}{\isacharparenright}}, where type \isa{{\isasymalpha}} is fixed \emph{after} term \isa{x}, and the constraint
  \isa{x\ {\isacharcolon}{\isacharcolon}\ {\isasymalpha}} is an implicit consequence of the occurrence of
  \isa{x\isactrlisub {\isasymalpha}} in the subsequent proposition.

  This twist of dependencies is also accommodated by the reverse
  operation of exporting results from a context: a type variable
  \isa{{\isasymalpha}} is considered fixed as long as it occurs in some fixed
  term variable of the context.  For example, exporting \isa{x{\isacharcolon}\ term{\isacharcomma}\ {\isasymalpha}{\isacharcolon}\ type\ {\isasymturnstile}\ x\isactrlisub {\isasymalpha}\ {\isacharequal}\ x\isactrlisub {\isasymalpha}} produces in the first step
  \isa{x{\isacharcolon}\ term\ {\isasymturnstile}\ x\isactrlisub {\isasymalpha}\ {\isacharequal}\ x\isactrlisub {\isasymalpha}} for fixed \isa{{\isasymalpha}},
  and only in the second step \isa{{\isasymturnstile}\ {\isacharquery}x\isactrlisub {\isacharquery}\isactrlisub {\isasymalpha}\ {\isacharequal}\ {\isacharquery}x\isactrlisub {\isacharquery}\isactrlisub {\isasymalpha}} for schematic \isa{{\isacharquery}x} and \isa{{\isacharquery}{\isasymalpha}}.

  \medskip The Isabelle/Isar proof context manages the gory details of
  term vs.\ type variables, with high-level principles for moving the
  frontier between fixed and schematic variables.

  The \isa{add{\isacharunderscore}fixes} operation explictly declares fixed
  variables; the \isa{declare{\isacharunderscore}term} operation absorbs a term into
  a context by fixing new type variables and adding syntactic
  constraints.

  The \isa{export} operation is able to perform the main work of
  generalizing term and type variables as sketched above, assuming
  that fixing variables and terms have been declared properly.

  There \isa{import} operation makes a generalized fact a genuine
  part of the context, by inventing fixed variables for the schematic
  ones.  The effect can be reversed by using \isa{export} later,
  potentially with an extended context; the result is equivalent to
  the original modulo renaming of schematic variables.

  The \isa{focus} operation provides a variant of \isa{import}
  for nested propositions (with explicit quantification): \isa{{\isasymAnd}x\isactrlisub {\isadigit{1}}\ {\isasymdots}\ x\isactrlisub n{\isachardot}\ B{\isacharparenleft}x\isactrlisub {\isadigit{1}}{\isacharcomma}\ {\isasymdots}{\isacharcomma}\ x\isactrlisub n{\isacharparenright}} is
  decomposed by inventing fixed variables \isa{x\isactrlisub {\isadigit{1}}{\isacharcomma}\ {\isasymdots}{\isacharcomma}\ x\isactrlisub n} for the body.%
\end{isamarkuptext}%
\isamarkuptrue%
%
\isadelimmlref
%
\endisadelimmlref
%
\isatagmlref
%
\begin{isamarkuptext}%
\begin{mldecls}
  \indexml{Variable.add\_fixes}\verb|Variable.add_fixes: |\isasep\isanewline%
\verb|  string list -> Proof.context -> string list * Proof.context| \\
  \indexml{Variable.variant\_fixes}\verb|Variable.variant_fixes: |\isasep\isanewline%
\verb|  string list -> Proof.context -> string list * Proof.context| \\
  \indexml{Variable.declare\_term}\verb|Variable.declare_term: term -> Proof.context -> Proof.context| \\
  \indexml{Variable.declare\_constraints}\verb|Variable.declare_constraints: term -> Proof.context -> Proof.context| \\
  \indexml{Variable.export}\verb|Variable.export: Proof.context -> Proof.context -> thm list -> thm list| \\
  \indexml{Variable.polymorphic}\verb|Variable.polymorphic: Proof.context -> term list -> term list| \\
  \indexml{Variable.import\_thms}\verb|Variable.import_thms: bool -> thm list -> Proof.context ->|\isasep\isanewline%
\verb|  ((ctyp list * cterm list) * thm list) * Proof.context| \\
  \indexml{Variable.focus}\verb|Variable.focus: cterm -> Proof.context -> (cterm list * cterm) * Proof.context| \\
  \end{mldecls}

  \begin{description}

  \item \verb|Variable.add_fixes|~\isa{xs\ ctxt} fixes term
  variables \isa{xs}, returning the resulting internal names.  By
  default, the internal representation coincides with the external
  one, which also means that the given variables must not be fixed
  already.  There is a different policy within a local proof body: the
  given names are just hints for newly invented Skolem variables.

  \item \verb|Variable.variant_fixes| is similar to \verb|Variable.add_fixes|, but always produces fresh variants of the given
  names.

  \item \verb|Variable.declare_term|~\isa{t\ ctxt} declares term
  \isa{t} to belong to the context.  This automatically fixes new
  type variables, but not term variables.  Syntactic constraints for
  type and term variables are declared uniformly, though.

  \item \verb|Variable.declare_constraints|~\isa{t\ ctxt} declares
  syntactic constraints from term \isa{t}, without making it part
  of the context yet.

  \item \verb|Variable.export|~\isa{inner\ outer\ thms} generalizes
  fixed type and term variables in \isa{thms} according to the
  difference of the \isa{inner} and \isa{outer} context,
  following the principles sketched above.

  \item \verb|Variable.polymorphic|~\isa{ctxt\ ts} generalizes type
  variables in \isa{ts} as far as possible, even those occurring
  in fixed term variables.  The default policy of type-inference is to
  fix newly introduced type variables, which is essentially reversed
  with \verb|Variable.polymorphic|: here the given terms are detached
  from the context as far as possible.

  \item \verb|Variable.import_thms|~\isa{open\ thms\ ctxt} invents fixed
  type and term variables for the schematic ones occurring in \isa{thms}.  The \isa{open} flag indicates whether the fixed names
  should be accessible to the user, otherwise newly introduced names
  are marked as ``internal'' (\secref{sec:names}).

  \item \verb|Variable.focus|~\isa{B} decomposes the outermost \isa{{\isasymAnd}} prefix of proposition \isa{B}.

  \end{description}%
\end{isamarkuptext}%
\isamarkuptrue%
%
\endisatagmlref
{\isafoldmlref}%
%
\isadelimmlref
%
\endisadelimmlref
%
\isamarkupsection{Assumptions \label{sec:assumptions}%
}
\isamarkuptrue%
%
\begin{isamarkuptext}%
An \emph{assumption} is a proposition that it is postulated in the
  current context.  Local conclusions may use assumptions as
  additional facts, but this imposes implicit hypotheses that weaken
  the overall statement.

  Assumptions are restricted to fixed non-schematic statements, i.e.\
  all generality needs to be expressed by explicit quantifiers.
  Nevertheless, the result will be in HHF normal form with outermost
  quantifiers stripped.  For example, by assuming \isa{{\isasymAnd}x\ {\isacharcolon}{\isacharcolon}\ {\isasymalpha}{\isachardot}\ P\ x} we get \isa{{\isasymAnd}x\ {\isacharcolon}{\isacharcolon}\ {\isasymalpha}{\isachardot}\ P\ x\ {\isasymturnstile}\ P\ {\isacharquery}x} for schematic \isa{{\isacharquery}x}
  of fixed type \isa{{\isasymalpha}}.  Local derivations accumulate more and
  more explicit references to hypotheses: \isa{A\isactrlisub {\isadigit{1}}{\isacharcomma}\ {\isasymdots}{\isacharcomma}\ A\isactrlisub n\ {\isasymturnstile}\ B} where \isa{A\isactrlisub {\isadigit{1}}{\isacharcomma}\ {\isasymdots}{\isacharcomma}\ A\isactrlisub n} needs to
  be covered by the assumptions of the current context.

  \medskip The \isa{add{\isacharunderscore}assms} operation augments the context by
  local assumptions, which are parameterized by an arbitrary \isa{export} rule (see below).

  The \isa{export} operation moves facts from a (larger) inner
  context into a (smaller) outer context, by discharging the
  difference of the assumptions as specified by the associated export
  rules.  Note that the discharged portion is determined by the
  difference contexts, not the facts being exported!  There is a
  separate flag to indicate a goal context, where the result is meant
  to refine an enclosing sub-goal of a structured proof state.

  \medskip The most basic export rule discharges assumptions directly
  by means of the \isa{{\isasymLongrightarrow}} introduction rule:
  \[
  \infer[(\isa{{\isasymLongrightarrow}{\isacharunderscore}intro})]{\isa{{\isasymGamma}\ {\isacharbackslash}\ A\ {\isasymturnstile}\ A\ {\isasymLongrightarrow}\ B}}{\isa{{\isasymGamma}\ {\isasymturnstile}\ B}}
  \]

  The variant for goal refinements marks the newly introduced
  premises, which causes the canonical Isar goal refinement scheme to
  enforce unification with local premises within the goal:
  \[
  \infer[(\isa{{\isacharhash}{\isasymLongrightarrow}{\isacharunderscore}intro})]{\isa{{\isasymGamma}\ {\isacharbackslash}\ A\ {\isasymturnstile}\ {\isacharhash}A\ {\isasymLongrightarrow}\ B}}{\isa{{\isasymGamma}\ {\isasymturnstile}\ B}}
  \]

  \medskip Alternative versions of assumptions may perform arbitrary
  transformations on export, as long as the corresponding portion of
  hypotheses is removed from the given facts.  For example, a local
  definition works by fixing \isa{x} and assuming \isa{x\ {\isasymequiv}\ t},
  with the following export rule to reverse the effect:
  \[
  \infer[(\isa{{\isasymequiv}{\isacharminus}expand})]{\isa{{\isasymGamma}\ {\isacharbackslash}\ x\ {\isasymequiv}\ t\ {\isasymturnstile}\ B\ t}}{\isa{{\isasymGamma}\ {\isasymturnstile}\ B\ x}}
  \]
  This works, because the assumption \isa{x\ {\isasymequiv}\ t} was introduced in
  a context with \isa{x} being fresh, so \isa{x} does not
  occur in \isa{{\isasymGamma}} here.%
\end{isamarkuptext}%
\isamarkuptrue%
%
\isadelimmlref
%
\endisadelimmlref
%
\isatagmlref
%
\begin{isamarkuptext}%
\begin{mldecls}
  \indexmltype{Assumption.export}\verb|type Assumption.export| \\
  \indexml{Assumption.assume}\verb|Assumption.assume: cterm -> thm| \\
  \indexml{Assumption.add\_assms}\verb|Assumption.add_assms: Assumption.export ->|\isasep\isanewline%
\verb|  cterm list -> Proof.context -> thm list * Proof.context| \\
  \indexml{Assumption.add\_assumes}\verb|Assumption.add_assumes: |\isasep\isanewline%
\verb|  cterm list -> Proof.context -> thm list * Proof.context| \\
  \indexml{Assumption.export}\verb|Assumption.export: bool -> Proof.context -> Proof.context -> thm -> thm| \\
  \end{mldecls}

  \begin{description}

  \item \verb|Assumption.export| represents arbitrary export
  rules, which is any function of type \verb|bool -> cterm list -> thm -> thm|,
  where the \verb|bool| indicates goal mode, and the \verb|cterm list| the collection of assumptions to be discharged
  simultaneously.

  \item \verb|Assumption.assume|~\isa{A} turns proposition \isa{A} into a raw assumption \isa{A\ {\isasymturnstile}\ A{\isacharprime}}, where the conclusion
  \isa{A{\isacharprime}} is in HHF normal form.

  \item \verb|Assumption.add_assms|~\isa{r\ As} augments the context
  by assumptions \isa{As} with export rule \isa{r}.  The
  resulting facts are hypothetical theorems as produced by the raw
  \verb|Assumption.assume|.

  \item \verb|Assumption.add_assumes|~\isa{As} is a special case of
  \verb|Assumption.add_assms| where the export rule performs \isa{{\isasymLongrightarrow}{\isacharunderscore}intro} or \isa{{\isacharhash}{\isasymLongrightarrow}{\isacharunderscore}intro}, depending on goal mode.

  \item \verb|Assumption.export|~\isa{is{\isacharunderscore}goal\ inner\ outer\ thm}
  exports result \isa{thm} from the the \isa{inner} context
  back into the \isa{outer} one; \isa{is{\isacharunderscore}goal\ {\isacharequal}\ true} means
  this is a goal context.  The result is in HHF normal form.  Note
  that \verb|ProofContext.export| combines \verb|Variable.export|
  and \verb|Assumption.export| in the canonical way.

  \end{description}%
\end{isamarkuptext}%
\isamarkuptrue%
%
\endisatagmlref
{\isafoldmlref}%
%
\isadelimmlref
%
\endisadelimmlref
%
\isamarkupsection{Results \label{sec:results}%
}
\isamarkuptrue%
%
\begin{isamarkuptext}%
Local results are established by monotonic reasoning from facts
  within a context.  This allows common combinations of theorems,
  e.g.\ via \isa{{\isasymAnd}{\isacharslash}{\isasymLongrightarrow}} elimination, resolution rules, or equational
  reasoning, see \secref{sec:thms}.  Unaccounted context manipulations
  should be avoided, notably raw \isa{{\isasymAnd}{\isacharslash}{\isasymLongrightarrow}} introduction or ad-hoc
  references to free variables or assumptions not present in the proof
  context.

  \medskip The \isa{SUBPROOF} combinator allows to structure a
  tactical proof recursively by decomposing a selected sub-goal:
  \isa{{\isacharparenleft}{\isasymAnd}x{\isachardot}\ A{\isacharparenleft}x{\isacharparenright}\ {\isasymLongrightarrow}\ B{\isacharparenleft}x{\isacharparenright}{\isacharparenright}\ {\isasymLongrightarrow}\ {\isasymdots}} is turned into \isa{B{\isacharparenleft}x{\isacharparenright}\ {\isasymLongrightarrow}\ {\isasymdots}}
  after fixing \isa{x} and assuming \isa{A{\isacharparenleft}x{\isacharparenright}}.  This means
  the tactic needs to solve the conclusion, but may use the premise as
  a local fact, for locally fixed variables.

  The \isa{prove} operation provides an interface for structured
  backwards reasoning under program control, with some explicit sanity
  checks of the result.  The goal context can be augmented by
  additional fixed variables (cf.\ \secref{sec:variables}) and
  assumptions (cf.\ \secref{sec:assumptions}), which will be available
  as local facts during the proof and discharged into implications in
  the result.  Type and term variables are generalized as usual,
  according to the context.

  The \isa{obtain} operation produces results by eliminating
  existing facts by means of a given tactic.  This acts like a dual
  conclusion: the proof demonstrates that the context may be augmented
  by certain fixed variables and assumptions.  See also
  \cite{isabelle-isar-ref} for the user-level \isa{{\isasymOBTAIN}} and
  \isa{{\isasymGUESS}} elements.  Final results, which may not refer to
  the parameters in the conclusion, need to exported explicitly into
  the original context.%
\end{isamarkuptext}%
\isamarkuptrue%
%
\isadelimmlref
%
\endisadelimmlref
%
\isatagmlref
%
\begin{isamarkuptext}%
\begin{mldecls}
  \indexml{SUBPROOF}\verb|SUBPROOF: ({context: Proof.context, schematics: ctyp list * cterm list,|\isasep\isanewline%
\verb|    params: cterm list, asms: cterm list, concl: cterm,|\isasep\isanewline%
\verb|    prems: thm list} -> tactic) -> Proof.context -> int -> tactic| \\
  \end{mldecls}
  \begin{mldecls}
  \indexml{Goal.prove}\verb|Goal.prove: Proof.context -> string list -> term list -> term ->|\isasep\isanewline%
\verb|  ({prems: thm list, context: Proof.context} -> tactic) -> thm| \\
  \indexml{Goal.prove\_multi}\verb|Goal.prove_multi: Proof.context -> string list -> term list -> term list ->|\isasep\isanewline%
\verb|  ({prems: thm list, context: Proof.context} -> tactic) -> thm list| \\
  \end{mldecls}
  \begin{mldecls}
  \indexml{Obtain.result}\verb|Obtain.result: (Proof.context -> tactic) ->|\isasep\isanewline%
\verb|  thm list -> Proof.context -> (cterm list * thm list) * Proof.context| \\
  \end{mldecls}

  \begin{description}

  \item \verb|SUBPROOF|~\isa{tac\ ctxt\ i} decomposes the structure
  of the specified sub-goal, producing an extended context and a
  reduced goal, which needs to be solved by the given tactic.  All
  schematic parameters of the goal are imported into the context as
  fixed ones, which may not be instantiated in the sub-proof.

  \item \verb|Goal.prove|~\isa{ctxt\ xs\ As\ C\ tac} states goal \isa{C} in the context augmented by fixed variables \isa{xs} and
  assumptions \isa{As}, and applies tactic \isa{tac} to solve
  it.  The latter may depend on the local assumptions being presented
  as facts.  The result is in HHF normal form.

  \item \verb|Goal.prove_multi| is simular to \verb|Goal.prove|, but
  states several conclusions simultaneously.  The goal is encoded by
  means of Pure conjunction; \verb|Goal.conjunction_tac| will turn this
  into a collection of individual subgoals.

  \item \verb|Obtain.result|~\isa{tac\ thms\ ctxt} eliminates the
  given facts using a tactic, which results in additional fixed
  variables and assumptions in the context.  Final results need to be
  exported explicitly.

  \end{description}%
\end{isamarkuptext}%
\isamarkuptrue%
%
\endisatagmlref
{\isafoldmlref}%
%
\isadelimmlref
%
\endisadelimmlref
%
\isadelimtheory
%
\endisadelimtheory
%
\isatagtheory
\isacommand{end}\isamarkupfalse%
%
\endisatagtheory
{\isafoldtheory}%
%
\isadelimtheory
%
\endisadelimtheory
\isanewline
\end{isabellebody}%
%%% Local Variables:
%%% mode: latex
%%% TeX-master: "root"
%%% End:

%
\begin{isabellebody}%
\def\isabellecontext{Isar}%
%
\isadelimtheory
%
\endisadelimtheory
%
\isatagtheory
\isacommand{theory}\isamarkupfalse%
\ Isar\isanewline
\isakeyword{imports}\ Base\isanewline
\isakeyword{begin}%
\endisatagtheory
{\isafoldtheory}%
%
\isadelimtheory
%
\endisadelimtheory
%
\isamarkupchapter{Isar proof texts%
}
\isamarkuptrue%
%
\isamarkupsection{Proof context%
}
\isamarkuptrue%
%
\begin{isamarkuptext}%
FIXME%
\end{isamarkuptext}%
\isamarkuptrue%
%
\isamarkupsection{Proof state \label{sec:isar-proof-state}%
}
\isamarkuptrue%
%
\begin{isamarkuptext}%
FIXME

\glossary{Proof state}{The whole configuration of a structured proof,
consisting of a \seeglossary{proof context} and an optional
\seeglossary{structured goal}.  Internally, an Isar proof state is
organized as a stack to accomodate block structure of proof texts.
For historical reasons, a low-level \seeglossary{tactical goal} is
occasionally called ``proof state'' as well.}

\glossary{Structured goal}{FIXME}

\glossary{Goal}{See \seeglossary{tactical goal} or \seeglossary{structured goal}. \norefpage}%
\end{isamarkuptext}%
\isamarkuptrue%
%
\isamarkupsection{Proof methods%
}
\isamarkuptrue%
%
\begin{isamarkuptext}%
FIXME%
\end{isamarkuptext}%
\isamarkuptrue%
%
\isamarkupsection{Attributes%
}
\isamarkuptrue%
%
\begin{isamarkuptext}%
FIXME ?!%
\end{isamarkuptext}%
\isamarkuptrue%
%
\isadelimtheory
%
\endisadelimtheory
%
\isatagtheory
\isacommand{end}\isamarkupfalse%
%
\endisatagtheory
{\isafoldtheory}%
%
\isadelimtheory
%
\endisadelimtheory
\isanewline
\end{isabellebody}%
%%% Local Variables:
%%% mode: latex
%%% TeX-master: "root"
%%% End:

%
\begin{isabellebody}%
\def\isabellecontext{locale}%
%
\isadelimtheory
\isanewline
\isanewline
\isanewline
%
\endisadelimtheory
%
\isatagtheory
\isacommand{theory}\isamarkupfalse%
\ {\isachardoublequoteopen}locale{\isachardoublequoteclose}\ \isakeyword{imports}\ base\ \isakeyword{begin}%
\endisatagtheory
{\isafoldtheory}%
%
\isadelimtheory
%
\endisadelimtheory
%
\isamarkupchapter{Structured specifications%
}
\isamarkuptrue%
%
\isamarkupsection{Specification elements%
}
\isamarkuptrue%
%
\begin{isamarkuptext}%
FIXME%
\end{isamarkuptext}%
\isamarkuptrue%
%
\isamarkupsection{Locales%
}
\isamarkuptrue%
%
\begin{isamarkuptext}%
FIXME%
\end{isamarkuptext}%
\isamarkuptrue%
%
\isadelimtheory
%
\endisadelimtheory
%
\isatagtheory
\isacommand{end}\isamarkupfalse%
%
\endisatagtheory
{\isafoldtheory}%
%
\isadelimtheory
%
\endisadelimtheory
\isanewline
\end{isabellebody}%
%%% Local Variables:
%%% mode: latex
%%% TeX-master: "root"
%%% End:

%
\begin{isabellebody}%
\def\isabellecontext{integration}%
%
\isadelimtheory
\isanewline
\isanewline
\isanewline
%
\endisadelimtheory
%
\isatagtheory
\isacommand{theory}\isamarkupfalse%
\ integration\ \isakeyword{imports}\ base\ \isakeyword{begin}%
\endisatagtheory
{\isafoldtheory}%
%
\isadelimtheory
%
\endisadelimtheory
%
\isamarkupchapter{System integration%
}
\isamarkuptrue%
%
\isamarkupsection{Isar toplevel \label{sec:isar-toplevel}%
}
\isamarkuptrue%
%
\begin{isamarkuptext}%
The Isar toplevel may be considered the centeral hub of the
  Isabelle/Isar system, where all key components and sub-systems are
  integrated into a single read-eval-print loop of Isar commands.  We
  shall even incorporate the existing {\ML} toplevel of the compiler
  and run-time system (cf.\ \secref{sec:ML-toplevel}).

  Isabelle/Isar departs from the original ``LCF system architecture''
  where {\ML} was really The Meta Language for defining theories and
  conducting proofs.  Instead, {\ML} now only serves as the
  implementation language for the system (and user extensions), while
  the specific Isar toplevel supports the concepts of theory and proof
  development natively.  This includes the graph structure of theories
  and the block structure of proofs, support for unlimited undo,
  facilities for tracing, debugging, timing, profiling etc.

  \medskip The toplevel maintains an implicit state, which is
  transformed by a sequence of transitions -- either interactively or
  in batch-mode.  In interactive mode, Isar state transitions are
  encapsulated as safe transactions, such that both failure and undo
  are handled conveniently without destroying the underlying draft
  theory (cf.~\secref{sec:context-theory}).  In batch mode,
  transitions operate in a linear (destructive) fashion, such that
  error conditions abort the present attempt to construct a theory or
  proof altogether.

  The toplevel state is a disjoint sum of empty \isa{toplevel}, or
  \isa{theory}, or \isa{proof}.  On entering the main Isar loop we
  start with an empty toplevel.  A theory is commenced by giving a
  \isa{{\isasymTHEORY}} header; within a theory we may issue theory
  commands such as \isa{{\isasymDEFINITION}}, or state a \isa{{\isasymTHEOREM}} to be proven.  Now we are within a proof state, with a
  rich collection of Isar proof commands for structured proof
  composition, or unstructured proof scripts.  When the proof is
  concluded we get back to the theory, which is then updated by
  storing the resulting fact.  Further theory declarations or theorem
  statements with proofs may follow, until we eventually conclude the
  theory development by issuing \isa{{\isasymEND}}.  The resulting theory
  is then stored within the theory database and we are back to the
  empty toplevel.

  In addition to these proper state transformations, there are also
  some diagnostic commands for peeking at the toplevel state without
  modifying it (e.g.\ \isakeyword{thm}, \isakeyword{term},
  \isakeyword{print-cases}).%
\end{isamarkuptext}%
\isamarkuptrue%
%
\isadelimmlref
%
\endisadelimmlref
%
\isatagmlref
%
\begin{isamarkuptext}%
\begin{mldecls}
  \indexmltype{Toplevel.state}\verb|type Toplevel.state| \\
  \indexml{Toplevel.UNDEF}\verb|Toplevel.UNDEF: exn| \\
  \indexml{Toplevel.is-toplevel}\verb|Toplevel.is_toplevel: Toplevel.state -> bool| \\
  \indexml{Toplevel.theory-of}\verb|Toplevel.theory_of: Toplevel.state -> theory| \\
  \indexml{Toplevel.proof-of}\verb|Toplevel.proof_of: Toplevel.state -> Proof.state| \\
  \indexml{Toplevel.debug}\verb|Toplevel.debug: bool ref| \\
  \indexml{Toplevel.timing}\verb|Toplevel.timing: bool ref| \\
  \indexml{Toplevel.profiling}\verb|Toplevel.profiling: int ref| \\
  \end{mldecls}

  \begin{description}

  \item \verb|Toplevel.state| represents Isar toplevel states,
  which are normally manipulated through the concept of toplevel
  transitions only (\secref{sec:toplevel-transition}).  Also note that
  a raw toplevel state is subject to the same linearity restrictions
  as a theory context (cf.~\secref{sec:context-theory}).

  \item \verb|Toplevel.UNDEF| is raised for undefined toplevel
  operations.  Many operations work only partially for certain cases,
  since \verb|Toplevel.state| is a sum type.

  \item \verb|Toplevel.is_toplevel|~\isa{state} checks for an empty
  toplevel state.

  \item \verb|Toplevel.theory_of|~\isa{state} selects the theory of
  a theory or proof (!), otherwise raises \verb|Toplevel.UNDEF|.

  \item \verb|Toplevel.proof_of|~\isa{state} selects the Isar proof
  state if available, otherwise raises \verb|Toplevel.UNDEF|.

  \item \verb|set Toplevel.debug| makes the toplevel print further
  details about internal error conditions, exceptions being raised
  etc.

  \item \verb|set Toplevel.timing| makes the toplevel print timing
  information for each Isar command being executed.

  \item \verb|Toplevel.profiling|~\verb|:=|~\isa{n} controls
  low-level profiling of the underlying {\ML} runtime system.  For
  Poly/ML, \isa{n\ {\isacharequal}\ {\isadigit{1}}} means time and \isa{n\ {\isacharequal}\ {\isadigit{2}}} space
  profiling.

  \end{description}%
\end{isamarkuptext}%
\isamarkuptrue%
%
\endisatagmlref
{\isafoldmlref}%
%
\isadelimmlref
%
\endisadelimmlref
%
\isamarkupsubsection{Toplevel transitions \label{sec:toplevel-transition}%
}
\isamarkuptrue%
%
\begin{isamarkuptext}%
An Isar toplevel transition consists of a partial function on the
  toplevel state, with additional information for diagnostics and
  error reporting: there are fields for command name, source position,
  optional source text, as well as flags for interactive-only commands
  (which issue a warning in batch-mode), printing of result state,
  etc.

  The operational part is represented as the sequential union of a
  list of partial functions, which are tried in turn until the first
  one succeeds.  This acts like an outer case-expression for various
  alternative state transitions.  For example, \isakeyword{qed} acts
  differently for a local proofs vs.\ the global ending of the main
  proof.

  Toplevel transitions are composed via transition transformers.
  Internally, Isar commands are put together from an empty transition
  extended by name and source position (and optional source text).  It
  is then left to the individual command parser to turn the given
  concrete syntax into a suitable transition transformer that adjoin
  actual operations on a theory or proof state etc.%
\end{isamarkuptext}%
\isamarkuptrue%
%
\isadelimmlref
%
\endisadelimmlref
%
\isatagmlref
%
\begin{isamarkuptext}%
\begin{mldecls}
  \indexml{Toplevel.print}\verb|Toplevel.print: Toplevel.transition -> Toplevel.transition| \\
  \indexml{Toplevel.no-timing}\verb|Toplevel.no_timing: Toplevel.transition -> Toplevel.transition| \\
  \indexml{Toplevel.keep}\verb|Toplevel.keep: (Toplevel.state -> unit) ->|\isasep\isanewline%
\verb|  Toplevel.transition -> Toplevel.transition| \\
  \indexml{Toplevel.theory}\verb|Toplevel.theory: (theory -> theory) ->|\isasep\isanewline%
\verb|  Toplevel.transition -> Toplevel.transition| \\
  \indexml{Toplevel.theory-to-proof}\verb|Toplevel.theory_to_proof: (theory -> Proof.state) ->|\isasep\isanewline%
\verb|  Toplevel.transition -> Toplevel.transition| \\
  \indexml{Toplevel.proof}\verb|Toplevel.proof: (Proof.state -> Proof.state) ->|\isasep\isanewline%
\verb|  Toplevel.transition -> Toplevel.transition| \\
  \indexml{Toplevel.proofs}\verb|Toplevel.proofs: (Proof.state -> Proof.state Seq.seq) ->|\isasep\isanewline%
\verb|  Toplevel.transition -> Toplevel.transition| \\
  \indexml{Toplevel.end-proof}\verb|Toplevel.end_proof: (bool -> Proof.state -> Proof.context) ->|\isasep\isanewline%
\verb|  Toplevel.transition -> Toplevel.transition| \\
  \end{mldecls}

  \begin{description}

  \item \verb|Toplevel.print|~\isa{tr} sets the print flag, which
  causes the toplevel loop to echo the result state (in interactive
  mode).

  \item \verb|Toplevel.no_timing|~\isa{tr} indicates that the
  transition should never show timing information, e.g.\ because it is
  a diagnostic command.

  \item \verb|Toplevel.keep|~\isa{tr} adjoins a diagnostic
  function.

  \item \verb|Toplevel.theory|~\isa{tr} adjoins a theory
  transformer.

  \item \verb|Toplevel.theory_to_proof|~\isa{tr} adjoins a global
  goal function, which turns a theory into a proof state.  The theory
  may be changed before entering the proof; the generic Isar goal
  setup includes an argument that specifies how to apply the proven
  result to the theory, when the proof is finished.

  \item \verb|Toplevel.proof|~\isa{tr} adjoins a deterministic
  proof command, with a singleton result.

  \item \verb|Toplevel.proofs|~\isa{tr} adjoins a general proof
  command, with zero or more result states (represented as a lazy
  list).

  \item \verb|Toplevel.end_proof|~\isa{tr} adjoins a concluding
  proof command, that returns the resulting theory, after storing the
  resulting facts in the context etc.

  \end{description}%
\end{isamarkuptext}%
\isamarkuptrue%
%
\endisatagmlref
{\isafoldmlref}%
%
\isadelimmlref
%
\endisadelimmlref
%
\isamarkupsubsection{Toplevel control%
}
\isamarkuptrue%
%
\begin{isamarkuptext}%
There are a few special control commands that modify the behavior
  the toplevel itself, and only make sense in interactive mode.  Under
  normal circumstances, the user encounters these only implicitly as
  part of the protocol between the Isabelle/Isar system and a
  user-interface such as ProofGeneral.

  \begin{description}

  \item \isacommand{undo} follows the three-level hierarchy of empty
  toplevel vs.\ theory vs.\ proof: undo within a proof reverts to the
  previous proof context, undo after a proof reverts to the theory
  before the initial goal statement, undo of a theory command reverts
  to the previous theory value, undo of a theory header discontinues
  the current theory development and removes it from the theory
  database (\secref{sec:theory-database}).

  \item \isacommand{kill} aborts the current level of development:
  kill in a proof context reverts to the theory before the initial
  goal statement, kill in a theory context aborts the current theory
  development, removing it from the database.

  \item \isacommand{exit} drops out of the Isar toplevel into the
  underlying {\ML} toplevel (\secref{sec:ML-toplevel}).  The Isar
  toplevel state is preserved and may be continued later.

  \item \isacommand{quit} terminates the Isabelle/Isar process without
  saving.

  \end{description}%
\end{isamarkuptext}%
\isamarkuptrue%
%
\isamarkupsection{ML toplevel \label{sec:ML-toplevel}%
}
\isamarkuptrue%
%
\begin{isamarkuptext}%
The {\ML} toplevel provides a read-compile-eval-print loop for {\ML}
  values, types, structures, and functors.  {\ML} declarations operate
  on the global system state, which consists of the compiler
  environment plus the values of {\ML} reference variables.  There is
  no clean way to undo {\ML} declarations, except for reverting to a
  previously saved state of the whole Isabelle process.  {\ML} input
  is either read interactively from a TTY, or from a string (usually
  within a theory text), or from a source file (usually loaded from a
  theory).

  Whenever the {\ML} toplevel is active, the current Isabelle theory
  context is passed as an internal reference variable.  Thus {\ML}
  code may access the theory context during compilation, it may even
  change the value of a theory being under construction --- while
  observing the usual linearity restrictions
  (cf.~\secref{sec:context-theory}).%
\end{isamarkuptext}%
\isamarkuptrue%
%
\isadelimmlref
%
\endisadelimmlref
%
\isatagmlref
%
\begin{isamarkuptext}%
\begin{mldecls}
  \indexml{the-context}\verb|the_context: unit -> theory| \\
  \indexml{Context.$>$$>$ }\verb|Context.>> : (Context.generic -> Context.generic) -> unit| \\
  \end{mldecls}

  \begin{description}

  \item \verb|the_context ()| refers to the theory context of the
  {\ML} toplevel --- at compile time!  {\ML} code needs to take care
  to refer to \verb|the_context ()| correctly.  Recall that
  evaluation of a function body is delayed until actual runtime.
  Moreover, persistent {\ML} toplevel bindings to an unfinished theory
  should be avoided: code should either project out the desired
  information immediately, or produce an explicit \verb|theory_ref| (cf.\ \secref{sec:context-theory}).

  \item \verb|Context.>>|~\isa{f} applies context transformation
  \isa{f} to the implicit context of the {\ML} toplevel.

  \end{description}

  It is very important to note that the above functions are really
  restricted to the compile time, even though the {\ML} compiler is
  invoked at runtime!  The majority of {\ML} code uses explicit
  functional arguments of a theory or proof context instead.  Thus it
  may be invoked for an arbitrary context later on, without having to
  worry about any operational details.

  \bigskip

  \begin{mldecls}
  \indexml{Isar.main}\verb|Isar.main: unit -> unit| \\
  \indexml{Isar.loop}\verb|Isar.loop: unit -> unit| \\
  \indexml{Isar.state}\verb|Isar.state: unit -> Toplevel.state| \\
  \indexml{Isar.exn}\verb|Isar.exn: unit -> (exn * string) option| \\
  \indexml{Isar.context}\verb|Isar.context: unit -> Proof.context| \\
  \indexml{Isar.goal}\verb|Isar.goal: unit -> thm| \\
  \end{mldecls}

  \begin{description}

  \item \verb|Isar.main ()| invokes the Isar toplevel from {\ML},
  initializing an empty toplevel state.

  \item \verb|Isar.loop ()| continues the Isar toplevel with the
  current state, after having dropped out of the Isar toplevel loop.

  \item \verb|Isar.state ()| and \verb|Isar.exn ()| get current
  toplevel state and error condition, respectively.  This only works
  after having dropped out of the Isar toplevel loop.

  \item \verb|Isar.context ()| produces the proof context from \verb|Isar.state ()|, analogous to \verb|Context.proof_of|
  (\secref{sec:generic-context}).

  \item \verb|Isar.goal ()| picks the tactical goal from \verb|Isar.state ()|, represented as a theorem according to
  \secref{sec:tactical-goals}.

  \end{description}%
\end{isamarkuptext}%
\isamarkuptrue%
%
\endisatagmlref
{\isafoldmlref}%
%
\isadelimmlref
%
\endisadelimmlref
%
\isamarkupsection{Theory database \label{sec:theory-database}%
}
\isamarkuptrue%
%
\begin{isamarkuptext}%
The theory database maintains a collection of theories, together
  with some administrative information about their original sources,
  which are held in an external store (i.e.\ some directory within the
  regular file system).

  The theory database is organized as a directed acyclic graph;
  entries are referenced by theory name.  Although some additional
  interfaces allow to include a directory specification as well, this
  is only a hint to the underlying theory loader.  The internal theory
  name space is flat!

  Theory \isa{A} is associated with the main theory file \isa{A}\verb,.thy,, which needs to be accessible through the theory
  loader path.  Any number of additional {\ML} source files may be
  associated with each theory, by declaring these dependencies in the
  theory header as \isa{{\isasymUSES}}, and loading them consecutively
  within the theory context.  The system keeps track of incoming {\ML}
  sources and associates them with the current theory.  The file
  \isa{A}\verb,.ML, is loaded after a theory has been concluded, in
  order to support legacy proof {\ML} proof scripts.

  The basic internal actions of the theory database are \isa{update}, \isa{outdate}, and \isa{remove}:

  \begin{itemize}

  \item \isa{update\ A} introduces a link of \isa{A} with a
  \isa{theory} value of the same name; it asserts that the theory
  sources are now consistent with that value;

  \item \isa{outdate\ A} invalidates the link of a theory database
  entry to its sources, but retains the present theory value;

  \item \isa{remove\ A} deletes entry \isa{A} from the theory
  database.
  
  \end{itemize}

  These actions are propagated to sub- or super-graphs of a theory
  entry as expected, in order to preserve global consistency of the
  state of all loaded theories with the sources of the external store.
  This implies certain causalities between actions: \isa{update}
  or \isa{outdate} of an entry will \isa{outdate} all
  descendants; \isa{remove} will \isa{remove} all descendants.

  \medskip There are separate user-level interfaces to operate on the
  theory database directly or indirectly.  The primitive actions then
  just happen automatically while working with the system.  In
  particular, processing a theory header \isa{{\isasymTHEORY}\ A\ {\isasymIMPORTS}\ B\isactrlsub {\isadigit{1}}\ {\isasymdots}\ B\isactrlsub n\ {\isasymBEGIN}} ensures that the
  sub-graph of the collective imports \isa{B\isactrlsub {\isadigit{1}}\ {\isasymdots}\ B\isactrlsub n}
  is up-to-date, too.  Earlier theories are reloaded as required, with
  \isa{update} actions proceeding in topological order according to
  theory dependencies.  There may be also a wave of implied \isa{outdate} actions for derived theory nodes until a stable situation
  is achieved eventually.%
\end{isamarkuptext}%
\isamarkuptrue%
%
\isadelimmlref
%
\endisadelimmlref
%
\isatagmlref
%
\begin{isamarkuptext}%
\begin{mldecls}
  \indexml{theory}\verb|theory: string -> theory| \\
  \indexml{use-thy}\verb|use_thy: string -> unit| \\
  \indexml{use-thys}\verb|use_thys: string list -> unit| \\
  \indexml{touch-thy}\verb|touch_thy: string -> unit| \\
  \indexml{remove-thy}\verb|remove_thy: string -> unit| \\[1ex]
  \indexml{ThyInfo.begin-theory}\verb|ThyInfo.begin_theory|\verb|: ... -> bool -> theory| \\
  \indexml{ThyInfo.end-theory}\verb|ThyInfo.end_theory: theory -> theory| \\
  \indexml{ThyInfo.register-theory}\verb|ThyInfo.register_theory: theory -> unit| \\[1ex]
  \verb|datatype action = Update |\verb,|,\verb| Outdate |\verb,|,\verb| Remove| \\
  \indexml{ThyInfo.add-hook}\verb|ThyInfo.add_hook: (ThyInfo.action -> string -> unit) -> unit| \\
  \end{mldecls}

  \begin{description}

  \item \verb|theory|~\isa{A} retrieves the theory value presently
  associated with name \isa{A}.  Note that the result might be
  outdated.

  \item \verb|use_thy|~\isa{A} ensures that theory \isa{A} is fully
  up-to-date wrt.\ the external file store, reloading outdated
  ancestors as required.

  \item \verb|use_thys| is similar to \verb|use_thy|, but handles
  several theories simultaneously.  Thus it acts like processing the
  import header of a theory, without performing the merge of the
  result, though.

  \item \verb|touch_thy|~\isa{A} performs and \isa{outdate} action
  on theory \isa{A} and all descendants.

  \item \verb|remove_thy|~\isa{A} deletes theory \isa{A} and all
  descendants from the theory database.

  \item \verb|ThyInfo.begin_theory| is the basic operation behind a
  \isa{{\isasymTHEORY}} header declaration.  This is {\ML} functions is
  normally not invoked directly.

  \item \verb|ThyInfo.end_theory| concludes the loading of a theory
  proper.  An attached theory {\ML} file may be still loaded later on.
  This is function is normally not invoked directly.

  \item \verb|ThyInfo.register_theory|~\isa{text\ thy} registers an
  existing theory value with the theory loader database.  There is no
  management of associated sources.

  \item \verb|ThyInfo.add_hook|~\isa{f} registers function \isa{f} as a hook for theory database actions.  The function will be
  invoked with the action and theory name being involved; thus derived
  actions may be performed in associated system components, e.g.\
  maintaining the state of an editor for the theory sources.

  The kind and order of actions occurring in practice depends both on
  user interactions and the internal process of resolving theory
  imports.  Hooks should not rely on a particular policy here!  Any
  exceptions raised by the hook are ignored.

  \end{description}%
\end{isamarkuptext}%
\isamarkuptrue%
%
\endisatagmlref
{\isafoldmlref}%
%
\isadelimmlref
%
\endisadelimmlref
%
\isadelimtheory
%
\endisadelimtheory
%
\isatagtheory
\isacommand{end}\isamarkupfalse%
%
\endisatagtheory
{\isafoldtheory}%
%
\isadelimtheory
%
\endisadelimtheory
\isanewline
\end{isabellebody}%
%%% Local Variables:
%%% mode: latex
%%% TeX-master: "root"
%%% End:


\appendix
%
\begin{isabellebody}%
\def\isabellecontext{ML}%
%
\isadelimtheory
\isanewline
\isanewline
%
\endisadelimtheory
%
\isatagtheory
\isacommand{theory}\isamarkupfalse%
\ {\isachardoublequoteopen}ML{\isachardoublequoteclose}\ \isakeyword{imports}\ base\ \isakeyword{begin}%
\endisatagtheory
{\isafoldtheory}%
%
\isadelimtheory
%
\endisadelimtheory
%
\isamarkupchapter{Aesthetics of ML programming%
}
\isamarkuptrue%
%
\begin{isamarkuptext}%
FIXME%
\end{isamarkuptext}%
\isamarkuptrue%
%
\begin{isamarkuptext}%
This style guide is loosely based on
  \url{http://caml.inria.fr/resources/doc/guides/guidelines.en.html}.
%  FIMXE \url{http://www.cs.cornell.edu/Courses/cs312/2003sp/handouts/style.htm}

  Like any style guide, it should not be interpreted dogmatically.
  Instead, it forms a collection of recommendations which,
  if obeyed, result in code that is not considered to be
  obfuscated.  In certain cases, derivations are encouraged,
  as far as you know what you are doing.

  \begin{description}

    \item[fundamental law of programming]
      Whenever writing code, keep in mind: A program is
      written once, modified ten times, and read
      100 times.  So simplify its writing,
      always keep future modifications in mind,
      and never jeopardize readability.  Every second you hesitate
      to spend on making your code more clear you will
      have to spend ten times understanding what you have
      written later on.

    \item[white space matters]
      Treat white space in your code as if it determines
      the meaning of code.

      \begin{itemize}

        \item The space bar is the easiest key to find on the keyboard,
          press it as often as necessary. {\ttfamily 2 + 2} is better
          than {\ttfamily 2+2}, likewise {\ttfamily f (x, y)}
          better than {\ttfamily f(x,y)}.

        \item Restrict your lines to \emph{at most} 80 characters.
          This will allow you to keep the beginning of a line
          in view while watching its end.

        \item Ban tabs; they are a context-sensitive formatting
          feature and likely to confuse anyone not using your
          favourite editor.

        \item Get rid of trailing whitespace.  Instead, do not
          surpess a trailing newline at the end of your files.

        \item Choose a generally accepted style of indentation,
          then use it systematically throughout the whole
          application.  An indentation of two spaces is appropriate.
          Avoid dangling indentation.

      \end{itemize}

    \item[cut-and-paste succeeds over copy-and-paste]
      \emph{Never} copy-and-paste code when programming.  If you
        need the same piece of code twice, introduce a
        reasonable auxiliary function (if there is no
        such function, very likely you got something wrong).
        Any copy-and-paste will turn out to be painful 
        when something has to be changed or fixed later on.

    \item[comments]
      are a device which requires careful thinking before using
      it.  The best comment for your code should be the code itself.
      Prefer efforts to write clear, understandable code
      over efforts to explain nasty code.

    \item[functional programming is based on functions]
      Avoid ``constructivisms'', e.g. pass a table lookup function,
      rather than an actual table with lookup in body.  Accustom
      your way of codeing to the level of expressiveness
      a functional programming language is giving onto you.

    \item[tuples]
      are often in the way.  When there is no striking argument
      to tuple function arguments, just write your function curried.

    \item[telling names]
      Any name should tell its purpose as exactly as possible,
      while keeping its length to the absolutely neccessary minimum.
      Always give the same name to function arguments which
      have the same meaning. Separate words by underscores
      (``\verb|int_of_string|'', not ``\verb|intOfString|'')

  \end{description}%
\end{isamarkuptext}%
\isamarkuptrue%
%
\isamarkupchapter{Basic library functions%
}
\isamarkuptrue%
%
\begin{isamarkuptext}%
Beyond the proposal of the SML/NJ basis library, Isabelle comes
  with its own library, from which selected parts are given here.
  See further files \emph{Pure/library.ML} and \emph{Pure/General/*.ML}.%
\end{isamarkuptext}%
\isamarkuptrue%
%
\isamarkupsection{Linear transformations%
}
\isamarkuptrue%
%
\isadelimmlref
%
\endisadelimmlref
%
\isatagmlref
%
\begin{isamarkuptext}%
\begin{mldecls}
  \indexml{(op |$>$)}\verb|(op |\verb,|,\verb|>): 'a * ('a -> 'b) -> 'b| \\
  \indexml{fold}\verb|fold: ('a -> 'b -> 'b) -> 'a list -> 'b -> 'b| \\
  \indexml{fold-rev}\verb|fold_rev: ('a -> 'b -> 'b) -> 'a list -> 'b -> 'b| \\
  \end{mldecls}%
\end{isamarkuptext}%
\isamarkuptrue%
%
\endisatagmlref
{\isafoldmlref}%
%
\isadelimmlref
%
\endisadelimmlref
%
\isadelimML
%
\endisadelimML
%
\isatagML
%
\endisatagML
{\isafoldML}%
%
\isadelimML
%
\endisadelimML
%
\begin{isamarkuptext}%
Many problems in functional programming can be thought of
  as linear transformations, i.e.~a caluclation starts with a
  particular value \isa{x\ {\isasymColon}\ foo} which is then transformed
  by application of a function \isa{f\ {\isasymColon}\ foo\ {\isasymRightarrow}\ foo},
  continued by an application of a function \isa{g\ {\isasymColon}\ foo\ {\isasymRightarrow}\ bar},
  and so on.  As a canoncial example, take primitive functions enriching
  theories by constants and definitions:
  \verb|Sign.add_consts_i: (string * typ * mixfix) list -> theory|\isasep\isanewline%
\verb|-> theory|
  and \verb|Theory.add_defs_i: bool -> bool|\isasep\isanewline%
\verb|-> (bstring * term) list -> theory -> theory|.
  Written with naive application, an addition of a constant with
  a corresponding definition would look like:
  \verb|Theory.add_defs_i false false [dummy_def]|\isasep\isanewline%
\verb|  (Sign.add_consts_i [dummy_const] thy)|.
  With increasing numbers of applications, this code gets quite unreadable.
  Using composition, at least the nesting of brackets may be reduced:
  \verb|(Theory.add_defs_i false false [dummy_def] o Sign.add_consts_i|\isasep\isanewline%
\verb|  [dummy_const]) thy|.
  What remains unsatisfactory is that things are written down in the opposite order
  as they actually ``happen''.%
\end{isamarkuptext}%
\isamarkuptrue%
%
\isadelimML
%
\endisadelimML
%
\isatagML
%
\endisatagML
{\isafoldML}%
%
\isadelimML
%
\endisadelimML
%
\begin{isamarkuptext}%
At this stage, Isabelle offers some combinators which allow for more convenient
  notation, most notably reverse application:
  \isasep\isanewline%
\verb|thy|\isasep\isanewline%
\verb||\verb,|,\verb|> Sign.add_consts_i [dummy_const]|\isasep\isanewline%
\verb||\verb,|,\verb|> Theory.add_defs_i false false [dummy_def]|%
\end{isamarkuptext}%
\isamarkuptrue%
%
\begin{isamarkuptext}%
\noindent When iterating over a list of parameters \isa{{\isacharbrackleft}x\isactrlisub {\isadigit{1}}{\isacharcomma}\ x\isactrlisub {\isadigit{2}}{\isacharcomma}\ {\isasymdots}\ x\isactrlisub n{\isacharbrackright}\ {\isasymColon}\ {\isacharprime}a\ list},
  the \verb|fold| combinator lifts a single function \isa{f\ {\isasymColon}\ {\isacharprime}a\ {\isacharminus}{\isachargreater}\ {\isacharprime}b\ {\isacharminus}{\isachargreater}\ {\isacharprime}b}:
  \isa{y\ {\isacharbar}{\isachargreater}\ fold\ f\ {\isacharbrackleft}x\isactrlisub {\isadigit{1}}{\isacharcomma}\ x\isactrlisub {\isadigit{2}}{\isacharcomma}\ {\isasymdots}\ x\isactrlisub n{\isacharbrackright}\ {\isasymequiv}\ y\ {\isacharbar}{\isachargreater}\ f\ x\isactrlisub {\isadigit{1}}\ {\isacharbar}{\isachargreater}\ f\ x\isactrlisub {\isadigit{2}}\ {\isacharbar}{\isachargreater}\ {\isasymdots}\ {\isacharbar}{\isachargreater}\ f\ x\isactrlisub n}%
\end{isamarkuptext}%
\isamarkuptrue%
%
\isadelimmlref
%
\endisadelimmlref
%
\isatagmlref
%
\begin{isamarkuptext}%
\begin{mldecls}
  \indexml{(op |-$>$)}\verb|(op |\verb,|,\verb|->): ('c * 'a) * ('c -> 'a -> 'b) -> 'b| \\
  \indexml{(op |$>$$>$)}\verb|(op |\verb,|,\verb|>>): ('a * 'c) * ('a -> 'b) -> 'b * 'c| \\
  \indexml{(op ||$>$)}\verb|(op |\verb,|,\verb||\verb,|,\verb|>): ('c * 'a) * ('a -> 'b) -> 'c * 'b| \\
  \indexml{(op ||$>$$>$)}\verb|(op |\verb,|,\verb||\verb,|,\verb|>>): ('c * 'a) * ('a -> 'd * 'b) -> ('c * 'd) * 'b| \\
  \indexml{fold-map}\verb|fold_map: ('a -> 'b -> 'c * 'b) -> 'a list -> 'b -> 'c list * 'b| \\
  \end{mldecls}%
\end{isamarkuptext}%
\isamarkuptrue%
%
\endisatagmlref
{\isafoldmlref}%
%
\isadelimmlref
%
\endisadelimmlref
%
\begin{isamarkuptext}%
\noindent FIXME transformations involving side results%
\end{isamarkuptext}%
\isamarkuptrue%
%
\isadelimmlref
%
\endisadelimmlref
%
\isatagmlref
%
\begin{isamarkuptext}%
\begin{mldecls}
  \indexml{(op \#$>$)}\verb|(op #>): ('a -> 'b) * ('b -> 'c) -> 'a -> 'c| \\
  \indexml{(op \#-$>$)}\verb|(op #->): ('a -> 'c * 'b) * ('c -> 'b -> 'd) -> 'a -> 'd| \\
  \indexml{(op \#$>$$>$)}\verb|(op #>>): ('a -> 'c * 'b) * ('c -> 'd) -> 'a -> 'd * 'b| \\
  \indexml{(op \#\#$>$)}\verb|(op ##>): ('a -> 'c * 'b) * ('b -> 'd) -> 'a -> 'c * 'd| \\
  \indexml{(op \#\#$>$$>$)}\verb|(op ##>>): ('a -> 'c * 'b) * ('b -> 'e * 'd) -> 'a -> ('c * 'e) * 'd| \\
  \end{mldecls}%
\end{isamarkuptext}%
\isamarkuptrue%
%
\endisatagmlref
{\isafoldmlref}%
%
\isadelimmlref
%
\endisadelimmlref
%
\begin{isamarkuptext}%
\noindent All those linear combinators also exist in higher-order
  variants which do not expect a value on the left hand side
  but a function.%
\end{isamarkuptext}%
\isamarkuptrue%
%
\isadelimmlref
%
\endisadelimmlref
%
\isatagmlref
%
\begin{isamarkuptext}%
\begin{mldecls}
  \indexml{(op `)}\verb|(op `): ('b -> 'a) -> 'b -> 'a * 'b| \\
  \indexml{tap}\verb|tap: ('b -> 'a) -> 'b -> 'b| \\
  \end{mldecls}%
\end{isamarkuptext}%
\isamarkuptrue%
%
\endisatagmlref
{\isafoldmlref}%
%
\isadelimmlref
%
\endisadelimmlref
%
\begin{isamarkuptext}%
\noindent FIXME%
\end{isamarkuptext}%
\isamarkuptrue%
%
\isamarkupsection{Options and partiality%
}
\isamarkuptrue%
%
\isadelimmlref
%
\endisadelimmlref
%
\isatagmlref
%
\begin{isamarkuptext}%
\begin{mldecls}
  \indexml{is-some}\verb|is_some: 'a option -> bool| \\
  \indexml{is-none}\verb|is_none: 'a option -> bool| \\
  \indexml{the}\verb|the: 'a option -> 'a| \\
  \indexml{these}\verb|these: 'a list option -> 'a list| \\
  \indexml{the-list}\verb|the_list: 'a option -> 'a list| \\
  \indexml{the-default}\verb|the_default: 'a -> 'a option -> 'a| \\
  \indexml{try}\verb|try: ('a -> 'b) -> 'a -> 'b option| \\
  \indexml{can}\verb|can: ('a -> 'b) -> 'a -> bool| \\
  \end{mldecls}%
\end{isamarkuptext}%
\isamarkuptrue%
%
\endisatagmlref
{\isafoldmlref}%
%
\isadelimmlref
%
\endisadelimmlref
%
\begin{isamarkuptext}%
Standard selector functions on \isa{option}s are provided.
  The \verb|try| and \verb|can| functions provide a convenient
  interface for handling exceptions -- both take as arguments
  a function \isa{f} together with a parameter \isa{x}
  and catch any exception during the evaluation of the application
  of \isa{f} to \isa{x}, either return a lifted result
  (\verb|NONE| on failure) or a boolean value (\verb|false| on failure).%
\end{isamarkuptext}%
\isamarkuptrue%
%
\isamarkupsection{Common data structures%
}
\isamarkuptrue%
%
\isamarkupsubsection{Lists (as set-like data structures)%
}
\isamarkuptrue%
%
\begin{isamarkuptext}%
\begin{mldecls}
  \indexml{member}\verb|member: ('b * 'a -> bool) -> 'a list -> 'b -> bool| \\
  \indexml{insert}\verb|insert: ('a * 'a -> bool) -> 'a -> 'a list -> 'a list| \\
  \indexml{remove}\verb|remove: ('b * 'a -> bool) -> 'b -> 'a list -> 'a list| \\
  \indexml{merge}\verb|merge: ('a * 'a -> bool) -> 'a list * 'a list -> 'a list| \\
  \end{mldecls}%
\end{isamarkuptext}%
\isamarkuptrue%
%
\begin{isamarkuptext}%
Lists are often used as set-like data structures -- set-like in
  then sense that they support notion of \verb|member|-ship,
  \verb|insert|-ing and \verb|remove|-ing, but are order-sensitive.
  This is convenient when implementing a history-like mechanism:
  \verb|insert| adds an element \emph{to the front} of a list,
  if not yet present; \verb|remove| removes \emph{all} occurences
  of a particular element.  Correspondingly \verb|merge| implements a 
  a merge on two lists suitable for merges of context data
  (\secref{sec:context-theory}).

  Functions are parametrized by an explicit equality function
  to accomplish overloaded equality;  in most cases of monomorphic
  equality, writing \verb|(op =)| should suffice.%
\end{isamarkuptext}%
\isamarkuptrue%
%
\isamarkupsubsection{Association lists%
}
\isamarkuptrue%
%
\begin{isamarkuptext}%
\begin{mldecls}
  \indexmlexception{AList.DUP}\verb|exception AList.DUP| \\
  \indexml{AList.lookup}\verb|AList.lookup: ('a * 'b -> bool) -> ('b * 'c) list -> 'a -> 'c option| \\
  \indexml{AList.defined}\verb|AList.defined: ('a * 'b -> bool) -> ('b * 'c) list -> 'a -> bool| \\
  \indexml{AList.update}\verb|AList.update: ('a * 'a -> bool) -> ('a * 'b) -> ('a * 'b) list -> ('a * 'b) list| \\
  \indexml{AList.default}\verb|AList.default: ('a * 'a -> bool) -> ('a * 'b) -> ('a * 'b) list -> ('a * 'b) list| \\
  \indexml{AList.delete}\verb|AList.delete: ('a * 'b -> bool) -> 'a -> ('b * 'c) list -> ('b * 'c) list| \\
  \indexml{AList.map-entry}\verb|AList.map_entry: ('a * 'b -> bool) -> 'a|\isasep\isanewline%
\verb|    -> ('c -> 'c) -> ('b * 'c) list -> ('b * 'c) list| \\
  \indexml{AList.map-default}\verb|AList.map_default: ('a * 'a -> bool) -> 'a * 'b -> ('b -> 'b)|\isasep\isanewline%
\verb|    -> ('a * 'b) list -> ('a * 'b) list| \\
  \indexml{AList.join}\verb|AList.join: ('a * 'a -> bool) -> ('a -> 'b * 'b -> 'b) (*exception DUP*)|\isasep\isanewline%
\verb|    -> ('a * 'b) list * ('a * 'b) list -> ('a * 'b) list (*exception AList.DUP*)| \\
  \indexml{AList.merge}\verb|AList.merge: ('a * 'a -> bool) -> ('b * 'b -> bool)|\isasep\isanewline%
\verb|    -> ('a * 'b) list * ('a * 'b) list -> ('a * 'b) list (*exception AList.DUP*)|
  \end{mldecls}%
\end{isamarkuptext}%
\isamarkuptrue%
%
\begin{isamarkuptext}%
Association lists can be seens as an extension of set-like lists:
  on the one hand, they may be used to implement finite mappings,
  on the other hand, they remain order-sensitive and allow for
  multiple key-value-pair with the same key: \verb|AList.lookup|
  returns the \emph{first} value corresponding to a particular
  key, if present.  \verb|AList.update| updates
  the \emph{first} occurence of a particular key; if no such
  key exists yet, the key-value-pair is added \emph{to the front}.
  \verb|AList.delete| only deletes the \emph{first} occurence of a key.
  \verb|AList.merge| provides an operation suitable for merges of context data
  (\secref{sec:context-theory}), where an equality parameter on
  values determines whether a merge should be considered a conflict.
  A slightly generalized operation if implementend by the \verb|AList.join|
  function which allows for explicit conflict resolution.%
\end{isamarkuptext}%
\isamarkuptrue%
%
\isamarkupsubsection{Tables%
}
\isamarkuptrue%
%
\begin{isamarkuptext}%
\begin{mldecls}
  \indexmltype{Symtab.key}\verb|type Symtab.key| \\
  \indexmltype{'a Symtab.table}\verb|type 'a Symtab.table| \\
  \indexmlexception{Symtab.DUP}\verb|exception Symtab.DUP of Symtab.key| \\
  \indexmlexception{Symtab.DUPS}\verb|exception Symtab.DUPS of Symtab.key list| \\
  \indexmlexception{Symtab.SAME}\verb|exception Symtab.SAME| \\
  \indexmlexception{Symtab.UNDEF}\verb|exception Symtab.UNDEF of Symtab.key| \\
  \indexml{Symtab.empty}\verb|Symtab.empty: 'a Symtab.table| \\
  \indexml{Symtab.dest}\verb|Symtab.dest: 'a Symtab.table -> (Symtab.key * 'a) list| \\
  \indexml{Symtab.keys}\verb|Symtab.keys: 'a Symtab.table -> Symtab.key list| \\
  \indexml{Symtab.lookup}\verb|Symtab.lookup: 'a Symtab.table -> Symtab.key -> 'a option| \\
  \indexml{Symtab.defined}\verb|Symtab.defined: 'a Symtab.table -> Symtab.key -> bool| \\
  \indexml{Symtab.update}\verb|Symtab.update: (Symtab.key * 'a) -> 'a Symtab.table -> 'a Symtab.table| \\
  \indexml{Symtab.default}\verb|Symtab.default: Symtab.key * 'a -> 'a Symtab.table -> 'a Symtab.table| \\
  \indexml{Symtab.delete}\verb|Symtab.delete: Symtab.key|\isasep\isanewline%
\verb|    -> 'a Symtab.table -> 'a Symtab.table (*exception Symtab.UNDEF*)| \\
  \indexml{Symtab.map-entry}\verb|Symtab.map_entry: Symtab.key -> ('a -> 'a)|\isasep\isanewline%
\verb|    -> 'a Symtab.table -> 'a Symtab.table| \\
  \indexml{Symtab.map-default}\verb|Symtab.map_default: (Symtab.key * 'a) -> ('a -> 'a)|\isasep\isanewline%
\verb|    -> 'a Symtab.table -> 'a Symtab.table| \\
  \indexml{Symtab.join}\verb|Symtab.join: (Symtab.key -> 'a * 'a -> 'a) (*exception Symtab.DUP/Symtab.SAME*)|\isasep\isanewline%
\verb|    -> 'a Symtab.table * 'a Symtab.table|\isasep\isanewline%
\verb|    -> 'a Symtab.table (*exception Symtab.DUPS*)| \\
  \indexml{Symtab.merge}\verb|Symtab.merge: ('a * 'a -> bool)|\isasep\isanewline%
\verb|    -> 'a Symtab.table * 'a Symtab.table|\isasep\isanewline%
\verb|    -> 'a Symtab.table (*exception Symtab.DUPS*)|
  \end{mldecls}%
\end{isamarkuptext}%
\isamarkuptrue%
%
\begin{isamarkuptext}%
Tables are an efficient representation of finite mappings without
  any notion of order;  due to their efficiency they should be used
  whenever such pure finite mappings are neccessary.

  The key type of tables must be given explicitly by instantiating
  the \verb|TableFun| functor which takes the key type
  together with its \verb|order|; for convience, we restrict
  here to the \verb|Symtab| instance with \verb|string|
  as key type.

  Most table functions correspond to those of association lists.%
\end{isamarkuptext}%
\isamarkuptrue%
%
\isamarkupchapter{Cookbook%
}
\isamarkuptrue%
%
\isamarkupsection{A method that depends on declarations in the context%
}
\isamarkuptrue%
%
\begin{isamarkuptext}%
FIXME%
\end{isamarkuptext}%
\isamarkuptrue%
%
\isadelimtheory
%
\endisadelimtheory
%
\isatagtheory
\isacommand{end}\isamarkupfalse%
%
\endisatagtheory
{\isafoldtheory}%
%
\isadelimtheory
%
\endisadelimtheory
\isanewline
\end{isabellebody}%
%%% Local Variables:
%%% mode: latex
%%% TeX-master: "root"
%%% End:


\begingroup
\tocentry{\bibname}
\bibliographystyle{plain} \small\raggedright\frenchspacing
\bibliography{../manual}
\endgroup

%FIXME
%\tocentry{\glossaryname}
%\printglossary

\tocentry{\indexname}
\printindex

\end{document}


%%% Local Variables: 
%%% mode: latex
%%% TeX-master: t
%%% End: 
